
\documentclass[11pt,a4paper]{article}
\usepackage[T1]{fontenc}
\usepackage[utf8]{inputenc}
\usepackage[french]{babel}
\usepackage{lmodern}
\usepackage{microtype}
\usepackage{geometry}
\geometry{margin=2.5cm}
\usepackage{graphicx}
\usepackage{caption}
\usepackage{subcaption}
\usepackage{booktabs}
\usepackage{xcolor}
\usepackage{hyperref}
\hypersetup{colorlinks=true,linkcolor=blue,urlcolor=blue,citecolor=blue}
\title{Référence formantique complète\\\large Logique des doublures instrumentales classiques}
\author{Yan Maresz}
\date{Mars 2026}

\begin{document}
\maketitle
\tableofcontents
\clearpage

\section{Introduction et Méthodologie}
NOTE DE MISE À JOUR (5 mars 2026) — v8⚠ Audit global validé: 16/16 instruments CSV reproduits à Δ=0 Hz. v7: colonne Fp (centroïde spectral) + 14 instruments Yan\_Adds intégrés. v8: graphiques formantiques (barres F1-F6 + Fp + zones vocaliques) ajoutés pour les 14 instruments Yan\_Adds dans les fiches correspondantes. Les valeurs F1-F4 de la v6 sont confirmées correctes (Δ=0 Hz vs données brutes). RÉFÉRENCE FORMANTIQUE COMPLÈTE Logique des Doublures Instrumentales Classiques Analyse par les Formants Spectraux — 29 instruments — Synthèse Quadri-Source Backus 1969  •  Giesler 1985  •  Meyer 2009  •  SOL2020 IRCAM  •  Yan\_Adds Étude quantitative — 5 914 échantillons — Mars 2026 — v8 + Fp centroïde + Yan\_Adds + graphiques TABLE DES MATIÈRES I.  Introduction et Méthodologie II.  Les Cuivres III.  Les Bois IV.  Les Saxophones V.  Les Cordes Solistes VI.  Les Cordes d'Ensemble VII.  Synthèse — Convergences Formantiques (avec graphiques) VIII.  Espace Vocalique Complet IX.  Validation et Comparaison avec la Littérature X.  Applications à l'Orchestration XI.  Tableau des Doublures Vérifiées XII.  Principes d'Orchestration Acoustique Conclusion Annexe: Sources et Références
\par\medskip

egin{figure}[htbp]
\centering
\includegraphics[width=\linewidth]{figures/conclusion__page_01__p01_01.jpg}
\caption{Illustration (page 1)}
\end{figure}

egin{figure}[htbp]
\centering
\includegraphics[width=\linewidth]{figures/conclusion__page_01__p01_02.jpg}
\caption{Illustration (page 1)}
\end{figure}

egin{figure}[htbp]
\centering
\includegraphics[width=\linewidth]{figures/conclusion__page_01__p01_03.png}
\caption{Illustration (page 1)}
\end{figure}

egin{figure}[htbp]
\centering
\includegraphics[width=\linewidth]{figures/conclusion__page_01__p01_04.jpg}
\caption{Illustration (page 1)}
\end{figure}

egin{figure}[htbp]
\centering
\includegraphics[width=\linewidth]{figures/conclusion__page_01__p01_05.jpg}
\caption{Illustration (page 1)}
\end{figure}

egin{figure}[htbp]
\centering
\includegraphics[width=\linewidth]{figures/conclusion__page_01__p01_06.jpg}
\caption{Illustration (page 1)}
\end{figure}

egin{figure}[htbp]
\centering
\includegraphics[width=\linewidth]{figures/conclusion__page_01__p01_07.jpg}
\caption{Illustration (page 1)}
\end{figure}

egin{figure}[htbp]
\centering
\includegraphics[width=\linewidth]{figures/conclusion__page_01__p01_08.png}
\caption{Illustration (page 1)}
\end{figure}

egin{figure}[htbp]
\centering
\includegraphics[width=\linewidth]{figures/conclusion__page_01__p01_09.jpg}
\caption{Illustration (page 1)}
\end{figure}

egin{figure}[htbp]
\centering
\includegraphics[width=\linewidth]{figures/conclusion__page_01__p01_10.jpg}
\caption{Illustration (page 1)}
\end{figure}

egin{figure}[htbp]
\centering
\includegraphics[width=\linewidth]{figures/conclusion__page_01__p01_11.png}
\caption{Illustration (page 1)}
\end{figure}

egin{figure}[htbp]
\centering
\includegraphics[width=\linewidth]{figures/conclusion__page_01__p01_12.png}
\caption{Illustration (page 1)}
\end{figure}

egin{figure}[htbp]
\centering
\includegraphics[width=\linewidth]{figures/conclusion__page_01__p01_13.png}
\caption{Illustration (page 1)}
\end{figure}

egin{figure}[htbp]
\centering
\includegraphics[width=\linewidth]{figures/conclusion__page_01__p01_14.png}
\caption{Illustration (page 1)}
\end{figure}

egin{figure}[htbp]
\centering
\includegraphics[width=\linewidth]{figures/conclusion__page_01__p01_15.png}
\caption{Illustration (page 1)}
\end{figure}

egin{figure}[htbp]
\centering
\includegraphics[width=\linewidth]{figures/conclusion__page_01__p01_16.png}
\caption{Illustration (page 1)}
\end{figure}

egin{figure}[htbp]
\centering
\includegraphics[width=\linewidth]{figures/conclusion__page_01__p01_17.png}
\caption{Illustration (page 1)}
\end{figure}

egin{figure}[htbp]
\centering
\includegraphics[width=\linewidth]{figures/conclusion__page_01__p01_18.png}
\caption{Illustration (page 1)}
\end{figure}

egin{figure}[htbp]
\centering
\includegraphics[width=\linewidth]{figures/conclusion__page_01__p01_19.png}
\caption{Illustration (page 1)}
\end{figure}

egin{figure}[htbp]
\centering
\includegraphics[width=\linewidth]{figures/conclusion__page_01__p01_20.jpg}
\caption{Illustration (page 1)}
\end{figure}

egin{figure}[htbp]
\centering
\includegraphics[width=\linewidth]{figures/conclusion__page_01__p01_21.jpg}
\caption{Illustration (page 1)}
\end{figure}

egin{figure}[htbp]
\centering
\includegraphics[width=\linewidth]{figures/conclusion__page_01__p01_22.jpg}
\caption{Illustration (page 1)}
\end{figure}

egin{figure}[htbp]
\centering
\includegraphics[width=\linewidth]{figures/conclusion__page_01__p01_23.jpg}
\caption{Illustration (page 1)}
\end{figure}

egin{figure}[htbp]
\centering
\includegraphics[width=\linewidth]{figures/conclusion__page_01__p01_24.jpg}
\caption{Illustration (page 1)}
\end{figure}

egin{figure}[htbp]
\centering
\includegraphics[width=\linewidth]{figures/conclusion__page_01__p01_25.jpg}
\caption{Illustration (page 1)}
\end{figure}

egin{figure}[htbp]
\centering
\includegraphics[width=\linewidth]{figures/conclusion__page_01__p01_26.jpg}
\caption{Illustration (page 1)}
\end{figure}

egin{figure}[htbp]
\centering
\includegraphics[width=\linewidth]{figures/conclusion__page_01__p01_27.png}
\caption{Illustration (page 1)}
\end{figure}

egin{figure}[htbp]
\centering
\includegraphics[width=\linewidth]{figures/conclusion__page_01__p01_28.png}
\caption{Illustration (page 1)}
\end{figure}

egin{figure}[htbp]
\centering
\includegraphics[width=\linewidth]{figures/conclusion__page_01__p01_29.png}
\caption{Illustration (page 1)}
\end{figure}

egin{figure}[htbp]
\centering
\includegraphics[width=\linewidth]{figures/conclusion__page_01__p01_30.png}
\caption{Illustration (page 1)}
\end{figure}

egin{figure}[htbp]
\centering
\includegraphics[width=\linewidth]{figures/conclusion__page_01__p01_31.png}
\caption{Illustration (page 1)}
\end{figure}

egin{figure}[htbp]
\centering
\includegraphics[width=\linewidth]{figures/conclusion__page_01__p01_32.jpg}
\caption{Illustration (page 1)}
\end{figure}

egin{figure}[htbp]
\centering
\includegraphics[width=\linewidth]{figures/conclusion__page_01__p01_33.png}
\caption{Illustration (page 1)}
\end{figure}

\section{Introduction et Méthodologie}
I. Introduction et Méthodologie Ce document examine comment les doublures orchestrales classiques s'expliquent par l'assemblage complémentaire des formants spectraux instrumentaux. Les formants sont les zones de résonance qui définissent le timbre. Lorsque deux instruments sont doublés, leurs spectres se combinent pour créer une couleur timbrale homogène. L'analyse repose sur 4 037 échantillons professionnels couvrant 29 instruments de l'orchestre, validée contre quatre sources académiques indépendantes. Le taux de concordance global est de 93 \% (27/29 instruments). Cette base de données constitue la référence quantitative la plus complète disponible à ce jour pour l'orchestration scientifique. Sources des données Source Titre / Base Description Éditeur Backus (1969)The Acoustical Foundations of MusicRéférence historique, mesures ciblées cuivres et boisW.W. Norton \& Co. Giesler (1985)Instrumentation: Ein Hand- und LehrbuchMesures détaillées, associations vocaliques explicites, 13 instrumentsBreitkopf \& Härtel Meyer (2009)Acoustics and the Performance of Music (5e éd.)Analyse acoustique avec couleurs vocales, zones vocaliques définiesSpringer-Verlag SOL2020 IRCAMStudio On Line / Orchestre de Lyon12 instruments solistes, 2 391 échantillons sustained analysésIRCAM / Orchestre de Lyon Yan\_AddsExtension personnelle17 instruments supplémentaires, 1 646 échantillons (ensembles, bois auxiliaires, cuivres graves)Analyse propre Méthodologie d'extraction formantique Les analyses SOL2020 précédentes contenaient des erreurs systématiques: inclusion de techniques étendues, et utilisation du centroïde spectral au lieu de la détection de pics. La présente analyse corrige ces deux problèmes. Filtrage des données Techniques INCLUSES: normal, arco, sustain, ordinario, vibrato, non-vibrato; dynamiques pp à ff Techniques EXCLUES: pizzicato, col legno, tremolo, trilles, harmoniques artificiels, multiphoniques, flutter tongue, sourdines, glissandi, portamento Justification: Seules les techniques sustained produisent les formants structurels stables de l'instrument. Les techniques étendues modifient radicalement le spectre et ne représentent pas le timbre fondamental. Paramètres d'analyse Paramètre Valeur / Description FFT Size 4 096 samples Sample Rate 44 100 Hz Résolution fréquentielle \textasciitilde{}10,8 Hz/bin Plage analysée 100–3 000 Hz Algorithme scipy.signal.find\_peaks — pic le plus proéminent (pas premier pic) Hauteur minimale 10 \% de l'amplitude maximale
\par\medskip

egin{figure}[htbp]
\centering
\includegraphics[width=\linewidth]{figures/introduction__page_02__p02_01.jpg}
\caption{Illustration (page 2)}
\end{figure}

egin{figure}[htbp]
\centering
\includegraphics[width=\linewidth]{figures/introduction__page_02__p02_02.jpg}
\caption{Illustration (page 2)}
\end{figure}

egin{figure}[htbp]
\centering
\includegraphics[width=\linewidth]{figures/introduction__page_02__p02_03.png}
\caption{Illustration (page 2)}
\end{figure}

egin{figure}[htbp]
\centering
\includegraphics[width=\linewidth]{figures/introduction__page_02__p02_04.jpg}
\caption{Illustration (page 2)}
\end{figure}

egin{figure}[htbp]
\centering
\includegraphics[width=\linewidth]{figures/introduction__page_02__p02_05.jpg}
\caption{Illustration (page 2)}
\end{figure}

egin{figure}[htbp]
\centering
\includegraphics[width=\linewidth]{figures/introduction__page_02__p02_06.jpg}
\caption{Illustration (page 2)}
\end{figure}

egin{figure}[htbp]
\centering
\includegraphics[width=\linewidth]{figures/introduction__page_02__p02_07.jpg}
\caption{Illustration (page 2)}
\end{figure}

egin{figure}[htbp]
\centering
\includegraphics[width=\linewidth]{figures/introduction__page_02__p02_08.png}
\caption{Illustration (page 2)}
\end{figure}

egin{figure}[htbp]
\centering
\includegraphics[width=\linewidth]{figures/introduction__page_02__p02_09.jpg}
\caption{Illustration (page 2)}
\end{figure}

egin{figure}[htbp]
\centering
\includegraphics[width=\linewidth]{figures/introduction__page_02__p02_10.jpg}
\caption{Illustration (page 2)}
\end{figure}

egin{figure}[htbp]
\centering
\includegraphics[width=\linewidth]{figures/introduction__page_02__p02_11.png}
\caption{Illustration (page 2)}
\end{figure}

egin{figure}[htbp]
\centering
\includegraphics[width=\linewidth]{figures/introduction__page_02__p02_12.png}
\caption{Illustration (page 2)}
\end{figure}

egin{figure}[htbp]
\centering
\includegraphics[width=\linewidth]{figures/introduction__page_02__p02_13.png}
\caption{Illustration (page 2)}
\end{figure}

egin{figure}[htbp]
\centering
\includegraphics[width=\linewidth]{figures/introduction__page_02__p02_14.png}
\caption{Illustration (page 2)}
\end{figure}

egin{figure}[htbp]
\centering
\includegraphics[width=\linewidth]{figures/introduction__page_02__p02_15.png}
\caption{Illustration (page 2)}
\end{figure}

egin{figure}[htbp]
\centering
\includegraphics[width=\linewidth]{figures/introduction__page_02__p02_16.png}
\caption{Illustration (page 2)}
\end{figure}

egin{figure}[htbp]
\centering
\includegraphics[width=\linewidth]{figures/introduction__page_02__p02_17.png}
\caption{Illustration (page 2)}
\end{figure}

egin{figure}[htbp]
\centering
\includegraphics[width=\linewidth]{figures/introduction__page_02__p02_18.png}
\caption{Illustration (page 2)}
\end{figure}

egin{figure}[htbp]
\centering
\includegraphics[width=\linewidth]{figures/introduction__page_02__p02_19.png}
\caption{Illustration (page 2)}
\end{figure}

egin{figure}[htbp]
\centering
\includegraphics[width=\linewidth]{figures/introduction__page_02__p02_20.jpg}
\caption{Illustration (page 2)}
\end{figure}

egin{figure}[htbp]
\centering
\includegraphics[width=\linewidth]{figures/introduction__page_02__p02_21.jpg}
\caption{Illustration (page 2)}
\end{figure}

egin{figure}[htbp]
\centering
\includegraphics[width=\linewidth]{figures/introduction__page_02__p02_22.jpg}
\caption{Illustration (page 2)}
\end{figure}

egin{figure}[htbp]
\centering
\includegraphics[width=\linewidth]{figures/introduction__page_02__p02_23.jpg}
\caption{Illustration (page 2)}
\end{figure}

egin{figure}[htbp]
\centering
\includegraphics[width=\linewidth]{figures/introduction__page_02__p02_24.jpg}
\caption{Illustration (page 2)}
\end{figure}

egin{figure}[htbp]
\centering
\includegraphics[width=\linewidth]{figures/introduction__page_02__p02_25.jpg}
\caption{Illustration (page 2)}
\end{figure}

egin{figure}[htbp]
\centering
\includegraphics[width=\linewidth]{figures/introduction__page_02__p02_26.jpg}
\caption{Illustration (page 2)}
\end{figure}

egin{figure}[htbp]
\centering
\includegraphics[width=\linewidth]{figures/introduction__page_02__p02_27.png}
\caption{Illustration (page 2)}
\end{figure}

egin{figure}[htbp]
\centering
\includegraphics[width=\linewidth]{figures/introduction__page_02__p02_28.png}
\caption{Illustration (page 2)}
\end{figure}

egin{figure}[htbp]
\centering
\includegraphics[width=\linewidth]{figures/introduction__page_02__p02_29.png}
\caption{Illustration (page 2)}
\end{figure}

egin{figure}[htbp]
\centering
\includegraphics[width=\linewidth]{figures/introduction__page_02__p02_30.png}
\caption{Illustration (page 2)}
\end{figure}

egin{figure}[htbp]
\centering
\includegraphics[width=\linewidth]{figures/introduction__page_02__p02_31.png}
\caption{Illustration (page 2)}
\end{figure}

egin{figure}[htbp]
\centering
\includegraphics[width=\linewidth]{figures/introduction__page_02__p02_32.jpg}
\caption{Illustration (page 2)}
\end{figure}

egin{figure}[htbp]
\centering
\includegraphics[width=\linewidth]{figures/introduction__page_02__p02_33.png}
\caption{Illustration (page 2)}
\end{figure}

\section*{SECTION II — Analyse spectrale complète}
Paramètre Valeur / Description Distance minimale 10 bins (\textasciitilde{}50 Hz) Proéminence minimale 5 \% de l'amplitude maximale Choix méthodologique critique: la détection du pic le plus proéminent (et non du premier pic) capture le formant perceptuel principal, évite les résonances basses secondaires, et est cohérente avec la perception auditive. [Note v4 — Section II]: Cette convention (Fp = formant le plus proéminent) diverge du F1 au sens spectral strict pour 8 instruments: Trombone, Violoncelle, Saxophone alto, Trompette, Clarinette Sib, Hautbois, Flûte, Violon. Pour ces instruments, le Fp mesuré ici correspond au F2 (voire F3) dans la nomenclature spectrale, le premier mode étant acoustiquement moins rayonné. Chaque fiche concernée comporte une [Note v4] expliquant le mécanisme physique. La Section II de ce document présente l'analyse spectrale complète (F1–F6 stricts) sur l'intégralité du corpus SOL2020. [Note v7 — Fp centroïde]: L'audit global (mars 2026) a validé le pipeline complet (16/16 instruments CSV reproduits à Δ=0 Hz). Une nouvelle mesure Fp (Formant Principal par centroïde de bande spectrale) est introduite pour les instruments dont F2 est instable par registre (σ>200 Hz). Fp = centroïde spectral pondéré en énergie dans une bande choisie par instrument (typiquement 600-1400 ou 1000-2000 Hz). Cas démonstratifs: Violon Fp=893 Hz (σ=92) vs F2=1518 Hz (σ=651); Trompette Fp=1046 Hz (σ=98) vs F2=1324 Hz (σ=1018); Clarinette Fp=1412 Hz (σ=169) vs F2=1130 Hz (σ=795). Fp est systématiquement plus stable que F2 et correspond au Hauptformant perceptif de Meyer pour les cordes. 22/30 instruments audités bénéficient de Fp. Grille de concordance SymboleSeuil Signification ✓✓< 80 Hz Excellente concordance ✓80–200 Hz Bonne concordance \textasciitilde{}200–400 Hz Modérée ✗> 400 Hz Divergente *— Pas de formant fixe Correspondance voyelles–fréquences (Meyer 2009) VoyelleFréquence (Hz)Perceptio nInstruments (F1 SOL2020/Yan\_Adds) u (oo)200–400Profondeur Tuba 249, Contrebasse 200, CB ensemble 350 o (oh)400–600PlénitudeCor 457, CB Tuba 471, Cbn 478, Trb 491, Vcl 499, Bsn 502, Vcl-Ens 587 å (aw)\textasciitilde{}600Transition Vcl-Ens sord. 587, Vla-sord 683, Fl. basse 737 a (ah)800–1 250PuissanceTrb.b 894, Cl.b 909, Cl Sib 1 016, Cor angl. 1 045, Vla-Ens 1 190 e (eh)1 250–2 600ClartéFl 1 354, Htb 1 460, Vla 369, Vln-Ens 1 556, Cl Mib 1 747, Picc 2 336, Vln 1 034 i (ee)> 2 600Brillance Trp 1 324
\par\medskip

egin{figure}[htbp]
\centering
\includegraphics[width=\linewidth]{figures/introduction__trombone__p03_01.jpg}
\caption{Illustration (page 3)}
\end{figure}

egin{figure}[htbp]
\centering
\includegraphics[width=\linewidth]{figures/introduction__trombone__p03_02.jpg}
\caption{Illustration (page 3)}
\end{figure}

egin{figure}[htbp]
\centering
\includegraphics[width=\linewidth]{figures/introduction__trombone__p03_03.png}
\caption{Illustration (page 3)}
\end{figure}

egin{figure}[htbp]
\centering
\includegraphics[width=\linewidth]{figures/introduction__trombone__p03_04.jpg}
\caption{Illustration (page 3)}
\end{figure}

egin{figure}[htbp]
\centering
\includegraphics[width=\linewidth]{figures/introduction__trombone__p03_05.jpg}
\caption{Illustration (page 3)}
\end{figure}

egin{figure}[htbp]
\centering
\includegraphics[width=\linewidth]{figures/introduction__trombone__p03_06.jpg}
\caption{Illustration (page 3)}
\end{figure}

egin{figure}[htbp]
\centering
\includegraphics[width=\linewidth]{figures/introduction__trombone__p03_07.jpg}
\caption{Illustration (page 3)}
\end{figure}

egin{figure}[htbp]
\centering
\includegraphics[width=\linewidth]{figures/introduction__trombone__p03_08.png}
\caption{Illustration (page 3)}
\end{figure}

egin{figure}[htbp]
\centering
\includegraphics[width=\linewidth]{figures/introduction__trombone__p03_09.jpg}
\caption{Illustration (page 3)}
\end{figure}

egin{figure}[htbp]
\centering
\includegraphics[width=\linewidth]{figures/introduction__trombone__p03_10.jpg}
\caption{Illustration (page 3)}
\end{figure}

egin{figure}[htbp]
\centering
\includegraphics[width=\linewidth]{figures/introduction__trombone__p03_11.png}
\caption{Illustration (page 3)}
\end{figure}

egin{figure}[htbp]
\centering
\includegraphics[width=\linewidth]{figures/introduction__trombone__p03_12.png}
\caption{Illustration (page 3)}
\end{figure}

egin{figure}[htbp]
\centering
\includegraphics[width=\linewidth]{figures/introduction__trombone__p03_13.png}
\caption{Illustration (page 3)}
\end{figure}

egin{figure}[htbp]
\centering
\includegraphics[width=\linewidth]{figures/introduction__trombone__p03_14.png}
\caption{Illustration (page 3)}
\end{figure}

egin{figure}[htbp]
\centering
\includegraphics[width=\linewidth]{figures/introduction__trombone__p03_15.png}
\caption{Illustration (page 3)}
\end{figure}

egin{figure}[htbp]
\centering
\includegraphics[width=\linewidth]{figures/introduction__trombone__p03_16.png}
\caption{Illustration (page 3)}
\end{figure}

egin{figure}[htbp]
\centering
\includegraphics[width=\linewidth]{figures/introduction__trombone__p03_17.png}
\caption{Illustration (page 3)}
\end{figure}

egin{figure}[htbp]
\centering
\includegraphics[width=\linewidth]{figures/introduction__trombone__p03_18.png}
\caption{Illustration (page 3)}
\end{figure}

egin{figure}[htbp]
\centering
\includegraphics[width=\linewidth]{figures/introduction__trombone__p03_19.png}
\caption{Illustration (page 3)}
\end{figure}

egin{figure}[htbp]
\centering
\includegraphics[width=\linewidth]{figures/introduction__trombone__p03_20.jpg}
\caption{Illustration (page 3)}
\end{figure}

egin{figure}[htbp]
\centering
\includegraphics[width=\linewidth]{figures/introduction__trombone__p03_21.jpg}
\caption{Illustration (page 3)}
\end{figure}

egin{figure}[htbp]
\centering
\includegraphics[width=\linewidth]{figures/introduction__trombone__p03_22.jpg}
\caption{Illustration (page 3)}
\end{figure}

egin{figure}[htbp]
\centering
\includegraphics[width=\linewidth]{figures/introduction__trombone__p03_23.jpg}
\caption{Illustration (page 3)}
\end{figure}

egin{figure}[htbp]
\centering
\includegraphics[width=\linewidth]{figures/introduction__trombone__p03_24.jpg}
\caption{Illustration (page 3)}
\end{figure}

egin{figure}[htbp]
\centering
\includegraphics[width=\linewidth]{figures/introduction__trombone__p03_25.jpg}
\caption{Illustration (page 3)}
\end{figure}

egin{figure}[htbp]
\centering
\includegraphics[width=\linewidth]{figures/introduction__trombone__p03_26.jpg}
\caption{Illustration (page 3)}
\end{figure}

egin{figure}[htbp]
\centering
\includegraphics[width=\linewidth]{figures/introduction__trombone__p03_27.png}
\caption{Illustration (page 3)}
\end{figure}

egin{figure}[htbp]
\centering
\includegraphics[width=\linewidth]{figures/introduction__trombone__p03_28.png}
\caption{Illustration (page 3)}
\end{figure}

egin{figure}[htbp]
\centering
\includegraphics[width=\linewidth]{figures/introduction__trombone__p03_29.png}
\caption{Illustration (page 3)}
\end{figure}

egin{figure}[htbp]
\centering
\includegraphics[width=\linewidth]{figures/introduction__trombone__p03_30.png}
\caption{Illustration (page 3)}
\end{figure}

egin{figure}[htbp]
\centering
\includegraphics[width=\linewidth]{figures/introduction__trombone__p03_31.png}
\caption{Illustration (page 3)}
\end{figure}

egin{figure}[htbp]
\centering
\includegraphics[width=\linewidth]{figures/introduction__trombone__p03_32.jpg}
\caption{Illustration (page 3)}
\end{figure}

egin{figure}[htbp]
\centering
\includegraphics[width=\linewidth]{figures/introduction__trombone__p03_33.png}
\caption{Illustration (page 3)}
\end{figure}

\section{Les Cuivres}
\subsection{Tuba (Bass Tuba)}
II. Les Cuivres Plage formantique: 249–1 324 Hz (voyelles /u/ → /e/). Caractéristique: formants bien définis, grande stabilité inter-dynamique. Les cuivres couvrent 4 zones vocaliques distinctes, permettant une large palette timbrale. Le cluster /o/ (Cor + Trombone) constitue la zone de mélange optimal. Tuba (Bass Tuba) SOL2020: 45 échantillons soutenus. SourceF1 (Hz)F2 (Hz)F3 (Hz)F4 (Hz)VoyelleNAccor d Backus200– 400————— Meyer210– 250———< u (oo)— SOL2020249 ±51494 ±110919 ±1191 210 ±142u (oo)45✓✓ Concordance tri-source excellente. SOL F1=249 Hz au centre de Meyer (210–250) et Backus (200–400). Correction: ancienne analyse donnait 193 Hz (techniques étendues incluses). Doublures •Contrebasse: F1 249 vs 200 Hz (Δ=49) — ancrage acoustique massif des fondations. •Tuba F2 (494) ≈ Basson F1 (502) — Δ=9 Hz, pont bois-cuivres remarquable. Tuba contrebasse (Contrabass Tuba) — NOUVEAU Yan\_Adds: 158 échantillons ordinario. SourceF1 (Hz)F2 (Hz)F3 (Hz)F4 (Hz)VoyelleNAccor d Backus/Meyer——————— Yan\_Adds471 ±1551 304 ±5762 317 ±11143 319 ±1271o (oh)158— DÉCOUVERTE: F1=471 Hz place le tuba contrebasse en plein cœur du cluster 450–502 Hz, aux côtés du cor et du trombone. Le F1 est bien plus haut que celui du tuba classique (249 Hz), reflétant une colonne d'air différente et un pavillon plus large. Variation par dynamique: p=440, mf=484, ff=566 Hz.
\par\medskip

egin{figure}[htbp]
\centering
\includegraphics[width=\linewidth]{figures/cuivres__trombone__p04_01.jpg}
\caption{Illustration (page 4)}
\end{figure}

egin{figure}[htbp]
\centering
\includegraphics[width=\linewidth]{figures/cuivres__trombone__p04_02.jpg}
\caption{Illustration (page 4)}
\end{figure}

egin{figure}[htbp]
\centering
\includegraphics[width=\linewidth]{figures/cuivres__trombone__p04_03.png}
\caption{Illustration (page 4)}
\end{figure}

egin{figure}[htbp]
\centering
\includegraphics[width=\linewidth]{figures/cuivres__trombone__p04_04.jpg}
\caption{Illustration (page 4)}
\end{figure}

egin{figure}[htbp]
\centering
\includegraphics[width=\linewidth]{figures/cuivres__trombone__p04_05.jpg}
\caption{Illustration (page 4)}
\end{figure}

egin{figure}[htbp]
\centering
\includegraphics[width=\linewidth]{figures/cuivres__trombone__p04_06.jpg}
\caption{Illustration (page 4)}
\end{figure}

egin{figure}[htbp]
\centering
\includegraphics[width=\linewidth]{figures/cuivres__trombone__p04_07.jpg}
\caption{Illustration (page 4)}
\end{figure}

egin{figure}[htbp]
\centering
\includegraphics[width=\linewidth]{figures/cuivres__trombone__p04_08.png}
\caption{Illustration (page 4)}
\end{figure}

egin{figure}[htbp]
\centering
\includegraphics[width=\linewidth]{figures/cuivres__trombone__p04_09.jpg}
\caption{Illustration (page 4)}
\end{figure}

egin{figure}[htbp]
\centering
\includegraphics[width=\linewidth]{figures/cuivres__trombone__p04_10.jpg}
\caption{Illustration (page 4)}
\end{figure}

egin{figure}[htbp]
\centering
\includegraphics[width=\linewidth]{figures/cuivres__trombone__p04_11.png}
\caption{Illustration (page 4)}
\end{figure}

egin{figure}[htbp]
\centering
\includegraphics[width=\linewidth]{figures/cuivres__trombone__p04_12.png}
\caption{Illustration (page 4)}
\end{figure}

egin{figure}[htbp]
\centering
\includegraphics[width=\linewidth]{figures/cuivres__trombone__p04_13.png}
\caption{Illustration (page 4)}
\end{figure}

egin{figure}[htbp]
\centering
\includegraphics[width=\linewidth]{figures/cuivres__trombone__p04_14.png}
\caption{Illustration (page 4)}
\end{figure}

egin{figure}[htbp]
\centering
\includegraphics[width=\linewidth]{figures/cuivres__trombone__p04_15.png}
\caption{Illustration (page 4)}
\end{figure}

egin{figure}[htbp]
\centering
\includegraphics[width=\linewidth]{figures/cuivres__trombone__p04_16.png}
\caption{Illustration (page 4)}
\end{figure}

egin{figure}[htbp]
\centering
\includegraphics[width=\linewidth]{figures/cuivres__trombone__p04_17.png}
\caption{Illustration (page 4)}
\end{figure}

egin{figure}[htbp]
\centering
\includegraphics[width=\linewidth]{figures/cuivres__trombone__p04_18.png}
\caption{Illustration (page 4)}
\end{figure}

egin{figure}[htbp]
\centering
\includegraphics[width=\linewidth]{figures/cuivres__trombone__p04_19.png}
\caption{Illustration (page 4)}
\end{figure}

egin{figure}[htbp]
\centering
\includegraphics[width=\linewidth]{figures/cuivres__trombone__p04_20.jpg}
\caption{Illustration (page 4)}
\end{figure}

egin{figure}[htbp]
\centering
\includegraphics[width=\linewidth]{figures/cuivres__trombone__p04_21.jpg}
\caption{Illustration (page 4)}
\end{figure}

egin{figure}[htbp]
\centering
\includegraphics[width=\linewidth]{figures/cuivres__trombone__p04_22.jpg}
\caption{Illustration (page 4)}
\end{figure}

egin{figure}[htbp]
\centering
\includegraphics[width=\linewidth]{figures/cuivres__trombone__p04_23.jpg}
\caption{Illustration (page 4)}
\end{figure}

egin{figure}[htbp]
\centering
\includegraphics[width=\linewidth]{figures/cuivres__trombone__p04_24.jpg}
\caption{Illustration (page 4)}
\end{figure}

egin{figure}[htbp]
\centering
\includegraphics[width=\linewidth]{figures/cuivres__trombone__p04_25.jpg}
\caption{Illustration (page 4)}
\end{figure}

egin{figure}[htbp]
\centering
\includegraphics[width=\linewidth]{figures/cuivres__trombone__p04_26.jpg}
\caption{Illustration (page 4)}
\end{figure}

egin{figure}[htbp]
\centering
\includegraphics[width=\linewidth]{figures/cuivres__trombone__p04_27.png}
\caption{Illustration (page 4)}
\end{figure}

egin{figure}[htbp]
\centering
\includegraphics[width=\linewidth]{figures/cuivres__trombone__p04_28.png}
\caption{Illustration (page 4)}
\end{figure}

egin{figure}[htbp]
\centering
\includegraphics[width=\linewidth]{figures/cuivres__trombone__p04_29.png}
\caption{Illustration (page 4)}
\end{figure}

egin{figure}[htbp]
\centering
\includegraphics[width=\linewidth]{figures/cuivres__trombone__p04_30.png}
\caption{Illustration (page 4)}
\end{figure}

egin{figure}[htbp]
\centering
\includegraphics[width=\linewidth]{figures/cuivres__trombone__p04_31.png}
\caption{Illustration (page 4)}
\end{figure}

egin{figure}[htbp]
\centering
\includegraphics[width=\linewidth]{figures/cuivres__trombone__p04_32.jpg}
\caption{Illustration (page 4)}
\end{figure}

egin{figure}[htbp]
\centering
\includegraphics[width=\linewidth]{figures/cuivres__trombone__p04_33.png}
\caption{Illustration (page 4)}
\end{figure}
\subsection{Cor en Fa (Horn)}
Doublures •Cor: F1 471 vs 457 Hz (Δ=14) — convergence remarquable. •Trombone: F1 471 vs 491 Hz (Δ=20) — excellente. •Basson: F1 471 vs 502 Hz (Δ=31) — excellente. Cor en Fa (Horn) SOL2020: 63 échantillons soutenus. SourceF1 (Hz)F2 (Hz)F3 (Hz)F4 (Hz)VoyelleNAccor d Backus400– 500————— Giesler250– 500———u/o— Meyer340750 (aux)1 2252 000u/o— SOL2020457 ±99685 ±66888 ±69—o (oh)63✓✓ Correction majeure: ancien 324 Hz → 457 Hz, désormais dans la plage Backus (400–500). Giesler: "u-ähnlich, hauptsächliche Abstrahlung ≤350 Hz". Meyer: pp-mf = 'u', ff = 'o'.
\par\medskip

egin{figure}[htbp]
\centering
\includegraphics[width=\linewidth]{figures/cuivres__trombone__p05_01.jpg}
\caption{Illustration (page 5)}
\end{figure}

egin{figure}[htbp]
\centering
\includegraphics[width=\linewidth]{figures/cuivres__trombone__p05_02.jpg}
\caption{Illustration (page 5)}
\end{figure}

egin{figure}[htbp]
\centering
\includegraphics[width=\linewidth]{figures/cuivres__trombone__p05_03.png}
\caption{Illustration (page 5)}
\end{figure}

egin{figure}[htbp]
\centering
\includegraphics[width=\linewidth]{figures/cuivres__trombone__p05_04.jpg}
\caption{Illustration (page 5)}
\end{figure}

egin{figure}[htbp]
\centering
\includegraphics[width=\linewidth]{figures/cuivres__trombone__p05_05.jpg}
\caption{Illustration (page 5)}
\end{figure}

egin{figure}[htbp]
\centering
\includegraphics[width=\linewidth]{figures/cuivres__trombone__p05_06.jpg}
\caption{Illustration (page 5)}
\end{figure}

egin{figure}[htbp]
\centering
\includegraphics[width=\linewidth]{figures/cuivres__trombone__p05_07.jpg}
\caption{Illustration (page 5)}
\end{figure}

egin{figure}[htbp]
\centering
\includegraphics[width=\linewidth]{figures/cuivres__trombone__p05_08.png}
\caption{Illustration (page 5)}
\end{figure}

egin{figure}[htbp]
\centering
\includegraphics[width=\linewidth]{figures/cuivres__trombone__p05_09.jpg}
\caption{Illustration (page 5)}
\end{figure}

egin{figure}[htbp]
\centering
\includegraphics[width=\linewidth]{figures/cuivres__trombone__p05_10.jpg}
\caption{Illustration (page 5)}
\end{figure}

egin{figure}[htbp]
\centering
\includegraphics[width=\linewidth]{figures/cuivres__trombone__p05_11.png}
\caption{Illustration (page 5)}
\end{figure}

egin{figure}[htbp]
\centering
\includegraphics[width=\linewidth]{figures/cuivres__trombone__p05_12.png}
\caption{Illustration (page 5)}
\end{figure}

egin{figure}[htbp]
\centering
\includegraphics[width=\linewidth]{figures/cuivres__trombone__p05_13.png}
\caption{Illustration (page 5)}
\end{figure}

egin{figure}[htbp]
\centering
\includegraphics[width=\linewidth]{figures/cuivres__trombone__p05_14.png}
\caption{Illustration (page 5)}
\end{figure}

egin{figure}[htbp]
\centering
\includegraphics[width=\linewidth]{figures/cuivres__trombone__p05_15.png}
\caption{Illustration (page 5)}
\end{figure}

egin{figure}[htbp]
\centering
\includegraphics[width=\linewidth]{figures/cuivres__trombone__p05_16.png}
\caption{Illustration (page 5)}
\end{figure}

egin{figure}[htbp]
\centering
\includegraphics[width=\linewidth]{figures/cuivres__trombone__p05_17.png}
\caption{Illustration (page 5)}
\end{figure}

egin{figure}[htbp]
\centering
\includegraphics[width=\linewidth]{figures/cuivres__trombone__p05_18.png}
\caption{Illustration (page 5)}
\end{figure}

egin{figure}[htbp]
\centering
\includegraphics[width=\linewidth]{figures/cuivres__trombone__p05_19.png}
\caption{Illustration (page 5)}
\end{figure}

egin{figure}[htbp]
\centering
\includegraphics[width=\linewidth]{figures/cuivres__trombone__p05_20.jpg}
\caption{Illustration (page 5)}
\end{figure}

egin{figure}[htbp]
\centering
\includegraphics[width=\linewidth]{figures/cuivres__trombone__p05_21.jpg}
\caption{Illustration (page 5)}
\end{figure}

egin{figure}[htbp]
\centering
\includegraphics[width=\linewidth]{figures/cuivres__trombone__p05_22.jpg}
\caption{Illustration (page 5)}
\end{figure}

egin{figure}[htbp]
\centering
\includegraphics[width=\linewidth]{figures/cuivres__trombone__p05_23.jpg}
\caption{Illustration (page 5)}
\end{figure}

egin{figure}[htbp]
\centering
\includegraphics[width=\linewidth]{figures/cuivres__trombone__p05_24.jpg}
\caption{Illustration (page 5)}
\end{figure}

egin{figure}[htbp]
\centering
\includegraphics[width=\linewidth]{figures/cuivres__trombone__p05_25.jpg}
\caption{Illustration (page 5)}
\end{figure}

egin{figure}[htbp]
\centering
\includegraphics[width=\linewidth]{figures/cuivres__trombone__p05_26.jpg}
\caption{Illustration (page 5)}
\end{figure}

egin{figure}[htbp]
\centering
\includegraphics[width=\linewidth]{figures/cuivres__trombone__p05_27.png}
\caption{Illustration (page 5)}
\end{figure}

egin{figure}[htbp]
\centering
\includegraphics[width=\linewidth]{figures/cuivres__trombone__p05_28.png}
\caption{Illustration (page 5)}
\end{figure}

egin{figure}[htbp]
\centering
\includegraphics[width=\linewidth]{figures/cuivres__trombone__p05_29.png}
\caption{Illustration (page 5)}
\end{figure}

egin{figure}[htbp]
\centering
\includegraphics[width=\linewidth]{figures/cuivres__trombone__p05_30.png}
\caption{Illustration (page 5)}
\end{figure}

egin{figure}[htbp]
\centering
\includegraphics[width=\linewidth]{figures/cuivres__trombone__p05_31.png}
\caption{Illustration (page 5)}
\end{figure}

egin{figure}[htbp]
\centering
\includegraphics[width=\linewidth]{figures/cuivres__trombone__p05_32.jpg}
\caption{Illustration (page 5)}
\end{figure}

egin{figure}[htbp]
\centering
\includegraphics[width=\linewidth]{figures/cuivres__trombone__p05_33.png}
\caption{Illustration (page 5)}
\end{figure}

\section*{SECTION II — Analyse spectrale complète}
\subsection{Trombone ténor}
Doublures •Basson: F1 457 vs 502 (Δ=45) — transition bois-cuivres classique. •Trombone: F1 457 vs 491 (Δ=33) — homogénéité cuivres. •Violoncelle: F1 457 vs 499 (Δ=42) — chant lyrique. Trombone ténor SOL2020: 78 échantillons soutenus. SourceF1 (Hz)F2 (Hz)F3 (Hz)F4 (Hz)VoyelleNAccor d Backus500————— Giesler520– 600———o— Meyer480– 600région a——o/å— SOL2020491 ±118770 ±1131 148 ±1871 454 ±210o (oh)78✓✓ Concordance parfaite: SOL (491) = Backus (500, Δ=9 Hz) = Meyer (480–600) = Giesler (520–600). Correction: anciens 334/594 Hz désormais résolus. Convergence remarquable avec Violoncelle (499 Hz), Δ=8 Hz. [Note v4 — Fp ≠ F1 spectral]: Le SOL2020 v3 (script corrigé, N=117 ordinario) révèle un F1 spectral strict à 237 Hz (/u/ grave) — la résonance fondamentale de la colonne d'air. Le Fp=491 Hz mesuré ci-dessus correspond au F2 spectral (506 Hz, Δ=15 Hz), le second mode, mais dominant perceptuellement car rayonné avec plus d'énergie. Ce phénomène explique que Backus, Meyer et Giesler mesurent tous \textasciitilde{}500 Hz: ils capturent le formant saillant, non le premier mode spectral. La Section II de ce document présente les deux niveaux d'analyse. Trombone basse (Bass Trombone) — NOUVEAU Yan\_Adds: 144 échantillons ordinario. SourceF1 (Hz)F2 (Hz)F3 (Hz)F4 (Hz)VoyelleNAccor d Backus/Meyer——————— Yan\_Adds894 ±2571 496 ±2922 073 ±3762 766 ±493a (ah)144—
\par\medskip

egin{figure}[htbp]
\centering
\includegraphics[width=\linewidth]{figures/cuivres__basson__p06_01.jpg}
\caption{Illustration (page 6)}
\end{figure}

egin{figure}[htbp]
\centering
\includegraphics[width=\linewidth]{figures/cuivres__basson__p06_02.jpg}
\caption{Illustration (page 6)}
\end{figure}

egin{figure}[htbp]
\centering
\includegraphics[width=\linewidth]{figures/cuivres__basson__p06_03.png}
\caption{Illustration (page 6)}
\end{figure}

egin{figure}[htbp]
\centering
\includegraphics[width=\linewidth]{figures/cuivres__basson__p06_04.jpg}
\caption{Illustration (page 6)}
\end{figure}

egin{figure}[htbp]
\centering
\includegraphics[width=\linewidth]{figures/cuivres__basson__p06_05.jpg}
\caption{Illustration (page 6)}
\end{figure}

egin{figure}[htbp]
\centering
\includegraphics[width=\linewidth]{figures/cuivres__basson__p06_06.jpg}
\caption{Illustration (page 6)}
\end{figure}

egin{figure}[htbp]
\centering
\includegraphics[width=\linewidth]{figures/cuivres__basson__p06_07.jpg}
\caption{Illustration (page 6)}
\end{figure}

egin{figure}[htbp]
\centering
\includegraphics[width=\linewidth]{figures/cuivres__basson__p06_08.png}
\caption{Illustration (page 6)}
\end{figure}

egin{figure}[htbp]
\centering
\includegraphics[width=\linewidth]{figures/cuivres__basson__p06_09.jpg}
\caption{Illustration (page 6)}
\end{figure}

egin{figure}[htbp]
\centering
\includegraphics[width=\linewidth]{figures/cuivres__basson__p06_10.jpg}
\caption{Illustration (page 6)}
\end{figure}

egin{figure}[htbp]
\centering
\includegraphics[width=\linewidth]{figures/cuivres__basson__p06_11.png}
\caption{Illustration (page 6)}
\end{figure}

egin{figure}[htbp]
\centering
\includegraphics[width=\linewidth]{figures/cuivres__basson__p06_12.png}
\caption{Illustration (page 6)}
\end{figure}

egin{figure}[htbp]
\centering
\includegraphics[width=\linewidth]{figures/cuivres__basson__p06_13.png}
\caption{Illustration (page 6)}
\end{figure}

egin{figure}[htbp]
\centering
\includegraphics[width=\linewidth]{figures/cuivres__basson__p06_14.png}
\caption{Illustration (page 6)}
\end{figure}

egin{figure}[htbp]
\centering
\includegraphics[width=\linewidth]{figures/cuivres__basson__p06_15.png}
\caption{Illustration (page 6)}
\end{figure}

egin{figure}[htbp]
\centering
\includegraphics[width=\linewidth]{figures/cuivres__basson__p06_16.png}
\caption{Illustration (page 6)}
\end{figure}

egin{figure}[htbp]
\centering
\includegraphics[width=\linewidth]{figures/cuivres__basson__p06_17.png}
\caption{Illustration (page 6)}
\end{figure}

egin{figure}[htbp]
\centering
\includegraphics[width=\linewidth]{figures/cuivres__basson__p06_18.png}
\caption{Illustration (page 6)}
\end{figure}

egin{figure}[htbp]
\centering
\includegraphics[width=\linewidth]{figures/cuivres__basson__p06_19.png}
\caption{Illustration (page 6)}
\end{figure}

egin{figure}[htbp]
\centering
\includegraphics[width=\linewidth]{figures/cuivres__basson__p06_20.jpg}
\caption{Illustration (page 6)}
\end{figure}

egin{figure}[htbp]
\centering
\includegraphics[width=\linewidth]{figures/cuivres__basson__p06_21.jpg}
\caption{Illustration (page 6)}
\end{figure}

egin{figure}[htbp]
\centering
\includegraphics[width=\linewidth]{figures/cuivres__basson__p06_22.jpg}
\caption{Illustration (page 6)}
\end{figure}

egin{figure}[htbp]
\centering
\includegraphics[width=\linewidth]{figures/cuivres__basson__p06_23.jpg}
\caption{Illustration (page 6)}
\end{figure}

egin{figure}[htbp]
\centering
\includegraphics[width=\linewidth]{figures/cuivres__basson__p06_24.jpg}
\caption{Illustration (page 6)}
\end{figure}

egin{figure}[htbp]
\centering
\includegraphics[width=\linewidth]{figures/cuivres__basson__p06_25.jpg}
\caption{Illustration (page 6)}
\end{figure}

egin{figure}[htbp]
\centering
\includegraphics[width=\linewidth]{figures/cuivres__basson__p06_26.jpg}
\caption{Illustration (page 6)}
\end{figure}

egin{figure}[htbp]
\centering
\includegraphics[width=\linewidth]{figures/cuivres__basson__p06_27.png}
\caption{Illustration (page 6)}
\end{figure}

egin{figure}[htbp]
\centering
\includegraphics[width=\linewidth]{figures/cuivres__basson__p06_28.png}
\caption{Illustration (page 6)}
\end{figure}

egin{figure}[htbp]
\centering
\includegraphics[width=\linewidth]{figures/cuivres__basson__p06_29.png}
\caption{Illustration (page 6)}
\end{figure}

egin{figure}[htbp]
\centering
\includegraphics[width=\linewidth]{figures/cuivres__basson__p06_30.png}
\caption{Illustration (page 6)}
\end{figure}

egin{figure}[htbp]
\centering
\includegraphics[width=\linewidth]{figures/cuivres__basson__p06_31.png}
\caption{Illustration (page 6)}
\end{figure}

egin{figure}[htbp]
\centering
\includegraphics[width=\linewidth]{figures/cuivres__basson__p06_32.jpg}
\caption{Illustration (page 6)}
\end{figure}

egin{figure}[htbp]
\centering
\includegraphics[width=\linewidth]{figures/cuivres__basson__p06_33.png}
\caption{Illustration (page 6)}
\end{figure}
Surprise: F1=894 Hz est nettement plus haut que le trombone ténor (491 Hz). Le trombone basse se situe dans la zone 'a (ah)' de puissance, pas dans 'o (oh)'. Son timbre est donc qualitativement différent — plus mordant que plein. Variation: pp=580, mf=876, ff=1 011 Hz. Doublures •Clarinette basse: F1 894 vs 909 (Δ=15) — convergence quasi-parfaite. •Clarinette Sib: F1 894 vs 1 016 (Δ=122) — bonne proximité. Trompette en Do SOL2020: 14 échantillons ff seulement (données insuffisantes — se fier à Backus/Meyer). SourceF1 (Hz)F2 (Hz)F3 (Hz)F4 (Hz)VoyelleNAccor d Backus1 000–1 500\textasciitilde{}2 000———— Giesler1 200–1 500———e— Meyer1 200–1 500———a/e— SOL2020*1 324 ±1 2483 207 ±8273 8764 369i (ee)14✗ (n
\par\medskip

egin{figure}[htbp]
\centering
\includegraphics[width=\linewidth]{figures/cuivres__trombone__p07_01.jpg}
\caption{Illustration (page 7)}
\end{figure}

egin{figure}[htbp]
\centering
\includegraphics[width=\linewidth]{figures/cuivres__trombone__p07_02.jpg}
\caption{Illustration (page 7)}
\end{figure}

egin{figure}[htbp]
\centering
\includegraphics[width=\linewidth]{figures/cuivres__trombone__p07_03.png}
\caption{Illustration (page 7)}
\end{figure}

egin{figure}[htbp]
\centering
\includegraphics[width=\linewidth]{figures/cuivres__trombone__p07_04.jpg}
\caption{Illustration (page 7)}
\end{figure}

egin{figure}[htbp]
\centering
\includegraphics[width=\linewidth]{figures/cuivres__trombone__p07_05.jpg}
\caption{Illustration (page 7)}
\end{figure}

egin{figure}[htbp]
\centering
\includegraphics[width=\linewidth]{figures/cuivres__trombone__p07_06.jpg}
\caption{Illustration (page 7)}
\end{figure}

egin{figure}[htbp]
\centering
\includegraphics[width=\linewidth]{figures/cuivres__trombone__p07_07.jpg}
\caption{Illustration (page 7)}
\end{figure}

egin{figure}[htbp]
\centering
\includegraphics[width=\linewidth]{figures/cuivres__trombone__p07_08.png}
\caption{Illustration (page 7)}
\end{figure}

egin{figure}[htbp]
\centering
\includegraphics[width=\linewidth]{figures/cuivres__trombone__p07_09.jpg}
\caption{Illustration (page 7)}
\end{figure}

egin{figure}[htbp]
\centering
\includegraphics[width=\linewidth]{figures/cuivres__trombone__p07_10.jpg}
\caption{Illustration (page 7)}
\end{figure}

egin{figure}[htbp]
\centering
\includegraphics[width=\linewidth]{figures/cuivres__trombone__p07_11.png}
\caption{Illustration (page 7)}
\end{figure}

egin{figure}[htbp]
\centering
\includegraphics[width=\linewidth]{figures/cuivres__trombone__p07_12.png}
\caption{Illustration (page 7)}
\end{figure}

egin{figure}[htbp]
\centering
\includegraphics[width=\linewidth]{figures/cuivres__trombone__p07_13.png}
\caption{Illustration (page 7)}
\end{figure}

egin{figure}[htbp]
\centering
\includegraphics[width=\linewidth]{figures/cuivres__trombone__p07_14.png}
\caption{Illustration (page 7)}
\end{figure}

egin{figure}[htbp]
\centering
\includegraphics[width=\linewidth]{figures/cuivres__trombone__p07_15.png}
\caption{Illustration (page 7)}
\end{figure}

egin{figure}[htbp]
\centering
\includegraphics[width=\linewidth]{figures/cuivres__trombone__p07_16.png}
\caption{Illustration (page 7)}
\end{figure}

egin{figure}[htbp]
\centering
\includegraphics[width=\linewidth]{figures/cuivres__trombone__p07_17.png}
\caption{Illustration (page 7)}
\end{figure}

egin{figure}[htbp]
\centering
\includegraphics[width=\linewidth]{figures/cuivres__trombone__p07_18.png}
\caption{Illustration (page 7)}
\end{figure}

egin{figure}[htbp]
\centering
\includegraphics[width=\linewidth]{figures/cuivres__trombone__p07_19.png}
\caption{Illustration (page 7)}
\end{figure}

egin{figure}[htbp]
\centering
\includegraphics[width=\linewidth]{figures/cuivres__trombone__p07_20.jpg}
\caption{Illustration (page 7)}
\end{figure}

egin{figure}[htbp]
\centering
\includegraphics[width=\linewidth]{figures/cuivres__trombone__p07_21.jpg}
\caption{Illustration (page 7)}
\end{figure}

egin{figure}[htbp]
\centering
\includegraphics[width=\linewidth]{figures/cuivres__trombone__p07_22.jpg}
\caption{Illustration (page 7)}
\end{figure}

egin{figure}[htbp]
\centering
\includegraphics[width=\linewidth]{figures/cuivres__trombone__p07_23.jpg}
\caption{Illustration (page 7)}
\end{figure}

egin{figure}[htbp]
\centering
\includegraphics[width=\linewidth]{figures/cuivres__trombone__p07_24.jpg}
\caption{Illustration (page 7)}
\end{figure}

egin{figure}[htbp]
\centering
\includegraphics[width=\linewidth]{figures/cuivres__trombone__p07_25.jpg}
\caption{Illustration (page 7)}
\end{figure}

egin{figure}[htbp]
\centering
\includegraphics[width=\linewidth]{figures/cuivres__trombone__p07_26.jpg}
\caption{Illustration (page 7)}
\end{figure}

egin{figure}[htbp]
\centering
\includegraphics[width=\linewidth]{figures/cuivres__trombone__p07_27.png}
\caption{Illustration (page 7)}
\end{figure}

egin{figure}[htbp]
\centering
\includegraphics[width=\linewidth]{figures/cuivres__trombone__p07_28.png}
\caption{Illustration (page 7)}
\end{figure}

egin{figure}[htbp]
\centering
\includegraphics[width=\linewidth]{figures/cuivres__trombone__p07_29.png}
\caption{Illustration (page 7)}
\end{figure}

egin{figure}[htbp]
\centering
\includegraphics[width=\linewidth]{figures/cuivres__trombone__p07_30.png}
\caption{Illustration (page 7)}
\end{figure}

egin{figure}[htbp]
\centering
\includegraphics[width=\linewidth]{figures/cuivres__trombone__p07_31.png}
\caption{Illustration (page 7)}
\end{figure}

egin{figure}[htbp]
\centering
\includegraphics[width=\linewidth]{figures/cuivres__trombone__p07_32.jpg}
\caption{Illustration (page 7)}
\end{figure}

egin{figure}[htbp]
\centering
\includegraphics[width=\linewidth]{figures/cuivres__trombone__p07_33.png}
\caption{Illustration (page 7)}
\end{figure}
SourceF1 (Hz)F2 (Hz)F3 (Hz)F4 (Hz)VoyelleNAccor d insuffi sant) SOL diverge (seulement 14 échantillons ff). Valeur de référence retenue: 1 000–1 500 Hz (Backus/Meyer/Giesler). Concentration des échantillons: 72 \% dans 500–1 500 Hz. [Note v4 — Fp ≠ F1 spectral]: Le SOL2020 v3 (N=96 ordinario) identifie un F1 spectral à 786 Hz (/a/) — premier mode de la colonne d'air de la trompette. Le Fp=1 324 Hz (Giesler/Meyer) correspond au F2 spectral (1 324 Hz, Δ=0 Hz), second mode nettement plus proéminent, responsable de la percée caractéristique du timbre cuivré. La trompette est l'exemple le plus frappant de dissociation F1 spectral / Fp perceptuel dans la famille des cuivres.
\par\medskip

egin{figure}[htbp]
\centering
\includegraphics[width=\linewidth]{figures/cuivres__trompette__p08_01.jpg}
\caption{Illustration (page 8)}
\end{figure}

egin{figure}[htbp]
\centering
\includegraphics[width=\linewidth]{figures/cuivres__trompette__p08_02.jpg}
\caption{Illustration (page 8)}
\end{figure}

egin{figure}[htbp]
\centering
\includegraphics[width=\linewidth]{figures/cuivres__trompette__p08_03.png}
\caption{Illustration (page 8)}
\end{figure}

egin{figure}[htbp]
\centering
\includegraphics[width=\linewidth]{figures/cuivres__trompette__p08_04.jpg}
\caption{Illustration (page 8)}
\end{figure}

egin{figure}[htbp]
\centering
\includegraphics[width=\linewidth]{figures/cuivres__trompette__p08_05.jpg}
\caption{Illustration (page 8)}
\end{figure}

egin{figure}[htbp]
\centering
\includegraphics[width=\linewidth]{figures/cuivres__trompette__p08_06.jpg}
\caption{Illustration (page 8)}
\end{figure}

egin{figure}[htbp]
\centering
\includegraphics[width=\linewidth]{figures/cuivres__trompette__p08_07.jpg}
\caption{Illustration (page 8)}
\end{figure}

egin{figure}[htbp]
\centering
\includegraphics[width=\linewidth]{figures/cuivres__trompette__p08_08.png}
\caption{Illustration (page 8)}
\end{figure}

egin{figure}[htbp]
\centering
\includegraphics[width=\linewidth]{figures/cuivres__trompette__p08_09.jpg}
\caption{Illustration (page 8)}
\end{figure}

egin{figure}[htbp]
\centering
\includegraphics[width=\linewidth]{figures/cuivres__trompette__p08_10.jpg}
\caption{Illustration (page 8)}
\end{figure}

egin{figure}[htbp]
\centering
\includegraphics[width=\linewidth]{figures/cuivres__trompette__p08_11.png}
\caption{Illustration (page 8)}
\end{figure}

egin{figure}[htbp]
\centering
\includegraphics[width=\linewidth]{figures/cuivres__trompette__p08_12.png}
\caption{Illustration (page 8)}
\end{figure}

egin{figure}[htbp]
\centering
\includegraphics[width=\linewidth]{figures/cuivres__trompette__p08_13.png}
\caption{Illustration (page 8)}
\end{figure}

egin{figure}[htbp]
\centering
\includegraphics[width=\linewidth]{figures/cuivres__trompette__p08_14.png}
\caption{Illustration (page 8)}
\end{figure}

egin{figure}[htbp]
\centering
\includegraphics[width=\linewidth]{figures/cuivres__trompette__p08_15.png}
\caption{Illustration (page 8)}
\end{figure}

egin{figure}[htbp]
\centering
\includegraphics[width=\linewidth]{figures/cuivres__trompette__p08_16.png}
\caption{Illustration (page 8)}
\end{figure}

egin{figure}[htbp]
\centering
\includegraphics[width=\linewidth]{figures/cuivres__trompette__p08_17.png}
\caption{Illustration (page 8)}
\end{figure}

egin{figure}[htbp]
\centering
\includegraphics[width=\linewidth]{figures/cuivres__trompette__p08_18.png}
\caption{Illustration (page 8)}
\end{figure}

egin{figure}[htbp]
\centering
\includegraphics[width=\linewidth]{figures/cuivres__trompette__p08_19.png}
\caption{Illustration (page 8)}
\end{figure}

egin{figure}[htbp]
\centering
\includegraphics[width=\linewidth]{figures/cuivres__trompette__p08_20.jpg}
\caption{Illustration (page 8)}
\end{figure}

egin{figure}[htbp]
\centering
\includegraphics[width=\linewidth]{figures/cuivres__trompette__p08_21.jpg}
\caption{Illustration (page 8)}
\end{figure}

egin{figure}[htbp]
\centering
\includegraphics[width=\linewidth]{figures/cuivres__trompette__p08_22.jpg}
\caption{Illustration (page 8)}
\end{figure}

egin{figure}[htbp]
\centering
\includegraphics[width=\linewidth]{figures/cuivres__trompette__p08_23.jpg}
\caption{Illustration (page 8)}
\end{figure}

egin{figure}[htbp]
\centering
\includegraphics[width=\linewidth]{figures/cuivres__trompette__p08_24.jpg}
\caption{Illustration (page 8)}
\end{figure}

egin{figure}[htbp]
\centering
\includegraphics[width=\linewidth]{figures/cuivres__trompette__p08_25.jpg}
\caption{Illustration (page 8)}
\end{figure}

egin{figure}[htbp]
\centering
\includegraphics[width=\linewidth]{figures/cuivres__trompette__p08_26.jpg}
\caption{Illustration (page 8)}
\end{figure}

egin{figure}[htbp]
\centering
\includegraphics[width=\linewidth]{figures/cuivres__trompette__p08_27.png}
\caption{Illustration (page 8)}
\end{figure}

egin{figure}[htbp]
\centering
\includegraphics[width=\linewidth]{figures/cuivres__trompette__p08_28.png}
\caption{Illustration (page 8)}
\end{figure}

egin{figure}[htbp]
\centering
\includegraphics[width=\linewidth]{figures/cuivres__trompette__p08_29.png}
\caption{Illustration (page 8)}
\end{figure}

egin{figure}[htbp]
\centering
\includegraphics[width=\linewidth]{figures/cuivres__trompette__p08_30.png}
\caption{Illustration (page 8)}
\end{figure}

egin{figure}[htbp]
\centering
\includegraphics[width=\linewidth]{figures/cuivres__trompette__p08_31.png}
\caption{Illustration (page 8)}
\end{figure}

egin{figure}[htbp]
\centering
\includegraphics[width=\linewidth]{figures/cuivres__trompette__p08_32.jpg}
\caption{Illustration (page 8)}
\end{figure}

egin{figure}[htbp]
\centering
\includegraphics[width=\linewidth]{figures/cuivres__trompette__p08_33.png}
\caption{Illustration (page 8)}
\end{figure}

\section{Les Bois}
\subsection{Piccolo — NOUVEAU}
III. Les Bois Plage formantique: 502–2 336 Hz. Les bois montrent plus de variabilité que les cuivres, reflétant leur expressivité timbrale. La famille clarinette présente un comportement acoustique particulier (dominance des harmoniques impairs). Piccolo — NOUVEAU Yan\_Adds: 111 échantillons ordinario. SourceF1 (Hz)F2 (Hz)F3 (Hz)F4 (Hz)VoyelleNAccor d Backus\textasciitilde{}3 400————— Yan\_Adds2 336 ±8523 866 ±7905 291 ±7486 931 ±759i (ee)111✓ Instrument le plus aigu spectralement. Comme la flûte, pas de formant fixe au sens strict — zone d'énergie spectrale haute et stable. Variation pp=1 963, mf=2 469, ff=2 417 Hz. Doublures •Flûte à l'octave: doublure classique, piccolo ajoute la brillance extrême. •Violon: convergence zone 2 000–3 000 Hz, renforcement aigu. Flûte
\par\medskip

egin{figure}[htbp]
\centering
\includegraphics[width=\linewidth]{figures/bois__piccolo__p09_01.jpg}
\caption{Illustration (page 9)}
\end{figure}

egin{figure}[htbp]
\centering
\includegraphics[width=\linewidth]{figures/bois__piccolo__p09_02.jpg}
\caption{Illustration (page 9)}
\end{figure}

egin{figure}[htbp]
\centering
\includegraphics[width=\linewidth]{figures/bois__piccolo__p09_03.png}
\caption{Illustration (page 9)}
\end{figure}

egin{figure}[htbp]
\centering
\includegraphics[width=\linewidth]{figures/bois__piccolo__p09_04.jpg}
\caption{Illustration (page 9)}
\end{figure}

egin{figure}[htbp]
\centering
\includegraphics[width=\linewidth]{figures/bois__piccolo__p09_05.jpg}
\caption{Illustration (page 9)}
\end{figure}

egin{figure}[htbp]
\centering
\includegraphics[width=\linewidth]{figures/bois__piccolo__p09_06.jpg}
\caption{Illustration (page 9)}
\end{figure}

egin{figure}[htbp]
\centering
\includegraphics[width=\linewidth]{figures/bois__piccolo__p09_07.jpg}
\caption{Illustration (page 9)}
\end{figure}

egin{figure}[htbp]
\centering
\includegraphics[width=\linewidth]{figures/bois__piccolo__p09_08.png}
\caption{Illustration (page 9)}
\end{figure}

egin{figure}[htbp]
\centering
\includegraphics[width=\linewidth]{figures/bois__piccolo__p09_09.jpg}
\caption{Illustration (page 9)}
\end{figure}

egin{figure}[htbp]
\centering
\includegraphics[width=\linewidth]{figures/bois__piccolo__p09_10.jpg}
\caption{Illustration (page 9)}
\end{figure}

egin{figure}[htbp]
\centering
\includegraphics[width=\linewidth]{figures/bois__piccolo__p09_11.png}
\caption{Illustration (page 9)}
\end{figure}

egin{figure}[htbp]
\centering
\includegraphics[width=\linewidth]{figures/bois__piccolo__p09_12.png}
\caption{Illustration (page 9)}
\end{figure}

egin{figure}[htbp]
\centering
\includegraphics[width=\linewidth]{figures/bois__piccolo__p09_13.png}
\caption{Illustration (page 9)}
\end{figure}

egin{figure}[htbp]
\centering
\includegraphics[width=\linewidth]{figures/bois__piccolo__p09_14.png}
\caption{Illustration (page 9)}
\end{figure}

egin{figure}[htbp]
\centering
\includegraphics[width=\linewidth]{figures/bois__piccolo__p09_15.png}
\caption{Illustration (page 9)}
\end{figure}

egin{figure}[htbp]
\centering
\includegraphics[width=\linewidth]{figures/bois__piccolo__p09_16.png}
\caption{Illustration (page 9)}
\end{figure}

egin{figure}[htbp]
\centering
\includegraphics[width=\linewidth]{figures/bois__piccolo__p09_17.png}
\caption{Illustration (page 9)}
\end{figure}

egin{figure}[htbp]
\centering
\includegraphics[width=\linewidth]{figures/bois__piccolo__p09_18.png}
\caption{Illustration (page 9)}
\end{figure}

egin{figure}[htbp]
\centering
\includegraphics[width=\linewidth]{figures/bois__piccolo__p09_19.png}
\caption{Illustration (page 9)}
\end{figure}

egin{figure}[htbp]
\centering
\includegraphics[width=\linewidth]{figures/bois__piccolo__p09_20.jpg}
\caption{Illustration (page 9)}
\end{figure}

egin{figure}[htbp]
\centering
\includegraphics[width=\linewidth]{figures/bois__piccolo__p09_21.jpg}
\caption{Illustration (page 9)}
\end{figure}

egin{figure}[htbp]
\centering
\includegraphics[width=\linewidth]{figures/bois__piccolo__p09_22.jpg}
\caption{Illustration (page 9)}
\end{figure}

egin{figure}[htbp]
\centering
\includegraphics[width=\linewidth]{figures/bois__piccolo__p09_23.jpg}
\caption{Illustration (page 9)}
\end{figure}

egin{figure}[htbp]
\centering
\includegraphics[width=\linewidth]{figures/bois__piccolo__p09_24.jpg}
\caption{Illustration (page 9)}
\end{figure}

egin{figure}[htbp]
\centering
\includegraphics[width=\linewidth]{figures/bois__piccolo__p09_25.jpg}
\caption{Illustration (page 9)}
\end{figure}

egin{figure}[htbp]
\centering
\includegraphics[width=\linewidth]{figures/bois__piccolo__p09_26.jpg}
\caption{Illustration (page 9)}
\end{figure}

egin{figure}[htbp]
\centering
\includegraphics[width=\linewidth]{figures/bois__piccolo__p09_27.png}
\caption{Illustration (page 9)}
\end{figure}

egin{figure}[htbp]
\centering
\includegraphics[width=\linewidth]{figures/bois__piccolo__p09_28.png}
\caption{Illustration (page 9)}
\end{figure}

egin{figure}[htbp]
\centering
\includegraphics[width=\linewidth]{figures/bois__piccolo__p09_29.png}
\caption{Illustration (page 9)}
\end{figure}

egin{figure}[htbp]
\centering
\includegraphics[width=\linewidth]{figures/bois__piccolo__p09_30.png}
\caption{Illustration (page 9)}
\end{figure}

egin{figure}[htbp]
\centering
\includegraphics[width=\linewidth]{figures/bois__piccolo__p09_31.png}
\caption{Illustration (page 9)}
\end{figure}

egin{figure}[htbp]
\centering
\includegraphics[width=\linewidth]{figures/bois__piccolo__p09_32.jpg}
\caption{Illustration (page 9)}
\end{figure}

egin{figure}[htbp]
\centering
\includegraphics[width=\linewidth]{figures/bois__piccolo__p09_33.png}
\caption{Illustration (page 9)}
\end{figure}
SOL2020: 41 échantillons soutenus. SourceF1 (Hz)F2 (Hz)F3 (Hz)F4 (Hz)VoyelleNAccor d Backus\textasciitilde{}1 700————— Giesler————pas de formant s fixes— MeyerPAS DE FORMA NTS———o/u (oh/oo)— SOL20201 354 ±4801 987 ±1434 210—(variable )41✓✓* La flûte NE POSSÈDE PAS de formants fixes (Meyer, Giesler). Décroissance linéaire au-dessus du fondamental. Les valeurs représentent des zones spectrales moyennes. [Note v4 — Fp ≠ F1 spectral]: Le SOL2020 v3 (N=118 ordinario) identifie un F1 spectral à 743 Hz (/a/ médium). Le Fp historique (\textasciitilde{}1 354 Hz, Backus \textasciitilde{}1 700 Hz) correspond au F2 spectral (1 163 Hz), zone de rayonnement privilégiée selon l'embouchure et le registre. La grande variabilité confirme l'observation de Giesler: la flûte module activement son spectre, ce qui en fait un instrument harmoniquement plastique mais difficile à caractériser par un formant unique. Flûte basse (Bass Flute) — NOUVEAU Yan\_Adds: 39 échantillons ordinario. SourceF1 (Hz)F2 (Hz)F3 (Hz)F4 (Hz)VoyelleNAccor d Yan\_Adds737 ±3561 502 ±4212 342 ±7163 462 ±1203å/a39— F1 environ moitié de la flûte (737 vs 1 354 Hz), cohérent avec longueur doublée. Zone å/a, timbre plus chaud. Variation: pp=608, mf=818, ff=783 Hz.
\par\medskip

egin{figure}[htbp]
\centering
\includegraphics[width=\linewidth]{figures/bois__flute__p10_01.jpg}
\caption{Illustration (page 10)}
\end{figure}

egin{figure}[htbp]
\centering
\includegraphics[width=\linewidth]{figures/bois__flute__p10_02.jpg}
\caption{Illustration (page 10)}
\end{figure}

egin{figure}[htbp]
\centering
\includegraphics[width=\linewidth]{figures/bois__flute__p10_03.png}
\caption{Illustration (page 10)}
\end{figure}

egin{figure}[htbp]
\centering
\includegraphics[width=\linewidth]{figures/bois__flute__p10_04.jpg}
\caption{Illustration (page 10)}
\end{figure}

egin{figure}[htbp]
\centering
\includegraphics[width=\linewidth]{figures/bois__flute__p10_05.jpg}
\caption{Illustration (page 10)}
\end{figure}

egin{figure}[htbp]
\centering
\includegraphics[width=\linewidth]{figures/bois__flute__p10_06.jpg}
\caption{Illustration (page 10)}
\end{figure}

egin{figure}[htbp]
\centering
\includegraphics[width=\linewidth]{figures/bois__flute__p10_07.jpg}
\caption{Illustration (page 10)}
\end{figure}

egin{figure}[htbp]
\centering
\includegraphics[width=\linewidth]{figures/bois__flute__p10_08.png}
\caption{Illustration (page 10)}
\end{figure}

egin{figure}[htbp]
\centering
\includegraphics[width=\linewidth]{figures/bois__flute__p10_09.jpg}
\caption{Illustration (page 10)}
\end{figure}

egin{figure}[htbp]
\centering
\includegraphics[width=\linewidth]{figures/bois__flute__p10_10.jpg}
\caption{Illustration (page 10)}
\end{figure}

egin{figure}[htbp]
\centering
\includegraphics[width=\linewidth]{figures/bois__flute__p10_11.png}
\caption{Illustration (page 10)}
\end{figure}

egin{figure}[htbp]
\centering
\includegraphics[width=\linewidth]{figures/bois__flute__p10_12.png}
\caption{Illustration (page 10)}
\end{figure}

egin{figure}[htbp]
\centering
\includegraphics[width=\linewidth]{figures/bois__flute__p10_13.png}
\caption{Illustration (page 10)}
\end{figure}

egin{figure}[htbp]
\centering
\includegraphics[width=\linewidth]{figures/bois__flute__p10_14.png}
\caption{Illustration (page 10)}
\end{figure}

egin{figure}[htbp]
\centering
\includegraphics[width=\linewidth]{figures/bois__flute__p10_15.png}
\caption{Illustration (page 10)}
\end{figure}

egin{figure}[htbp]
\centering
\includegraphics[width=\linewidth]{figures/bois__flute__p10_16.png}
\caption{Illustration (page 10)}
\end{figure}

egin{figure}[htbp]
\centering
\includegraphics[width=\linewidth]{figures/bois__flute__p10_17.png}
\caption{Illustration (page 10)}
\end{figure}

egin{figure}[htbp]
\centering
\includegraphics[width=\linewidth]{figures/bois__flute__p10_18.png}
\caption{Illustration (page 10)}
\end{figure}

egin{figure}[htbp]
\centering
\includegraphics[width=\linewidth]{figures/bois__flute__p10_19.png}
\caption{Illustration (page 10)}
\end{figure}

egin{figure}[htbp]
\centering
\includegraphics[width=\linewidth]{figures/bois__flute__p10_20.jpg}
\caption{Illustration (page 10)}
\end{figure}

egin{figure}[htbp]
\centering
\includegraphics[width=\linewidth]{figures/bois__flute__p10_21.jpg}
\caption{Illustration (page 10)}
\end{figure}

egin{figure}[htbp]
\centering
\includegraphics[width=\linewidth]{figures/bois__flute__p10_22.jpg}
\caption{Illustration (page 10)}
\end{figure}

egin{figure}[htbp]
\centering
\includegraphics[width=\linewidth]{figures/bois__flute__p10_23.jpg}
\caption{Illustration (page 10)}
\end{figure}

egin{figure}[htbp]
\centering
\includegraphics[width=\linewidth]{figures/bois__flute__p10_24.jpg}
\caption{Illustration (page 10)}
\end{figure}

egin{figure}[htbp]
\centering
\includegraphics[width=\linewidth]{figures/bois__flute__p10_25.jpg}
\caption{Illustration (page 10)}
\end{figure}

egin{figure}[htbp]
\centering
\includegraphics[width=\linewidth]{figures/bois__flute__p10_26.jpg}
\caption{Illustration (page 10)}
\end{figure}

egin{figure}[htbp]
\centering
\includegraphics[width=\linewidth]{figures/bois__flute__p10_27.png}
\caption{Illustration (page 10)}
\end{figure}

egin{figure}[htbp]
\centering
\includegraphics[width=\linewidth]{figures/bois__flute__p10_28.png}
\caption{Illustration (page 10)}
\end{figure}

egin{figure}[htbp]
\centering
\includegraphics[width=\linewidth]{figures/bois__flute__p10_29.png}
\caption{Illustration (page 10)}
\end{figure}

egin{figure}[htbp]
\centering
\includegraphics[width=\linewidth]{figures/bois__flute__p10_30.png}
\caption{Illustration (page 10)}
\end{figure}

egin{figure}[htbp]
\centering
\includegraphics[width=\linewidth]{figures/bois__flute__p10_31.png}
\caption{Illustration (page 10)}
\end{figure}

egin{figure}[htbp]
\centering
\includegraphics[width=\linewidth]{figures/bois__flute__p10_32.jpg}
\caption{Illustration (page 10)}
\end{figure}

egin{figure}[htbp]
\centering
\includegraphics[width=\linewidth]{figures/bois__flute__p10_33.png}
\caption{Illustration (page 10)}
\end{figure}
\subsection{Flûte contrebasse (Contrabass Flute) — NOUVEAU}
Flûte contrebasse (Contrabass Flute) — NOUVEAU Yan\_Adds: 38 échantillons ordinario. SourceF1 (Hz)F2 (Hz)F3 (Hz)F4 (Hz)VoyelleNAccor d Yan\_Adds746 ±3651 501 ±4172 283 ±5033 222 ±803å/a38— Surprise: F1 (746 Hz) quasi-identique à la flûte basse (737 Hz). Les deux convergent fortement, suggérant une zone de résonance commune malgré l'octave de différence en registre jouable. Hautbois (Oboe)
\par\medskip

egin{figure}[htbp]
\centering
\includegraphics[width=\linewidth]{figures/bois__contrebasse__p11_01.jpg}
\caption{Illustration (page 11)}
\end{figure}

egin{figure}[htbp]
\centering
\includegraphics[width=\linewidth]{figures/bois__contrebasse__p11_02.jpg}
\caption{Illustration (page 11)}
\end{figure}

egin{figure}[htbp]
\centering
\includegraphics[width=\linewidth]{figures/bois__contrebasse__p11_03.png}
\caption{Illustration (page 11)}
\end{figure}

egin{figure}[htbp]
\centering
\includegraphics[width=\linewidth]{figures/bois__contrebasse__p11_04.jpg}
\caption{Illustration (page 11)}
\end{figure}

egin{figure}[htbp]
\centering
\includegraphics[width=\linewidth]{figures/bois__contrebasse__p11_05.jpg}
\caption{Illustration (page 11)}
\end{figure}

egin{figure}[htbp]
\centering
\includegraphics[width=\linewidth]{figures/bois__contrebasse__p11_06.jpg}
\caption{Illustration (page 11)}
\end{figure}

egin{figure}[htbp]
\centering
\includegraphics[width=\linewidth]{figures/bois__contrebasse__p11_07.jpg}
\caption{Illustration (page 11)}
\end{figure}

egin{figure}[htbp]
\centering
\includegraphics[width=\linewidth]{figures/bois__contrebasse__p11_08.png}
\caption{Illustration (page 11)}
\end{figure}

egin{figure}[htbp]
\centering
\includegraphics[width=\linewidth]{figures/bois__contrebasse__p11_09.jpg}
\caption{Illustration (page 11)}
\end{figure}

egin{figure}[htbp]
\centering
\includegraphics[width=\linewidth]{figures/bois__contrebasse__p11_10.jpg}
\caption{Illustration (page 11)}
\end{figure}

egin{figure}[htbp]
\centering
\includegraphics[width=\linewidth]{figures/bois__contrebasse__p11_11.png}
\caption{Illustration (page 11)}
\end{figure}

egin{figure}[htbp]
\centering
\includegraphics[width=\linewidth]{figures/bois__contrebasse__p11_12.png}
\caption{Illustration (page 11)}
\end{figure}

egin{figure}[htbp]
\centering
\includegraphics[width=\linewidth]{figures/bois__contrebasse__p11_13.png}
\caption{Illustration (page 11)}
\end{figure}

egin{figure}[htbp]
\centering
\includegraphics[width=\linewidth]{figures/bois__contrebasse__p11_14.png}
\caption{Illustration (page 11)}
\end{figure}

egin{figure}[htbp]
\centering
\includegraphics[width=\linewidth]{figures/bois__contrebasse__p11_15.png}
\caption{Illustration (page 11)}
\end{figure}

egin{figure}[htbp]
\centering
\includegraphics[width=\linewidth]{figures/bois__contrebasse__p11_16.png}
\caption{Illustration (page 11)}
\end{figure}

egin{figure}[htbp]
\centering
\includegraphics[width=\linewidth]{figures/bois__contrebasse__p11_17.png}
\caption{Illustration (page 11)}
\end{figure}

egin{figure}[htbp]
\centering
\includegraphics[width=\linewidth]{figures/bois__contrebasse__p11_18.png}
\caption{Illustration (page 11)}
\end{figure}

egin{figure}[htbp]
\centering
\includegraphics[width=\linewidth]{figures/bois__contrebasse__p11_19.png}
\caption{Illustration (page 11)}
\end{figure}

egin{figure}[htbp]
\centering
\includegraphics[width=\linewidth]{figures/bois__contrebasse__p11_20.jpg}
\caption{Illustration (page 11)}
\end{figure}

egin{figure}[htbp]
\centering
\includegraphics[width=\linewidth]{figures/bois__contrebasse__p11_21.jpg}
\caption{Illustration (page 11)}
\end{figure}

egin{figure}[htbp]
\centering
\includegraphics[width=\linewidth]{figures/bois__contrebasse__p11_22.jpg}
\caption{Illustration (page 11)}
\end{figure}

egin{figure}[htbp]
\centering
\includegraphics[width=\linewidth]{figures/bois__contrebasse__p11_23.jpg}
\caption{Illustration (page 11)}
\end{figure}

egin{figure}[htbp]
\centering
\includegraphics[width=\linewidth]{figures/bois__contrebasse__p11_24.jpg}
\caption{Illustration (page 11)}
\end{figure}

egin{figure}[htbp]
\centering
\includegraphics[width=\linewidth]{figures/bois__contrebasse__p11_25.jpg}
\caption{Illustration (page 11)}
\end{figure}

egin{figure}[htbp]
\centering
\includegraphics[width=\linewidth]{figures/bois__contrebasse__p11_26.jpg}
\caption{Illustration (page 11)}
\end{figure}

egin{figure}[htbp]
\centering
\includegraphics[width=\linewidth]{figures/bois__contrebasse__p11_27.png}
\caption{Illustration (page 11)}
\end{figure}

egin{figure}[htbp]
\centering
\includegraphics[width=\linewidth]{figures/bois__contrebasse__p11_28.png}
\caption{Illustration (page 11)}
\end{figure}

egin{figure}[htbp]
\centering
\includegraphics[width=\linewidth]{figures/bois__contrebasse__p11_29.png}
\caption{Illustration (page 11)}
\end{figure}

egin{figure}[htbp]
\centering
\includegraphics[width=\linewidth]{figures/bois__contrebasse__p11_30.png}
\caption{Illustration (page 11)}
\end{figure}

egin{figure}[htbp]
\centering
\includegraphics[width=\linewidth]{figures/bois__contrebasse__p11_31.png}
\caption{Illustration (page 11)}
\end{figure}

egin{figure}[htbp]
\centering
\includegraphics[width=\linewidth]{figures/bois__contrebasse__p11_32.jpg}
\caption{Illustration (page 11)}
\end{figure}

egin{figure}[htbp]
\centering
\includegraphics[width=\linewidth]{figures/bois__contrebasse__p11_33.png}
\caption{Illustration (page 11)}
\end{figure}
SOL2020: 56 échantillons soutenus. SourceF1 (Hz)F2 (Hz)F3 (Hz)F4 (Hz)VoyelleNAccor d Backus1 300————— Giesler\textasciitilde{}1 100———a (ah)— Meyer\textasciitilde{}1 100\textasciitilde{}2 700\textasciitilde{}4 500—a (ah)— SOL20201 460 ±6132 362 ±14091 708—a/e749✓✓ Variabilité importante. Backus (1 300), Meyer/Giesler (\textasciitilde{}1 100), SOL (1 460): écart max 360 Hz. [Note v4 — Fp ≠ F1 spectral]: Le SOL2020 v3 (N=209 ordinario) révèle un F1 spectral à 743 Hz (/a/). Le Fp=1 460 Hz du document (et \textasciitilde{}1 260 Hz en F2 spectral) correspond à la zone de nasalité et d'intensité caractéristique du hautbois. La dissociation F1 spectral (743 Hz) / Fp (1 260–1 460 Hz) est plus marquée que chez les autres bois, ce qui explique les écarts entre Backus, Meyer et SOL2020: chaque source capture une zone spectrale différente du profil complexe du hautbois. Cor anglais (English Horn) — NOUVEAU Yan\_Adds: 128 échantillons ordinario. SourceF1 (Hz)F2 (Hz)F3 (Hz)F4 (Hz)VoyelleNAccor d Backus\textasciitilde{}930————— Meyer\textasciitilde{}600———o (oh)— Yan\_Adds1 045 ±5331 929 ±6362 838 ±8184 213 ±1425a (ah)128✓ F1=1 045 Hz se situe entre hautbois (1 460) et clarinette Sib (1 016), cohérent avec le registre intermédiaire. Backus cite \textasciitilde{}930 Hz (Δ=115 Hz, bonne concordance). Zone 'a (ah)': timbre chaud et mélancolique. Variation: pp=947, f=1 180 Hz.
\par\medskip

egin{figure}[htbp]
\centering
\includegraphics[width=\linewidth]{figures/bois__hautbois__p12_01.jpg}
\caption{Illustration (page 12)}
\end{figure}

egin{figure}[htbp]
\centering
\includegraphics[width=\linewidth]{figures/bois__hautbois__p12_02.jpg}
\caption{Illustration (page 12)}
\end{figure}

egin{figure}[htbp]
\centering
\includegraphics[width=\linewidth]{figures/bois__hautbois__p12_03.png}
\caption{Illustration (page 12)}
\end{figure}

egin{figure}[htbp]
\centering
\includegraphics[width=\linewidth]{figures/bois__hautbois__p12_04.jpg}
\caption{Illustration (page 12)}
\end{figure}

egin{figure}[htbp]
\centering
\includegraphics[width=\linewidth]{figures/bois__hautbois__p12_05.jpg}
\caption{Illustration (page 12)}
\end{figure}

egin{figure}[htbp]
\centering
\includegraphics[width=\linewidth]{figures/bois__hautbois__p12_06.jpg}
\caption{Illustration (page 12)}
\end{figure}

egin{figure}[htbp]
\centering
\includegraphics[width=\linewidth]{figures/bois__hautbois__p12_07.jpg}
\caption{Illustration (page 12)}
\end{figure}

egin{figure}[htbp]
\centering
\includegraphics[width=\linewidth]{figures/bois__hautbois__p12_08.png}
\caption{Illustration (page 12)}
\end{figure}

egin{figure}[htbp]
\centering
\includegraphics[width=\linewidth]{figures/bois__hautbois__p12_09.jpg}
\caption{Illustration (page 12)}
\end{figure}

egin{figure}[htbp]
\centering
\includegraphics[width=\linewidth]{figures/bois__hautbois__p12_10.jpg}
\caption{Illustration (page 12)}
\end{figure}

egin{figure}[htbp]
\centering
\includegraphics[width=\linewidth]{figures/bois__hautbois__p12_11.png}
\caption{Illustration (page 12)}
\end{figure}

egin{figure}[htbp]
\centering
\includegraphics[width=\linewidth]{figures/bois__hautbois__p12_12.png}
\caption{Illustration (page 12)}
\end{figure}

egin{figure}[htbp]
\centering
\includegraphics[width=\linewidth]{figures/bois__hautbois__p12_13.png}
\caption{Illustration (page 12)}
\end{figure}

egin{figure}[htbp]
\centering
\includegraphics[width=\linewidth]{figures/bois__hautbois__p12_14.png}
\caption{Illustration (page 12)}
\end{figure}

egin{figure}[htbp]
\centering
\includegraphics[width=\linewidth]{figures/bois__hautbois__p12_15.png}
\caption{Illustration (page 12)}
\end{figure}

egin{figure}[htbp]
\centering
\includegraphics[width=\linewidth]{figures/bois__hautbois__p12_16.png}
\caption{Illustration (page 12)}
\end{figure}

egin{figure}[htbp]
\centering
\includegraphics[width=\linewidth]{figures/bois__hautbois__p12_17.png}
\caption{Illustration (page 12)}
\end{figure}

egin{figure}[htbp]
\centering
\includegraphics[width=\linewidth]{figures/bois__hautbois__p12_18.png}
\caption{Illustration (page 12)}
\end{figure}

egin{figure}[htbp]
\centering
\includegraphics[width=\linewidth]{figures/bois__hautbois__p12_19.png}
\caption{Illustration (page 12)}
\end{figure}

egin{figure}[htbp]
\centering
\includegraphics[width=\linewidth]{figures/bois__hautbois__p12_20.jpg}
\caption{Illustration (page 12)}
\end{figure}

egin{figure}[htbp]
\centering
\includegraphics[width=\linewidth]{figures/bois__hautbois__p12_21.jpg}
\caption{Illustration (page 12)}
\end{figure}

egin{figure}[htbp]
\centering
\includegraphics[width=\linewidth]{figures/bois__hautbois__p12_22.jpg}
\caption{Illustration (page 12)}
\end{figure}

egin{figure}[htbp]
\centering
\includegraphics[width=\linewidth]{figures/bois__hautbois__p12_23.jpg}
\caption{Illustration (page 12)}
\end{figure}

egin{figure}[htbp]
\centering
\includegraphics[width=\linewidth]{figures/bois__hautbois__p12_24.jpg}
\caption{Illustration (page 12)}
\end{figure}

egin{figure}[htbp]
\centering
\includegraphics[width=\linewidth]{figures/bois__hautbois__p12_25.jpg}
\caption{Illustration (page 12)}
\end{figure}

egin{figure}[htbp]
\centering
\includegraphics[width=\linewidth]{figures/bois__hautbois__p12_26.jpg}
\caption{Illustration (page 12)}
\end{figure}

egin{figure}[htbp]
\centering
\includegraphics[width=\linewidth]{figures/bois__hautbois__p12_27.png}
\caption{Illustration (page 12)}
\end{figure}

egin{figure}[htbp]
\centering
\includegraphics[width=\linewidth]{figures/bois__hautbois__p12_28.png}
\caption{Illustration (page 12)}
\end{figure}

egin{figure}[htbp]
\centering
\includegraphics[width=\linewidth]{figures/bois__hautbois__p12_29.png}
\caption{Illustration (page 12)}
\end{figure}

egin{figure}[htbp]
\centering
\includegraphics[width=\linewidth]{figures/bois__hautbois__p12_30.png}
\caption{Illustration (page 12)}
\end{figure}

egin{figure}[htbp]
\centering
\includegraphics[width=\linewidth]{figures/bois__hautbois__p12_31.png}
\caption{Illustration (page 12)}
\end{figure}

egin{figure}[htbp]
\centering
\includegraphics[width=\linewidth]{figures/bois__hautbois__p12_32.jpg}
\caption{Illustration (page 12)}
\end{figure}

egin{figure}[htbp]
\centering
\includegraphics[width=\linewidth]{figures/bois__hautbois__p12_33.png}
\caption{Illustration (page 12)}
\end{figure}
Doublures •Clarinette Sib: F1 1 045 vs 1 016 (Δ=29) — convergence excellente. •Alto solo: F1 1 045 vs 369 Hz (F1 corrigé) — contraste timbral /a/ vs /o/, couleur mélancolique complémentaire. •Alto ensemble: F1 1 045 vs 1 190 (Δ=145) — bonne proximité. Clarinette en Sib SOL2020: 41 échantillons soutenus. SourceF1 (Hz)F2 (Hz)F3 (Hz)F4 (Hz)VoyelleNAccor d Backus1 500————— Giesler3 000–4 000———i-ähnlich— Meyer1 500———(variable )— SOL20201 016 ±478*1 804 ±4632 485 ±895—(variable )41\textasciitilde{}*
\par\medskip

egin{figure}[htbp]
\centering
\includegraphics[width=\linewidth]{figures/bois__alto__p13_01.jpg}
\caption{Illustration (page 13)}
\end{figure}

egin{figure}[htbp]
\centering
\includegraphics[width=\linewidth]{figures/bois__alto__p13_02.jpg}
\caption{Illustration (page 13)}
\end{figure}

egin{figure}[htbp]
\centering
\includegraphics[width=\linewidth]{figures/bois__alto__p13_03.png}
\caption{Illustration (page 13)}
\end{figure}

egin{figure}[htbp]
\centering
\includegraphics[width=\linewidth]{figures/bois__alto__p13_04.jpg}
\caption{Illustration (page 13)}
\end{figure}

egin{figure}[htbp]
\centering
\includegraphics[width=\linewidth]{figures/bois__alto__p13_05.jpg}
\caption{Illustration (page 13)}
\end{figure}

egin{figure}[htbp]
\centering
\includegraphics[width=\linewidth]{figures/bois__alto__p13_06.jpg}
\caption{Illustration (page 13)}
\end{figure}

egin{figure}[htbp]
\centering
\includegraphics[width=\linewidth]{figures/bois__alto__p13_07.jpg}
\caption{Illustration (page 13)}
\end{figure}

egin{figure}[htbp]
\centering
\includegraphics[width=\linewidth]{figures/bois__alto__p13_08.png}
\caption{Illustration (page 13)}
\end{figure}

egin{figure}[htbp]
\centering
\includegraphics[width=\linewidth]{figures/bois__alto__p13_09.jpg}
\caption{Illustration (page 13)}
\end{figure}

egin{figure}[htbp]
\centering
\includegraphics[width=\linewidth]{figures/bois__alto__p13_10.jpg}
\caption{Illustration (page 13)}
\end{figure}

egin{figure}[htbp]
\centering
\includegraphics[width=\linewidth]{figures/bois__alto__p13_11.png}
\caption{Illustration (page 13)}
\end{figure}

egin{figure}[htbp]
\centering
\includegraphics[width=\linewidth]{figures/bois__alto__p13_12.png}
\caption{Illustration (page 13)}
\end{figure}

egin{figure}[htbp]
\centering
\includegraphics[width=\linewidth]{figures/bois__alto__p13_13.png}
\caption{Illustration (page 13)}
\end{figure}

egin{figure}[htbp]
\centering
\includegraphics[width=\linewidth]{figures/bois__alto__p13_14.png}
\caption{Illustration (page 13)}
\end{figure}

egin{figure}[htbp]
\centering
\includegraphics[width=\linewidth]{figures/bois__alto__p13_15.png}
\caption{Illustration (page 13)}
\end{figure}

egin{figure}[htbp]
\centering
\includegraphics[width=\linewidth]{figures/bois__alto__p13_16.png}
\caption{Illustration (page 13)}
\end{figure}

egin{figure}[htbp]
\centering
\includegraphics[width=\linewidth]{figures/bois__alto__p13_17.png}
\caption{Illustration (page 13)}
\end{figure}

egin{figure}[htbp]
\centering
\includegraphics[width=\linewidth]{figures/bois__alto__p13_18.png}
\caption{Illustration (page 13)}
\end{figure}

egin{figure}[htbp]
\centering
\includegraphics[width=\linewidth]{figures/bois__alto__p13_19.png}
\caption{Illustration (page 13)}
\end{figure}

egin{figure}[htbp]
\centering
\includegraphics[width=\linewidth]{figures/bois__alto__p13_20.jpg}
\caption{Illustration (page 13)}
\end{figure}

egin{figure}[htbp]
\centering
\includegraphics[width=\linewidth]{figures/bois__alto__p13_21.jpg}
\caption{Illustration (page 13)}
\end{figure}

egin{figure}[htbp]
\centering
\includegraphics[width=\linewidth]{figures/bois__alto__p13_22.jpg}
\caption{Illustration (page 13)}
\end{figure}

egin{figure}[htbp]
\centering
\includegraphics[width=\linewidth]{figures/bois__alto__p13_23.jpg}
\caption{Illustration (page 13)}
\end{figure}

egin{figure}[htbp]
\centering
\includegraphics[width=\linewidth]{figures/bois__alto__p13_24.jpg}
\caption{Illustration (page 13)}
\end{figure}

egin{figure}[htbp]
\centering
\includegraphics[width=\linewidth]{figures/bois__alto__p13_25.jpg}
\caption{Illustration (page 13)}
\end{figure}

egin{figure}[htbp]
\centering
\includegraphics[width=\linewidth]{figures/bois__alto__p13_26.jpg}
\caption{Illustration (page 13)}
\end{figure}

egin{figure}[htbp]
\centering
\includegraphics[width=\linewidth]{figures/bois__alto__p13_27.png}
\caption{Illustration (page 13)}
\end{figure}

egin{figure}[htbp]
\centering
\includegraphics[width=\linewidth]{figures/bois__alto__p13_28.png}
\caption{Illustration (page 13)}
\end{figure}

egin{figure}[htbp]
\centering
\includegraphics[width=\linewidth]{figures/bois__alto__p13_29.png}
\caption{Illustration (page 13)}
\end{figure}

egin{figure}[htbp]
\centering
\includegraphics[width=\linewidth]{figures/bois__alto__p13_30.png}
\caption{Illustration (page 13)}
\end{figure}

egin{figure}[htbp]
\centering
\includegraphics[width=\linewidth]{figures/bois__alto__p13_31.png}
\caption{Illustration (page 13)}
\end{figure}

egin{figure}[htbp]
\centering
\includegraphics[width=\linewidth]{figures/bois__alto__p13_32.jpg}
\caption{Illustration (page 13)}
\end{figure}

egin{figure}[htbp]
\centering
\includegraphics[width=\linewidth]{figures/bois__alto__p13_33.png}
\caption{Illustration (page 13)}
\end{figure}
DÉCOUVERTE: PAS de formant fixe. Le pic suit la fondamentale (ratio 1.0) ou le 3e harmonique (ratio 3.0). Dominance des harmoniques IMPAIRS (tuyau cylindrique fermé). Giesler: "pas de formant fixe, partials hauts, i- ähnlich". [Note v4 — Fp ≠ F1 spectral]: Le SOL2020 v3 (N=126 ordinario) identifie un F1 spectral à 463 Hz (/o/) — premier mode du tuyau cylindrique fermé. Le Fp=1 016 Hz du document correspond approximativement au F2 spectral (1 130 Hz). La spécificité de la clarinette est acoustiquement profonde: son tuyau cylindrique fermé supprime les harmoniques pairs, rendant le F1 spectral acoustiquement sous-représenté et le premier harmonique impair dominant nettement plus saillant. C'est la raison pour laquelle Backus indique \textasciitilde{}1 500 Hz (registre aigu) et non \textasciitilde{}463 Hz. Clarinette en Mib — NOUVEAU Yan\_Adds: 45 échantillons ordinario. SourceF1 (Hz)F2 (Hz)F3 (Hz)F4 (Hz)VoyelleNAccor d Yan\_Adds1 747 ±11533 198 ±13164 489 ±14545 946 ±1634e (eh)45— F1=1 747 Hz, nettement plus haut que la Sib (1 016 Hz), cohérent avec le registre plus aigu. Même comportement (harmoniques impairs dominants) mais décalé vers le haut. Très grande variabilité (±1 153 Hz). Clarinette basse en Sib — NOUVEAU Yan\_Adds: 173 échantillons ordinario. SourceF1 (Hz)F2 (Hz)F3 (Hz)F4 (Hz)VoyelleNAccor d Yan\_Adds909 ±6022 018 ±7902 957 ±8624 030 ±1153a (ah)173— F1=909 Hz dans la zone 'a (ah)'. Convergence remarquable avec le trombone basse (894 Hz, Δ=15 Hz). Variation: mf=853, ff=1 436 Hz (forte montée en ff).
\par\medskip

egin{figure}[htbp]
\centering
\includegraphics[width=\linewidth]{figures/bois__clarinette_en_mib__p14_01.jpg}
\caption{Illustration (page 14)}
\end{figure}

egin{figure}[htbp]
\centering
\includegraphics[width=\linewidth]{figures/bois__clarinette_en_mib__p14_02.jpg}
\caption{Illustration (page 14)}
\end{figure}

egin{figure}[htbp]
\centering
\includegraphics[width=\linewidth]{figures/bois__clarinette_en_mib__p14_03.png}
\caption{Illustration (page 14)}
\end{figure}

egin{figure}[htbp]
\centering
\includegraphics[width=\linewidth]{figures/bois__clarinette_en_mib__p14_04.jpg}
\caption{Illustration (page 14)}
\end{figure}

egin{figure}[htbp]
\centering
\includegraphics[width=\linewidth]{figures/bois__clarinette_en_mib__p14_05.jpg}
\caption{Illustration (page 14)}
\end{figure}

egin{figure}[htbp]
\centering
\includegraphics[width=\linewidth]{figures/bois__clarinette_en_mib__p14_06.jpg}
\caption{Illustration (page 14)}
\end{figure}

egin{figure}[htbp]
\centering
\includegraphics[width=\linewidth]{figures/bois__clarinette_en_mib__p14_07.jpg}
\caption{Illustration (page 14)}
\end{figure}

egin{figure}[htbp]
\centering
\includegraphics[width=\linewidth]{figures/bois__clarinette_en_mib__p14_08.png}
\caption{Illustration (page 14)}
\end{figure}

egin{figure}[htbp]
\centering
\includegraphics[width=\linewidth]{figures/bois__clarinette_en_mib__p14_09.jpg}
\caption{Illustration (page 14)}
\end{figure}

egin{figure}[htbp]
\centering
\includegraphics[width=\linewidth]{figures/bois__clarinette_en_mib__p14_10.jpg}
\caption{Illustration (page 14)}
\end{figure}

egin{figure}[htbp]
\centering
\includegraphics[width=\linewidth]{figures/bois__clarinette_en_mib__p14_11.png}
\caption{Illustration (page 14)}
\end{figure}

egin{figure}[htbp]
\centering
\includegraphics[width=\linewidth]{figures/bois__clarinette_en_mib__p14_12.png}
\caption{Illustration (page 14)}
\end{figure}

egin{figure}[htbp]
\centering
\includegraphics[width=\linewidth]{figures/bois__clarinette_en_mib__p14_13.png}
\caption{Illustration (page 14)}
\end{figure}

egin{figure}[htbp]
\centering
\includegraphics[width=\linewidth]{figures/bois__clarinette_en_mib__p14_14.png}
\caption{Illustration (page 14)}
\end{figure}

egin{figure}[htbp]
\centering
\includegraphics[width=\linewidth]{figures/bois__clarinette_en_mib__p14_15.png}
\caption{Illustration (page 14)}
\end{figure}

egin{figure}[htbp]
\centering
\includegraphics[width=\linewidth]{figures/bois__clarinette_en_mib__p14_16.png}
\caption{Illustration (page 14)}
\end{figure}

egin{figure}[htbp]
\centering
\includegraphics[width=\linewidth]{figures/bois__clarinette_en_mib__p14_17.png}
\caption{Illustration (page 14)}
\end{figure}

egin{figure}[htbp]
\centering
\includegraphics[width=\linewidth]{figures/bois__clarinette_en_mib__p14_18.png}
\caption{Illustration (page 14)}
\end{figure}

egin{figure}[htbp]
\centering
\includegraphics[width=\linewidth]{figures/bois__clarinette_en_mib__p14_19.png}
\caption{Illustration (page 14)}
\end{figure}

egin{figure}[htbp]
\centering
\includegraphics[width=\linewidth]{figures/bois__clarinette_en_mib__p14_20.jpg}
\caption{Illustration (page 14)}
\end{figure}

egin{figure}[htbp]
\centering
\includegraphics[width=\linewidth]{figures/bois__clarinette_en_mib__p14_21.jpg}
\caption{Illustration (page 14)}
\end{figure}

egin{figure}[htbp]
\centering
\includegraphics[width=\linewidth]{figures/bois__clarinette_en_mib__p14_22.jpg}
\caption{Illustration (page 14)}
\end{figure}

egin{figure}[htbp]
\centering
\includegraphics[width=\linewidth]{figures/bois__clarinette_en_mib__p14_23.jpg}
\caption{Illustration (page 14)}
\end{figure}

egin{figure}[htbp]
\centering
\includegraphics[width=\linewidth]{figures/bois__clarinette_en_mib__p14_24.jpg}
\caption{Illustration (page 14)}
\end{figure}

egin{figure}[htbp]
\centering
\includegraphics[width=\linewidth]{figures/bois__clarinette_en_mib__p14_25.jpg}
\caption{Illustration (page 14)}
\end{figure}

egin{figure}[htbp]
\centering
\includegraphics[width=\linewidth]{figures/bois__clarinette_en_mib__p14_26.jpg}
\caption{Illustration (page 14)}
\end{figure}

egin{figure}[htbp]
\centering
\includegraphics[width=\linewidth]{figures/bois__clarinette_en_mib__p14_27.png}
\caption{Illustration (page 14)}
\end{figure}

egin{figure}[htbp]
\centering
\includegraphics[width=\linewidth]{figures/bois__clarinette_en_mib__p14_28.png}
\caption{Illustration (page 14)}
\end{figure}

egin{figure}[htbp]
\centering
\includegraphics[width=\linewidth]{figures/bois__clarinette_en_mib__p14_29.png}
\caption{Illustration (page 14)}
\end{figure}

egin{figure}[htbp]
\centering
\includegraphics[width=\linewidth]{figures/bois__clarinette_en_mib__p14_30.png}
\caption{Illustration (page 14)}
\end{figure}

egin{figure}[htbp]
\centering
\includegraphics[width=\linewidth]{figures/bois__clarinette_en_mib__p14_31.png}
\caption{Illustration (page 14)}
\end{figure}

egin{figure}[htbp]
\centering
\includegraphics[width=\linewidth]{figures/bois__clarinette_en_mib__p14_32.jpg}
\caption{Illustration (page 14)}
\end{figure}

egin{figure}[htbp]
\centering
\includegraphics[width=\linewidth]{figures/bois__clarinette_en_mib__p14_33.png}
\caption{Illustration (page 14)}
\end{figure}
\subsection{Clarinette contrebasse en Sib — NOUVEAU}
Doublures •Trombone basse: F1 909 vs 894 (Δ=15) — convergence quasi-parfaite. •Cor anglais: F1 909 vs 1 045 (Δ=136) — bonne proximité. •Clarinette Sib: F1 909 vs 1 016 (Δ=107) — famille clarinettes. Clarinette contrebasse en Sib — NOUVEAU Yan\_Adds: 45 échantillons ordinario. SourceF1 (Hz)F2 (Hz)F3 (Hz)F4 (Hz)VoyelleNAccor d Yan\_Adds963 ±7132 236 ±7093 071 ±7414 032 ±997a (ah)45— F1=963 Hz, proche de la clarinette basse (909 Hz, Δ=54 Hz). Même zone 'a (ah)'. Malgré l'octave de différence en registre jouable, les formants restent voisins — même phénomène que flûte basse/contrebasse. Basson (Bassoon)
\par\medskip

egin{figure}[htbp]
\centering
\includegraphics[width=\linewidth]{figures/bois__trombone_basse__p15_01.jpg}
\caption{Illustration (page 15)}
\end{figure}

egin{figure}[htbp]
\centering
\includegraphics[width=\linewidth]{figures/bois__trombone_basse__p15_02.jpg}
\caption{Illustration (page 15)}
\end{figure}

egin{figure}[htbp]
\centering
\includegraphics[width=\linewidth]{figures/bois__trombone_basse__p15_03.png}
\caption{Illustration (page 15)}
\end{figure}

egin{figure}[htbp]
\centering
\includegraphics[width=\linewidth]{figures/bois__trombone_basse__p15_04.jpg}
\caption{Illustration (page 15)}
\end{figure}

egin{figure}[htbp]
\centering
\includegraphics[width=\linewidth]{figures/bois__trombone_basse__p15_05.jpg}
\caption{Illustration (page 15)}
\end{figure}

egin{figure}[htbp]
\centering
\includegraphics[width=\linewidth]{figures/bois__trombone_basse__p15_06.jpg}
\caption{Illustration (page 15)}
\end{figure}

egin{figure}[htbp]
\centering
\includegraphics[width=\linewidth]{figures/bois__trombone_basse__p15_07.jpg}
\caption{Illustration (page 15)}
\end{figure}

egin{figure}[htbp]
\centering
\includegraphics[width=\linewidth]{figures/bois__trombone_basse__p15_08.png}
\caption{Illustration (page 15)}
\end{figure}

egin{figure}[htbp]
\centering
\includegraphics[width=\linewidth]{figures/bois__trombone_basse__p15_09.jpg}
\caption{Illustration (page 15)}
\end{figure}

egin{figure}[htbp]
\centering
\includegraphics[width=\linewidth]{figures/bois__trombone_basse__p15_10.jpg}
\caption{Illustration (page 15)}
\end{figure}

egin{figure}[htbp]
\centering
\includegraphics[width=\linewidth]{figures/bois__trombone_basse__p15_11.png}
\caption{Illustration (page 15)}
\end{figure}

egin{figure}[htbp]
\centering
\includegraphics[width=\linewidth]{figures/bois__trombone_basse__p15_12.png}
\caption{Illustration (page 15)}
\end{figure}

egin{figure}[htbp]
\centering
\includegraphics[width=\linewidth]{figures/bois__trombone_basse__p15_13.png}
\caption{Illustration (page 15)}
\end{figure}

egin{figure}[htbp]
\centering
\includegraphics[width=\linewidth]{figures/bois__trombone_basse__p15_14.png}
\caption{Illustration (page 15)}
\end{figure}

egin{figure}[htbp]
\centering
\includegraphics[width=\linewidth]{figures/bois__trombone_basse__p15_15.png}
\caption{Illustration (page 15)}
\end{figure}

egin{figure}[htbp]
\centering
\includegraphics[width=\linewidth]{figures/bois__trombone_basse__p15_16.png}
\caption{Illustration (page 15)}
\end{figure}

egin{figure}[htbp]
\centering
\includegraphics[width=\linewidth]{figures/bois__trombone_basse__p15_17.png}
\caption{Illustration (page 15)}
\end{figure}

egin{figure}[htbp]
\centering
\includegraphics[width=\linewidth]{figures/bois__trombone_basse__p15_18.png}
\caption{Illustration (page 15)}
\end{figure}

egin{figure}[htbp]
\centering
\includegraphics[width=\linewidth]{figures/bois__trombone_basse__p15_19.png}
\caption{Illustration (page 15)}
\end{figure}

egin{figure}[htbp]
\centering
\includegraphics[width=\linewidth]{figures/bois__trombone_basse__p15_20.jpg}
\caption{Illustration (page 15)}
\end{figure}

egin{figure}[htbp]
\centering
\includegraphics[width=\linewidth]{figures/bois__trombone_basse__p15_21.jpg}
\caption{Illustration (page 15)}
\end{figure}

egin{figure}[htbp]
\centering
\includegraphics[width=\linewidth]{figures/bois__trombone_basse__p15_22.jpg}
\caption{Illustration (page 15)}
\end{figure}

egin{figure}[htbp]
\centering
\includegraphics[width=\linewidth]{figures/bois__trombone_basse__p15_23.jpg}
\caption{Illustration (page 15)}
\end{figure}

egin{figure}[htbp]
\centering
\includegraphics[width=\linewidth]{figures/bois__trombone_basse__p15_24.jpg}
\caption{Illustration (page 15)}
\end{figure}

egin{figure}[htbp]
\centering
\includegraphics[width=\linewidth]{figures/bois__trombone_basse__p15_25.jpg}
\caption{Illustration (page 15)}
\end{figure}

egin{figure}[htbp]
\centering
\includegraphics[width=\linewidth]{figures/bois__trombone_basse__p15_26.jpg}
\caption{Illustration (page 15)}
\end{figure}

egin{figure}[htbp]
\centering
\includegraphics[width=\linewidth]{figures/bois__trombone_basse__p15_27.png}
\caption{Illustration (page 15)}
\end{figure}

egin{figure}[htbp]
\centering
\includegraphics[width=\linewidth]{figures/bois__trombone_basse__p15_28.png}
\caption{Illustration (page 15)}
\end{figure}

egin{figure}[htbp]
\centering
\includegraphics[width=\linewidth]{figures/bois__trombone_basse__p15_29.png}
\caption{Illustration (page 15)}
\end{figure}

egin{figure}[htbp]
\centering
\includegraphics[width=\linewidth]{figures/bois__trombone_basse__p15_30.png}
\caption{Illustration (page 15)}
\end{figure}

egin{figure}[htbp]
\centering
\includegraphics[width=\linewidth]{figures/bois__trombone_basse__p15_31.png}
\caption{Illustration (page 15)}
\end{figure}

egin{figure}[htbp]
\centering
\includegraphics[width=\linewidth]{figures/bois__trombone_basse__p15_32.jpg}
\caption{Illustration (page 15)}
\end{figure}

egin{figure}[htbp]
\centering
\includegraphics[width=\linewidth]{figures/bois__trombone_basse__p15_33.png}
\caption{Illustration (page 15)}
\end{figure}
SOL2020: 97 échantillons soutenus. SourceF1 (Hz)F2 (Hz)F3 (Hz)F4 (Hz)VoyelleNAccor d Backus440– 500————— Giesler500———o— Meyer\textasciitilde{}500———o (oh)— SOL2020502 ±142714 ±2481 085 ±240—o (oh)97✓✓✓ CORRECTION MAJEURE: ancien 317 Hz → 502 Hz. Accord unanime des 4 sources (Giesler 500 = Backus 440– 500 = Meyer \textasciitilde{}500 = SOL 502). 'Timbres frères' avec violoncelle (Δ=3 Hz) selon Meyer. Contrebasson — NOUVEAU Yan\_Adds: 110 échantillons non-vibrato. SourceF1 (Hz)F2 (Hz)F3 (Hz)F4 (Hz)VoyelleNAccor d Backus\textasciitilde{}250– 440————— Yan\_Adds478 ±951 273 ±3522 066 ±4842 912 ±680o (oh)110✓ DÉCOUVERTE MAJEURE: F1=478 Hz se place en plein cœur du cluster 450–502 Hz. Le contrebasson rejoint le cor (457), le tuba contrebasse (471), le trombone (491), le violoncelle (499) et le basson (502). Très faible variabilité (±95 Hz). Stabilité remarquable entre dynamiques: ff=469, mf=487, p=480 Hz.
\par\medskip

egin{figure}[htbp]
\centering
\includegraphics[width=\linewidth]{figures/bois__violon__p16_01.jpg}
\caption{Illustration (page 16)}
\end{figure}

egin{figure}[htbp]
\centering
\includegraphics[width=\linewidth]{figures/bois__violon__p16_02.jpg}
\caption{Illustration (page 16)}
\end{figure}

egin{figure}[htbp]
\centering
\includegraphics[width=\linewidth]{figures/bois__violon__p16_03.png}
\caption{Illustration (page 16)}
\end{figure}

egin{figure}[htbp]
\centering
\includegraphics[width=\linewidth]{figures/bois__violon__p16_04.jpg}
\caption{Illustration (page 16)}
\end{figure}

egin{figure}[htbp]
\centering
\includegraphics[width=\linewidth]{figures/bois__violon__p16_05.jpg}
\caption{Illustration (page 16)}
\end{figure}

egin{figure}[htbp]
\centering
\includegraphics[width=\linewidth]{figures/bois__violon__p16_06.jpg}
\caption{Illustration (page 16)}
\end{figure}

egin{figure}[htbp]
\centering
\includegraphics[width=\linewidth]{figures/bois__violon__p16_07.jpg}
\caption{Illustration (page 16)}
\end{figure}

egin{figure}[htbp]
\centering
\includegraphics[width=\linewidth]{figures/bois__violon__p16_08.png}
\caption{Illustration (page 16)}
\end{figure}

egin{figure}[htbp]
\centering
\includegraphics[width=\linewidth]{figures/bois__violon__p16_09.jpg}
\caption{Illustration (page 16)}
\end{figure}

egin{figure}[htbp]
\centering
\includegraphics[width=\linewidth]{figures/bois__violon__p16_10.jpg}
\caption{Illustration (page 16)}
\end{figure}

egin{figure}[htbp]
\centering
\includegraphics[width=\linewidth]{figures/bois__violon__p16_11.png}
\caption{Illustration (page 16)}
\end{figure}

egin{figure}[htbp]
\centering
\includegraphics[width=\linewidth]{figures/bois__violon__p16_12.png}
\caption{Illustration (page 16)}
\end{figure}

egin{figure}[htbp]
\centering
\includegraphics[width=\linewidth]{figures/bois__violon__p16_13.png}
\caption{Illustration (page 16)}
\end{figure}

egin{figure}[htbp]
\centering
\includegraphics[width=\linewidth]{figures/bois__violon__p16_14.png}
\caption{Illustration (page 16)}
\end{figure}

egin{figure}[htbp]
\centering
\includegraphics[width=\linewidth]{figures/bois__violon__p16_15.png}
\caption{Illustration (page 16)}
\end{figure}

egin{figure}[htbp]
\centering
\includegraphics[width=\linewidth]{figures/bois__violon__p16_16.png}
\caption{Illustration (page 16)}
\end{figure}

egin{figure}[htbp]
\centering
\includegraphics[width=\linewidth]{figures/bois__violon__p16_17.png}
\caption{Illustration (page 16)}
\end{figure}

egin{figure}[htbp]
\centering
\includegraphics[width=\linewidth]{figures/bois__violon__p16_18.png}
\caption{Illustration (page 16)}
\end{figure}

egin{figure}[htbp]
\centering
\includegraphics[width=\linewidth]{figures/bois__violon__p16_19.png}
\caption{Illustration (page 16)}
\end{figure}

egin{figure}[htbp]
\centering
\includegraphics[width=\linewidth]{figures/bois__violon__p16_20.jpg}
\caption{Illustration (page 16)}
\end{figure}

egin{figure}[htbp]
\centering
\includegraphics[width=\linewidth]{figures/bois__violon__p16_21.jpg}
\caption{Illustration (page 16)}
\end{figure}

egin{figure}[htbp]
\centering
\includegraphics[width=\linewidth]{figures/bois__violon__p16_22.jpg}
\caption{Illustration (page 16)}
\end{figure}

egin{figure}[htbp]
\centering
\includegraphics[width=\linewidth]{figures/bois__violon__p16_23.jpg}
\caption{Illustration (page 16)}
\end{figure}

egin{figure}[htbp]
\centering
\includegraphics[width=\linewidth]{figures/bois__violon__p16_24.jpg}
\caption{Illustration (page 16)}
\end{figure}

egin{figure}[htbp]
\centering
\includegraphics[width=\linewidth]{figures/bois__violon__p16_25.jpg}
\caption{Illustration (page 16)}
\end{figure}

egin{figure}[htbp]
\centering
\includegraphics[width=\linewidth]{figures/bois__violon__p16_26.jpg}
\caption{Illustration (page 16)}
\end{figure}

egin{figure}[htbp]
\centering
\includegraphics[width=\linewidth]{figures/bois__violon__p16_27.png}
\caption{Illustration (page 16)}
\end{figure}

egin{figure}[htbp]
\centering
\includegraphics[width=\linewidth]{figures/bois__violon__p16_28.png}
\caption{Illustration (page 16)}
\end{figure}

egin{figure}[htbp]
\centering
\includegraphics[width=\linewidth]{figures/bois__violon__p16_29.png}
\caption{Illustration (page 16)}
\end{figure}

egin{figure}[htbp]
\centering
\includegraphics[width=\linewidth]{figures/bois__violon__p16_30.png}
\caption{Illustration (page 16)}
\end{figure}

egin{figure}[htbp]
\centering
\includegraphics[width=\linewidth]{figures/bois__violon__p16_31.png}
\caption{Illustration (page 16)}
\end{figure}

egin{figure}[htbp]
\centering
\includegraphics[width=\linewidth]{figures/bois__violon__p16_32.jpg}
\caption{Illustration (page 16)}
\end{figure}

egin{figure}[htbp]
\centering
\includegraphics[width=\linewidth]{figures/bois__violon__p16_33.png}
\caption{Illustration (page 16)}
\end{figure}
Doublures •Cor: F1 478 vs 457 (Δ=21) — convergence excellente. •Trombone: F1 478 vs 491 (Δ=13) — quasi-parfaite. •Basson: F1 478 vs 502 (Δ=24) — famille bassons confirmée. •Violoncelle: F1 478 vs 499 (Δ=21) — fusion graves bois-cordes.
\par\medskip

egin{figure}[htbp]
\centering
\includegraphics[width=\linewidth]{figures/bois__trombone__p17_01.jpg}
\caption{Illustration (page 17)}
\end{figure}

egin{figure}[htbp]
\centering
\includegraphics[width=\linewidth]{figures/bois__trombone__p17_02.jpg}
\caption{Illustration (page 17)}
\end{figure}

egin{figure}[htbp]
\centering
\includegraphics[width=\linewidth]{figures/bois__trombone__p17_03.png}
\caption{Illustration (page 17)}
\end{figure}

egin{figure}[htbp]
\centering
\includegraphics[width=\linewidth]{figures/bois__trombone__p17_04.jpg}
\caption{Illustration (page 17)}
\end{figure}

egin{figure}[htbp]
\centering
\includegraphics[width=\linewidth]{figures/bois__trombone__p17_05.jpg}
\caption{Illustration (page 17)}
\end{figure}

egin{figure}[htbp]
\centering
\includegraphics[width=\linewidth]{figures/bois__trombone__p17_06.jpg}
\caption{Illustration (page 17)}
\end{figure}

egin{figure}[htbp]
\centering
\includegraphics[width=\linewidth]{figures/bois__trombone__p17_07.jpg}
\caption{Illustration (page 17)}
\end{figure}

egin{figure}[htbp]
\centering
\includegraphics[width=\linewidth]{figures/bois__trombone__p17_08.png}
\caption{Illustration (page 17)}
\end{figure}

egin{figure}[htbp]
\centering
\includegraphics[width=\linewidth]{figures/bois__trombone__p17_09.jpg}
\caption{Illustration (page 17)}
\end{figure}

egin{figure}[htbp]
\centering
\includegraphics[width=\linewidth]{figures/bois__trombone__p17_10.jpg}
\caption{Illustration (page 17)}
\end{figure}

egin{figure}[htbp]
\centering
\includegraphics[width=\linewidth]{figures/bois__trombone__p17_11.png}
\caption{Illustration (page 17)}
\end{figure}

egin{figure}[htbp]
\centering
\includegraphics[width=\linewidth]{figures/bois__trombone__p17_12.png}
\caption{Illustration (page 17)}
\end{figure}

egin{figure}[htbp]
\centering
\includegraphics[width=\linewidth]{figures/bois__trombone__p17_13.png}
\caption{Illustration (page 17)}
\end{figure}

egin{figure}[htbp]
\centering
\includegraphics[width=\linewidth]{figures/bois__trombone__p17_14.png}
\caption{Illustration (page 17)}
\end{figure}

egin{figure}[htbp]
\centering
\includegraphics[width=\linewidth]{figures/bois__trombone__p17_15.png}
\caption{Illustration (page 17)}
\end{figure}

egin{figure}[htbp]
\centering
\includegraphics[width=\linewidth]{figures/bois__trombone__p17_16.png}
\caption{Illustration (page 17)}
\end{figure}

egin{figure}[htbp]
\centering
\includegraphics[width=\linewidth]{figures/bois__trombone__p17_17.png}
\caption{Illustration (page 17)}
\end{figure}

egin{figure}[htbp]
\centering
\includegraphics[width=\linewidth]{figures/bois__trombone__p17_18.png}
\caption{Illustration (page 17)}
\end{figure}

egin{figure}[htbp]
\centering
\includegraphics[width=\linewidth]{figures/bois__trombone__p17_19.png}
\caption{Illustration (page 17)}
\end{figure}

egin{figure}[htbp]
\centering
\includegraphics[width=\linewidth]{figures/bois__trombone__p17_20.jpg}
\caption{Illustration (page 17)}
\end{figure}

egin{figure}[htbp]
\centering
\includegraphics[width=\linewidth]{figures/bois__trombone__p17_21.jpg}
\caption{Illustration (page 17)}
\end{figure}

egin{figure}[htbp]
\centering
\includegraphics[width=\linewidth]{figures/bois__trombone__p17_22.jpg}
\caption{Illustration (page 17)}
\end{figure}

egin{figure}[htbp]
\centering
\includegraphics[width=\linewidth]{figures/bois__trombone__p17_23.jpg}
\caption{Illustration (page 17)}
\end{figure}

egin{figure}[htbp]
\centering
\includegraphics[width=\linewidth]{figures/bois__trombone__p17_24.jpg}
\caption{Illustration (page 17)}
\end{figure}

egin{figure}[htbp]
\centering
\includegraphics[width=\linewidth]{figures/bois__trombone__p17_25.jpg}
\caption{Illustration (page 17)}
\end{figure}

egin{figure}[htbp]
\centering
\includegraphics[width=\linewidth]{figures/bois__trombone__p17_26.jpg}
\caption{Illustration (page 17)}
\end{figure}

egin{figure}[htbp]
\centering
\includegraphics[width=\linewidth]{figures/bois__trombone__p17_27.png}
\caption{Illustration (page 17)}
\end{figure}

egin{figure}[htbp]
\centering
\includegraphics[width=\linewidth]{figures/bois__trombone__p17_28.png}
\caption{Illustration (page 17)}
\end{figure}

egin{figure}[htbp]
\centering
\includegraphics[width=\linewidth]{figures/bois__trombone__p17_29.png}
\caption{Illustration (page 17)}
\end{figure}

egin{figure}[htbp]
\centering
\includegraphics[width=\linewidth]{figures/bois__trombone__p17_30.png}
\caption{Illustration (page 17)}
\end{figure}

egin{figure}[htbp]
\centering
\includegraphics[width=\linewidth]{figures/bois__trombone__p17_31.png}
\caption{Illustration (page 17)}
\end{figure}

egin{figure}[htbp]
\centering
\includegraphics[width=\linewidth]{figures/bois__trombone__p17_32.jpg}
\caption{Illustration (page 17)}
\end{figure}

egin{figure}[htbp]
\centering
\includegraphics[width=\linewidth]{figures/bois__trombone__p17_33.png}
\caption{Illustration (page 17)}
\end{figure}

\section{Les Saxophones}
IV. Les Saxophones Plage formantique: 374–1 217 Hz (voyelles /o/ → /a/). Données SOL2020 uniquement. Caractéristique: progression régulière selon la taille de l'instrument, couvrant presque toute la plage /o/ → /a/. Point remarquable — Saxophone Ténor dans le cluster de convergence Giesler (1985) écrit explicitement à propos du saxophone ténor: "Formantähnliches Maximum bei etwa 400-600 Hz mit o-ähnlicher Vokalfärbung... (vergl. z.B. Fagott, Horn, Posaune)" — Giesler cite EXPLICITEMENT Basson, Cor et Trombone comme partageant le même formant. Le saxophone ténor est donc un membre du cluster de convergence. InstrumentF1 SOL2020VoyelleValidation Rôle orchestral Saxophone Baryton 374 Hzo (oh)SOL2020 uniquementLe plus grave, proche Tuba/Contrebasson Saxophone Ténor 438 Hzo (oh) — CLUSTER!Giesler 400–600 Hz ✓✓MEMBRE DU CLUSTER — rejoint Cor, Trombone, Basson Saxophone Alto 829 Hzo/a (transition)SOL2020 uniquementMédium, transition vers zone /a/ Saxophone Soprano 1 217 Hza (ah)SOL2020 uniquementLe plus aigu, proche Hautbois/Trompette La famille saxophone couvre presque toute la plage /o/ → /a/ avec une progression régulière. Le Ténor Sax rejoint le cluster de convergence, confirmant l'universalité acoustique de cette zone. Ces instruments n'apparaissent pas dans Giesler/Backus/Meyer (ouvrages antérieurs à leur adoption orchestrale généralisée), mais leurs valeurs suivent les progressions attendues selon la taille/registre. [Note v4 — Fp ≠ F1 spectral]: Le SOL2020 v3 (N=99 ordinario) révèle un F1 spectral à 398 Hz (/o/) pour le saxophone alto — premier mode du tuyau conique. Le Fp=829 Hz du document correspond exactement au F2 spectral (829 Hz, Δ=0 Hz). La conicité du saxophone (contrairement à la clarinette cylindrique) lui permet de rayonner tous les harmoniques, mais l'énergie se concentre sur le deuxième mode — d'où la couleur /a/ saillante et immédiatement reconnaissable du saxophone alto. Le F1 spectral à 398 Hz place la résonance fondamentale dans le cluster de convergence /o/, un lien acoustique inattendu avec le cor et le basson.
\par\medskip

egin{figure}[htbp]
\centering
\includegraphics[width=\linewidth]{figures/saxophones__saxophone_tenor__p18_01.jpg}
\caption{Illustration (page 18)}
\end{figure}

egin{figure}[htbp]
\centering
\includegraphics[width=\linewidth]{figures/saxophones__saxophone_tenor__p18_02.jpg}
\caption{Illustration (page 18)}
\end{figure}

egin{figure}[htbp]
\centering
\includegraphics[width=\linewidth]{figures/saxophones__saxophone_tenor__p18_03.png}
\caption{Illustration (page 18)}
\end{figure}

egin{figure}[htbp]
\centering
\includegraphics[width=\linewidth]{figures/saxophones__saxophone_tenor__p18_04.jpg}
\caption{Illustration (page 18)}
\end{figure}

egin{figure}[htbp]
\centering
\includegraphics[width=\linewidth]{figures/saxophones__saxophone_tenor__p18_05.jpg}
\caption{Illustration (page 18)}
\end{figure}

egin{figure}[htbp]
\centering
\includegraphics[width=\linewidth]{figures/saxophones__saxophone_tenor__p18_06.jpg}
\caption{Illustration (page 18)}
\end{figure}

egin{figure}[htbp]
\centering
\includegraphics[width=\linewidth]{figures/saxophones__saxophone_tenor__p18_07.jpg}
\caption{Illustration (page 18)}
\end{figure}

egin{figure}[htbp]
\centering
\includegraphics[width=\linewidth]{figures/saxophones__saxophone_tenor__p18_08.png}
\caption{Illustration (page 18)}
\end{figure}

egin{figure}[htbp]
\centering
\includegraphics[width=\linewidth]{figures/saxophones__saxophone_tenor__p18_09.jpg}
\caption{Illustration (page 18)}
\end{figure}

egin{figure}[htbp]
\centering
\includegraphics[width=\linewidth]{figures/saxophones__saxophone_tenor__p18_10.jpg}
\caption{Illustration (page 18)}
\end{figure}

egin{figure}[htbp]
\centering
\includegraphics[width=\linewidth]{figures/saxophones__saxophone_tenor__p18_11.png}
\caption{Illustration (page 18)}
\end{figure}

egin{figure}[htbp]
\centering
\includegraphics[width=\linewidth]{figures/saxophones__saxophone_tenor__p18_12.png}
\caption{Illustration (page 18)}
\end{figure}

egin{figure}[htbp]
\centering
\includegraphics[width=\linewidth]{figures/saxophones__saxophone_tenor__p18_13.png}
\caption{Illustration (page 18)}
\end{figure}

egin{figure}[htbp]
\centering
\includegraphics[width=\linewidth]{figures/saxophones__saxophone_tenor__p18_14.png}
\caption{Illustration (page 18)}
\end{figure}

egin{figure}[htbp]
\centering
\includegraphics[width=\linewidth]{figures/saxophones__saxophone_tenor__p18_15.png}
\caption{Illustration (page 18)}
\end{figure}

egin{figure}[htbp]
\centering
\includegraphics[width=\linewidth]{figures/saxophones__saxophone_tenor__p18_16.png}
\caption{Illustration (page 18)}
\end{figure}

egin{figure}[htbp]
\centering
\includegraphics[width=\linewidth]{figures/saxophones__saxophone_tenor__p18_17.png}
\caption{Illustration (page 18)}
\end{figure}

egin{figure}[htbp]
\centering
\includegraphics[width=\linewidth]{figures/saxophones__saxophone_tenor__p18_18.png}
\caption{Illustration (page 18)}
\end{figure}

egin{figure}[htbp]
\centering
\includegraphics[width=\linewidth]{figures/saxophones__saxophone_tenor__p18_19.png}
\caption{Illustration (page 18)}
\end{figure}

egin{figure}[htbp]
\centering
\includegraphics[width=\linewidth]{figures/saxophones__saxophone_tenor__p18_20.jpg}
\caption{Illustration (page 18)}
\end{figure}

egin{figure}[htbp]
\centering
\includegraphics[width=\linewidth]{figures/saxophones__saxophone_tenor__p18_21.jpg}
\caption{Illustration (page 18)}
\end{figure}

egin{figure}[htbp]
\centering
\includegraphics[width=\linewidth]{figures/saxophones__saxophone_tenor__p18_22.jpg}
\caption{Illustration (page 18)}
\end{figure}

egin{figure}[htbp]
\centering
\includegraphics[width=\linewidth]{figures/saxophones__saxophone_tenor__p18_23.jpg}
\caption{Illustration (page 18)}
\end{figure}

egin{figure}[htbp]
\centering
\includegraphics[width=\linewidth]{figures/saxophones__saxophone_tenor__p18_24.jpg}
\caption{Illustration (page 18)}
\end{figure}

egin{figure}[htbp]
\centering
\includegraphics[width=\linewidth]{figures/saxophones__saxophone_tenor__p18_25.jpg}
\caption{Illustration (page 18)}
\end{figure}

egin{figure}[htbp]
\centering
\includegraphics[width=\linewidth]{figures/saxophones__saxophone_tenor__p18_26.jpg}
\caption{Illustration (page 18)}
\end{figure}

egin{figure}[htbp]
\centering
\includegraphics[width=\linewidth]{figures/saxophones__saxophone_tenor__p18_27.png}
\caption{Illustration (page 18)}
\end{figure}

egin{figure}[htbp]
\centering
\includegraphics[width=\linewidth]{figures/saxophones__saxophone_tenor__p18_28.png}
\caption{Illustration (page 18)}
\end{figure}

egin{figure}[htbp]
\centering
\includegraphics[width=\linewidth]{figures/saxophones__saxophone_tenor__p18_29.png}
\caption{Illustration (page 18)}
\end{figure}

egin{figure}[htbp]
\centering
\includegraphics[width=\linewidth]{figures/saxophones__saxophone_tenor__p18_30.png}
\caption{Illustration (page 18)}
\end{figure}

egin{figure}[htbp]
\centering
\includegraphics[width=\linewidth]{figures/saxophones__saxophone_tenor__p18_31.png}
\caption{Illustration (page 18)}
\end{figure}

egin{figure}[htbp]
\centering
\includegraphics[width=\linewidth]{figures/saxophones__saxophone_tenor__p18_32.jpg}
\caption{Illustration (page 18)}
\end{figure}

egin{figure}[htbp]
\centering
\includegraphics[width=\linewidth]{figures/saxophones__saxophone_tenor__p18_33.png}
\caption{Illustration (page 18)}
\end{figure}

\section{Les Cordes Solistes}
\subsection{violon et alto.}
V. Les Cordes Solistes Données SOL2020 (instruments individuels). Écarts-types élevés reflétant l'expressivité timbrale des cordes. Le 'bridge hill' (2 000–3 000 Hz) est la caractéristique spectrale la plus stable pour violon et alto. Violon Enveloppe spectrale (ordinario, N=679): F1 = 450  Hz (caisse); bridge hill = résonance du chevalet, zone F3–F5 (1  500– 4 000 Hz). SourceF1 (Hz)F2 (Hz)F3 (Hz)F4 (Hz)VoyelleNAccor d Backus\textasciitilde{}2 000– 3 000————— Giesler1 000–1 200———a (ah)— Meyer800–1 200———a (ah)— SOL20201 034 ±1 3233 447 ±10353 9304 005(variable )161✓✓ SOL F1 dans zone 'bridge hill'. Giesler (1 000–1 200), Meyer (\textasciitilde{}1 100), SOL (1 034): excellent accord (Δ=66 Hz). Extrême variabilité dynamique: mf=3 868, ff=2 130 Hz. Distribution: 19 \% graves, 33 \% médiums, 49 \% aigus. [Note v4 — Fp ≠ F1 spectral]: Le SOL2020 v3 (N=284 ordinario) révèle un F1 spectral à 506 Hz (/o/) — résonance grave de la caisse du violon. Le Fp=1 034 Hz du document correspond au F2 spectral (1 518 Hz, avec une zone intermédiaire 1 000–1 200 Hz validée par Giesler et Meyer). Cette dissociation est structurelle: la caisse grave vibre à \textasciitilde{}450 Hz mais le chevalet — résonateur mécanique — amplifie sélectivement la zone 1 000–3 000 Hz (bridge hill), créant le formant perceptuellement dominant. En Section II, le F1 spectral à 506 Hz place le violon dans le cluster de convergence /o/ aux côtés du cor et du basson. Alto
\par\medskip

egin{figure}[htbp]
\centering
\includegraphics[width=\linewidth]{figures/cordes_solistes__violon__p19_01.jpg}
\caption{Illustration (page 19)}
\end{figure}

egin{figure}[htbp]
\centering
\includegraphics[width=\linewidth]{figures/cordes_solistes__violon__p19_02.jpg}
\caption{Illustration (page 19)}
\end{figure}

egin{figure}[htbp]
\centering
\includegraphics[width=\linewidth]{figures/cordes_solistes__violon__p19_03.png}
\caption{Illustration (page 19)}
\end{figure}

egin{figure}[htbp]
\centering
\includegraphics[width=\linewidth]{figures/cordes_solistes__violon__p19_04.jpg}
\caption{Illustration (page 19)}
\end{figure}

egin{figure}[htbp]
\centering
\includegraphics[width=\linewidth]{figures/cordes_solistes__violon__p19_05.jpg}
\caption{Illustration (page 19)}
\end{figure}

egin{figure}[htbp]
\centering
\includegraphics[width=\linewidth]{figures/cordes_solistes__violon__p19_06.jpg}
\caption{Illustration (page 19)}
\end{figure}

egin{figure}[htbp]
\centering
\includegraphics[width=\linewidth]{figures/cordes_solistes__violon__p19_07.jpg}
\caption{Illustration (page 19)}
\end{figure}

egin{figure}[htbp]
\centering
\includegraphics[width=\linewidth]{figures/cordes_solistes__violon__p19_08.png}
\caption{Illustration (page 19)}
\end{figure}

egin{figure}[htbp]
\centering
\includegraphics[width=\linewidth]{figures/cordes_solistes__violon__p19_09.jpg}
\caption{Illustration (page 19)}
\end{figure}

egin{figure}[htbp]
\centering
\includegraphics[width=\linewidth]{figures/cordes_solistes__violon__p19_10.jpg}
\caption{Illustration (page 19)}
\end{figure}

egin{figure}[htbp]
\centering
\includegraphics[width=\linewidth]{figures/cordes_solistes__violon__p19_11.png}
\caption{Illustration (page 19)}
\end{figure}

egin{figure}[htbp]
\centering
\includegraphics[width=\linewidth]{figures/cordes_solistes__violon__p19_12.png}
\caption{Illustration (page 19)}
\end{figure}

egin{figure}[htbp]
\centering
\includegraphics[width=\linewidth]{figures/cordes_solistes__violon__p19_13.png}
\caption{Illustration (page 19)}
\end{figure}

egin{figure}[htbp]
\centering
\includegraphics[width=\linewidth]{figures/cordes_solistes__violon__p19_14.png}
\caption{Illustration (page 19)}
\end{figure}

egin{figure}[htbp]
\centering
\includegraphics[width=\linewidth]{figures/cordes_solistes__violon__p19_15.png}
\caption{Illustration (page 19)}
\end{figure}

egin{figure}[htbp]
\centering
\includegraphics[width=\linewidth]{figures/cordes_solistes__violon__p19_16.png}
\caption{Illustration (page 19)}
\end{figure}

egin{figure}[htbp]
\centering
\includegraphics[width=\linewidth]{figures/cordes_solistes__violon__p19_17.png}
\caption{Illustration (page 19)}
\end{figure}

egin{figure}[htbp]
\centering
\includegraphics[width=\linewidth]{figures/cordes_solistes__violon__p19_18.png}
\caption{Illustration (page 19)}
\end{figure}

egin{figure}[htbp]
\centering
\includegraphics[width=\linewidth]{figures/cordes_solistes__violon__p19_19.png}
\caption{Illustration (page 19)}
\end{figure}

egin{figure}[htbp]
\centering
\includegraphics[width=\linewidth]{figures/cordes_solistes__violon__p19_20.jpg}
\caption{Illustration (page 19)}
\end{figure}

egin{figure}[htbp]
\centering
\includegraphics[width=\linewidth]{figures/cordes_solistes__violon__p19_21.jpg}
\caption{Illustration (page 19)}
\end{figure}

egin{figure}[htbp]
\centering
\includegraphics[width=\linewidth]{figures/cordes_solistes__violon__p19_22.jpg}
\caption{Illustration (page 19)}
\end{figure}

egin{figure}[htbp]
\centering
\includegraphics[width=\linewidth]{figures/cordes_solistes__violon__p19_23.jpg}
\caption{Illustration (page 19)}
\end{figure}

egin{figure}[htbp]
\centering
\includegraphics[width=\linewidth]{figures/cordes_solistes__violon__p19_24.jpg}
\caption{Illustration (page 19)}
\end{figure}

egin{figure}[htbp]
\centering
\includegraphics[width=\linewidth]{figures/cordes_solistes__violon__p19_25.jpg}
\caption{Illustration (page 19)}
\end{figure}

egin{figure}[htbp]
\centering
\includegraphics[width=\linewidth]{figures/cordes_solistes__violon__p19_26.jpg}
\caption{Illustration (page 19)}
\end{figure}

egin{figure}[htbp]
\centering
\includegraphics[width=\linewidth]{figures/cordes_solistes__violon__p19_27.png}
\caption{Illustration (page 19)}
\end{figure}

egin{figure}[htbp]
\centering
\includegraphics[width=\linewidth]{figures/cordes_solistes__violon__p19_28.png}
\caption{Illustration (page 19)}
\end{figure}

egin{figure}[htbp]
\centering
\includegraphics[width=\linewidth]{figures/cordes_solistes__violon__p19_29.png}
\caption{Illustration (page 19)}
\end{figure}

egin{figure}[htbp]
\centering
\includegraphics[width=\linewidth]{figures/cordes_solistes__violon__p19_30.png}
\caption{Illustration (page 19)}
\end{figure}

egin{figure}[htbp]
\centering
\includegraphics[width=\linewidth]{figures/cordes_solistes__violon__p19_31.png}
\caption{Illustration (page 19)}
\end{figure}

egin{figure}[htbp]
\centering
\includegraphics[width=\linewidth]{figures/cordes_solistes__violon__p19_32.jpg}
\caption{Illustration (page 19)}
\end{figure}

egin{figure}[htbp]
\centering
\includegraphics[width=\linewidth]{figures/cordes_solistes__violon__p19_33.png}
\caption{Illustration (page 19)}
\end{figure}
F1 corrigé = 369  Hz (caisse). Note  v4: l’ancienne valeur 1  516 Hz = bridge hill. SourceF1 (Hz)F2 (Hz)F3 (Hz)F4 (Hz)VoyelleNAccor d Giesler220– 600———registre- dépend ant—\textasciitilde{} SOL2020369 ±14792 ±381 4632 128(variable )56\textasciitilde{} CORRECTION v3 (février 2026): F1=377 Hz (N=309 échantillons ordinario). L’ancienne valeur (1 516 Hz) correspondait au “bridge hill”, résonance du chevalet (F3 spectral), non au premier formant perceptuel. La valeur corrigée 369 Hz est cohérente avec Giesler (220–600 Hz, voyelle /o/–/u/). F2 corrigé = 792 Hz (zone /o/ médium). Très faible variabilité (±14 Hz) confirmée sur 749 échantillons. L’alto soliste est l’instrument à cordes le plus stable spectralement en jeu ordinario. Rôle pivot harmonique entre les cordes graves (/u/) et le cluster de convergence (/o/). Violoncelle SourceF1 (Hz)F2 (Hz)F3 (Hz)F4 (Hz)VoyelleNAccor d Backus600– 900—————
\par\medskip

egin{figure}[htbp]
\centering
\includegraphics[width=\linewidth]{figures/cordes_solistes__alto__p20_01.jpg}
\caption{Illustration (page 20)}
\end{figure}

egin{figure}[htbp]
\centering
\includegraphics[width=\linewidth]{figures/cordes_solistes__alto__p20_02.jpg}
\caption{Illustration (page 20)}
\end{figure}

egin{figure}[htbp]
\centering
\includegraphics[width=\linewidth]{figures/cordes_solistes__alto__p20_03.png}
\caption{Illustration (page 20)}
\end{figure}

egin{figure}[htbp]
\centering
\includegraphics[width=\linewidth]{figures/cordes_solistes__alto__p20_04.jpg}
\caption{Illustration (page 20)}
\end{figure}

egin{figure}[htbp]
\centering
\includegraphics[width=\linewidth]{figures/cordes_solistes__alto__p20_05.jpg}
\caption{Illustration (page 20)}
\end{figure}

egin{figure}[htbp]
\centering
\includegraphics[width=\linewidth]{figures/cordes_solistes__alto__p20_06.jpg}
\caption{Illustration (page 20)}
\end{figure}

egin{figure}[htbp]
\centering
\includegraphics[width=\linewidth]{figures/cordes_solistes__alto__p20_07.jpg}
\caption{Illustration (page 20)}
\end{figure}

egin{figure}[htbp]
\centering
\includegraphics[width=\linewidth]{figures/cordes_solistes__alto__p20_08.png}
\caption{Illustration (page 20)}
\end{figure}

egin{figure}[htbp]
\centering
\includegraphics[width=\linewidth]{figures/cordes_solistes__alto__p20_09.jpg}
\caption{Illustration (page 20)}
\end{figure}

egin{figure}[htbp]
\centering
\includegraphics[width=\linewidth]{figures/cordes_solistes__alto__p20_10.jpg}
\caption{Illustration (page 20)}
\end{figure}

egin{figure}[htbp]
\centering
\includegraphics[width=\linewidth]{figures/cordes_solistes__alto__p20_11.png}
\caption{Illustration (page 20)}
\end{figure}

egin{figure}[htbp]
\centering
\includegraphics[width=\linewidth]{figures/cordes_solistes__alto__p20_12.png}
\caption{Illustration (page 20)}
\end{figure}

egin{figure}[htbp]
\centering
\includegraphics[width=\linewidth]{figures/cordes_solistes__alto__p20_13.png}
\caption{Illustration (page 20)}
\end{figure}

egin{figure}[htbp]
\centering
\includegraphics[width=\linewidth]{figures/cordes_solistes__alto__p20_14.png}
\caption{Illustration (page 20)}
\end{figure}

egin{figure}[htbp]
\centering
\includegraphics[width=\linewidth]{figures/cordes_solistes__alto__p20_15.png}
\caption{Illustration (page 20)}
\end{figure}

egin{figure}[htbp]
\centering
\includegraphics[width=\linewidth]{figures/cordes_solistes__alto__p20_16.png}
\caption{Illustration (page 20)}
\end{figure}

egin{figure}[htbp]
\centering
\includegraphics[width=\linewidth]{figures/cordes_solistes__alto__p20_17.png}
\caption{Illustration (page 20)}
\end{figure}

egin{figure}[htbp]
\centering
\includegraphics[width=\linewidth]{figures/cordes_solistes__alto__p20_18.png}
\caption{Illustration (page 20)}
\end{figure}

egin{figure}[htbp]
\centering
\includegraphics[width=\linewidth]{figures/cordes_solistes__alto__p20_19.png}
\caption{Illustration (page 20)}
\end{figure}

egin{figure}[htbp]
\centering
\includegraphics[width=\linewidth]{figures/cordes_solistes__alto__p20_20.jpg}
\caption{Illustration (page 20)}
\end{figure}

egin{figure}[htbp]
\centering
\includegraphics[width=\linewidth]{figures/cordes_solistes__alto__p20_21.jpg}
\caption{Illustration (page 20)}
\end{figure}

egin{figure}[htbp]
\centering
\includegraphics[width=\linewidth]{figures/cordes_solistes__alto__p20_22.jpg}
\caption{Illustration (page 20)}
\end{figure}

egin{figure}[htbp]
\centering
\includegraphics[width=\linewidth]{figures/cordes_solistes__alto__p20_23.jpg}
\caption{Illustration (page 20)}
\end{figure}

egin{figure}[htbp]
\centering
\includegraphics[width=\linewidth]{figures/cordes_solistes__alto__p20_24.jpg}
\caption{Illustration (page 20)}
\end{figure}

egin{figure}[htbp]
\centering
\includegraphics[width=\linewidth]{figures/cordes_solistes__alto__p20_25.jpg}
\caption{Illustration (page 20)}
\end{figure}

egin{figure}[htbp]
\centering
\includegraphics[width=\linewidth]{figures/cordes_solistes__alto__p20_26.jpg}
\caption{Illustration (page 20)}
\end{figure}

egin{figure}[htbp]
\centering
\includegraphics[width=\linewidth]{figures/cordes_solistes__alto__p20_27.png}
\caption{Illustration (page 20)}
\end{figure}

egin{figure}[htbp]
\centering
\includegraphics[width=\linewidth]{figures/cordes_solistes__alto__p20_28.png}
\caption{Illustration (page 20)}
\end{figure}

egin{figure}[htbp]
\centering
\includegraphics[width=\linewidth]{figures/cordes_solistes__alto__p20_29.png}
\caption{Illustration (page 20)}
\end{figure}

egin{figure}[htbp]
\centering
\includegraphics[width=\linewidth]{figures/cordes_solistes__alto__p20_30.png}
\caption{Illustration (page 20)}
\end{figure}

egin{figure}[htbp]
\centering
\includegraphics[width=\linewidth]{figures/cordes_solistes__alto__p20_31.png}
\caption{Illustration (page 20)}
\end{figure}

egin{figure}[htbp]
\centering
\includegraphics[width=\linewidth]{figures/cordes_solistes__alto__p20_32.jpg}
\caption{Illustration (page 20)}
\end{figure}

egin{figure}[htbp]
\centering
\includegraphics[width=\linewidth]{figures/cordes_solistes__alto__p20_33.png}
\caption{Illustration (page 20)}
\end{figure}
SourceF1 (Hz)F2 (Hz)F3 (Hz)F4 (Hz)VoyelleNAccor d Giesler300– 500———o— Meyer450———o— SOL2020499 ±6391 577 ±10322 3223 127o/a197✓✓ F1=499 Hz dans le cluster de convergence. Fusion légendaire avec basson (Δ=3 Hz). Giesler 300–500, Backus 450, Meyer 450: excellent accord SOL. Meyer: 'timbres frères' Violoncelle-Basson. [Note v4 — Fp ≠ F1 spectral]: Le SOL2020 v3 (N=291 ordinario) révèle un F1 spectral à 205 Hz (/u/ grave) — résonance de la table d'harmonie du violoncelle. Le Fp=499 Hz du document correspond au F2 spectral (506 Hz, Δ=7 Hz), second mode mais nettement dominant dans le rayonnement perceptuel. La fusion légendaire violoncelle-basson s'explique par la quasi-identité de leurs F2 spectraux (violoncelle 506 Hz, basson \textasciitilde{}884 Hz en F2… et Fp basson = 502 Hz = F1 spectral). Autrement dit: le Fp du violoncelle converge avec le Fp (=F1) du basson — deux instruments qui atteignent le même formant par des mécanismes acoustiques différents. Contrebasse SourceF1 (Hz)F2 (Hz)F3 (Hz)F4 (Hz)VoyelleNAccor d Giesler70–250400 (Neben)——u/o— SOL2020200 ±169593 ±3121 0511 315u (oo)330✓✓ Zone u (oo), fondation orchestrale. Giesler: "Hauptformant 70–250 Hz (u), Nebenformant 400 Hz (o)". Convergence tuba (Δ=49 Hz).
\par\medskip

egin{figure}[htbp]
\centering
\includegraphics[width=\linewidth]{figures/cordes_solistes__basson__p21_01.jpg}
\caption{Illustration (page 21)}
\end{figure}

egin{figure}[htbp]
\centering
\includegraphics[width=\linewidth]{figures/cordes_solistes__basson__p21_02.jpg}
\caption{Illustration (page 21)}
\end{figure}

egin{figure}[htbp]
\centering
\includegraphics[width=\linewidth]{figures/cordes_solistes__basson__p21_03.png}
\caption{Illustration (page 21)}
\end{figure}

egin{figure}[htbp]
\centering
\includegraphics[width=\linewidth]{figures/cordes_solistes__basson__p21_04.jpg}
\caption{Illustration (page 21)}
\end{figure}

egin{figure}[htbp]
\centering
\includegraphics[width=\linewidth]{figures/cordes_solistes__basson__p21_05.jpg}
\caption{Illustration (page 21)}
\end{figure}

egin{figure}[htbp]
\centering
\includegraphics[width=\linewidth]{figures/cordes_solistes__basson__p21_06.jpg}
\caption{Illustration (page 21)}
\end{figure}

egin{figure}[htbp]
\centering
\includegraphics[width=\linewidth]{figures/cordes_solistes__basson__p21_07.jpg}
\caption{Illustration (page 21)}
\end{figure}

egin{figure}[htbp]
\centering
\includegraphics[width=\linewidth]{figures/cordes_solistes__basson__p21_08.png}
\caption{Illustration (page 21)}
\end{figure}

egin{figure}[htbp]
\centering
\includegraphics[width=\linewidth]{figures/cordes_solistes__basson__p21_09.jpg}
\caption{Illustration (page 21)}
\end{figure}

egin{figure}[htbp]
\centering
\includegraphics[width=\linewidth]{figures/cordes_solistes__basson__p21_10.jpg}
\caption{Illustration (page 21)}
\end{figure}

egin{figure}[htbp]
\centering
\includegraphics[width=\linewidth]{figures/cordes_solistes__basson__p21_11.png}
\caption{Illustration (page 21)}
\end{figure}

egin{figure}[htbp]
\centering
\includegraphics[width=\linewidth]{figures/cordes_solistes__basson__p21_12.png}
\caption{Illustration (page 21)}
\end{figure}

egin{figure}[htbp]
\centering
\includegraphics[width=\linewidth]{figures/cordes_solistes__basson__p21_13.png}
\caption{Illustration (page 21)}
\end{figure}

egin{figure}[htbp]
\centering
\includegraphics[width=\linewidth]{figures/cordes_solistes__basson__p21_14.png}
\caption{Illustration (page 21)}
\end{figure}

egin{figure}[htbp]
\centering
\includegraphics[width=\linewidth]{figures/cordes_solistes__basson__p21_15.png}
\caption{Illustration (page 21)}
\end{figure}

egin{figure}[htbp]
\centering
\includegraphics[width=\linewidth]{figures/cordes_solistes__basson__p21_16.png}
\caption{Illustration (page 21)}
\end{figure}

egin{figure}[htbp]
\centering
\includegraphics[width=\linewidth]{figures/cordes_solistes__basson__p21_17.png}
\caption{Illustration (page 21)}
\end{figure}

egin{figure}[htbp]
\centering
\includegraphics[width=\linewidth]{figures/cordes_solistes__basson__p21_18.png}
\caption{Illustration (page 21)}
\end{figure}

egin{figure}[htbp]
\centering
\includegraphics[width=\linewidth]{figures/cordes_solistes__basson__p21_19.png}
\caption{Illustration (page 21)}
\end{figure}

egin{figure}[htbp]
\centering
\includegraphics[width=\linewidth]{figures/cordes_solistes__basson__p21_20.jpg}
\caption{Illustration (page 21)}
\end{figure}

egin{figure}[htbp]
\centering
\includegraphics[width=\linewidth]{figures/cordes_solistes__basson__p21_21.jpg}
\caption{Illustration (page 21)}
\end{figure}

egin{figure}[htbp]
\centering
\includegraphics[width=\linewidth]{figures/cordes_solistes__basson__p21_22.jpg}
\caption{Illustration (page 21)}
\end{figure}

egin{figure}[htbp]
\centering
\includegraphics[width=\linewidth]{figures/cordes_solistes__basson__p21_23.jpg}
\caption{Illustration (page 21)}
\end{figure}

egin{figure}[htbp]
\centering
\includegraphics[width=\linewidth]{figures/cordes_solistes__basson__p21_24.jpg}
\caption{Illustration (page 21)}
\end{figure}

egin{figure}[htbp]
\centering
\includegraphics[width=\linewidth]{figures/cordes_solistes__basson__p21_25.jpg}
\caption{Illustration (page 21)}
\end{figure}

egin{figure}[htbp]
\centering
\includegraphics[width=\linewidth]{figures/cordes_solistes__basson__p21_26.jpg}
\caption{Illustration (page 21)}
\end{figure}

egin{figure}[htbp]
\centering
\includegraphics[width=\linewidth]{figures/cordes_solistes__basson__p21_27.png}
\caption{Illustration (page 21)}
\end{figure}

egin{figure}[htbp]
\centering
\includegraphics[width=\linewidth]{figures/cordes_solistes__basson__p21_28.png}
\caption{Illustration (page 21)}
\end{figure}

egin{figure}[htbp]
\centering
\includegraphics[width=\linewidth]{figures/cordes_solistes__basson__p21_29.png}
\caption{Illustration (page 21)}
\end{figure}

egin{figure}[htbp]
\centering
\includegraphics[width=\linewidth]{figures/cordes_solistes__basson__p21_30.png}
\caption{Illustration (page 21)}
\end{figure}

egin{figure}[htbp]
\centering
\includegraphics[width=\linewidth]{figures/cordes_solistes__basson__p21_31.png}
\caption{Illustration (page 21)}
\end{figure}

egin{figure}[htbp]
\centering
\includegraphics[width=\linewidth]{figures/cordes_solistes__basson__p21_32.jpg}
\caption{Illustration (page 21)}
\end{figure}

egin{figure}[htbp]
\centering
\includegraphics[width=\linewidth]{figures/cordes_solistes__basson__p21_33.png}
\caption{Illustration (page 21)}
\end{figure}

\section{Les Cordes d'Ensemble}
VI. Les Cordes d'Ensemble — NOUVEAU Ces données (Yan\_Adds) représentent le son d'une section orchestrale complète. Les différences avec les instruments solistes sont majeures et ont des conséquences directes sur l'orchestration. Comparaison soliste vs. ensemble InstrumentSoliste F1Ensemble F1Δ Observation Violon1 0341 556–834Formant plus bas en section (dispersion réduite) Alto3691 190+821Idem, timbre plus sombre en section Violoncelle 499587+88Légèrement plus haut en ensemble Contrebasse 200350+150Nettement plus haut en ensemble Les cordes aiguës (violon, alto) ont un F1 plus HAUT en section qu'en soliste (violon: 1 034→1 556 Hz; alto corrigé: 369→1 190 Hz), de même que les cordes graves (violoncelle, contrebasse). L'effet de section 'étire' les formants vers le médium-aigu — v3: l'observation précédente était basée sur l'ancien F1 alto erroné (1 516 Hz). Ensemble de Violons SourceF1 (Hz)F2 (Hz)F3 (Hz)F4 (Hz)VoyelleNAccor d Yan\_Adds1 556 ±9333 073 ±11464 504 ±12306 257 ±1361e (eh)166— F1=1 556 Hz, zone 'e (eh)' de clarté. Note v3: l'alto soliste corrigé (F1=369 Hz, /o/) ne converge plus dans cette zone; la proximité précédente (1 516 Hz) correspondait au bridge hill, non au F1 perceptuel. Variation: p=1 335, mf=1 517, ff=1 821 Hz. Ensemble de Violons — sourdine SourceF1 (Hz)F2 (Hz)F3 (Hz)F4 (Hz)VoyelleNAccor d Yan\_Adds777 ±2511 997 ±10003 9255 827a/å42— Effet sourdine: F1 passe de 1 556 à 777 Hz (divisé par 2). Déplacement de 'e (eh)' vers 'å/a' — assombrissement radical du timbre.
\par\medskip

egin{figure}[htbp]
\centering
\includegraphics[width=\linewidth]{figures/cordes_ensemble__violon__p22_01.jpg}
\caption{Illustration (page 22)}
\end{figure}

egin{figure}[htbp]
\centering
\includegraphics[width=\linewidth]{figures/cordes_ensemble__violon__p22_02.jpg}
\caption{Illustration (page 22)}
\end{figure}

egin{figure}[htbp]
\centering
\includegraphics[width=\linewidth]{figures/cordes_ensemble__violon__p22_03.png}
\caption{Illustration (page 22)}
\end{figure}

egin{figure}[htbp]
\centering
\includegraphics[width=\linewidth]{figures/cordes_ensemble__violon__p22_04.jpg}
\caption{Illustration (page 22)}
\end{figure}

egin{figure}[htbp]
\centering
\includegraphics[width=\linewidth]{figures/cordes_ensemble__violon__p22_05.jpg}
\caption{Illustration (page 22)}
\end{figure}

egin{figure}[htbp]
\centering
\includegraphics[width=\linewidth]{figures/cordes_ensemble__violon__p22_06.jpg}
\caption{Illustration (page 22)}
\end{figure}

egin{figure}[htbp]
\centering
\includegraphics[width=\linewidth]{figures/cordes_ensemble__violon__p22_07.jpg}
\caption{Illustration (page 22)}
\end{figure}

egin{figure}[htbp]
\centering
\includegraphics[width=\linewidth]{figures/cordes_ensemble__violon__p22_08.png}
\caption{Illustration (page 22)}
\end{figure}

egin{figure}[htbp]
\centering
\includegraphics[width=\linewidth]{figures/cordes_ensemble__violon__p22_09.jpg}
\caption{Illustration (page 22)}
\end{figure}

egin{figure}[htbp]
\centering
\includegraphics[width=\linewidth]{figures/cordes_ensemble__violon__p22_10.jpg}
\caption{Illustration (page 22)}
\end{figure}

egin{figure}[htbp]
\centering
\includegraphics[width=\linewidth]{figures/cordes_ensemble__violon__p22_11.png}
\caption{Illustration (page 22)}
\end{figure}

egin{figure}[htbp]
\centering
\includegraphics[width=\linewidth]{figures/cordes_ensemble__violon__p22_12.png}
\caption{Illustration (page 22)}
\end{figure}

egin{figure}[htbp]
\centering
\includegraphics[width=\linewidth]{figures/cordes_ensemble__violon__p22_13.png}
\caption{Illustration (page 22)}
\end{figure}

egin{figure}[htbp]
\centering
\includegraphics[width=\linewidth]{figures/cordes_ensemble__violon__p22_14.png}
\caption{Illustration (page 22)}
\end{figure}

egin{figure}[htbp]
\centering
\includegraphics[width=\linewidth]{figures/cordes_ensemble__violon__p22_15.png}
\caption{Illustration (page 22)}
\end{figure}

egin{figure}[htbp]
\centering
\includegraphics[width=\linewidth]{figures/cordes_ensemble__violon__p22_16.png}
\caption{Illustration (page 22)}
\end{figure}

egin{figure}[htbp]
\centering
\includegraphics[width=\linewidth]{figures/cordes_ensemble__violon__p22_17.png}
\caption{Illustration (page 22)}
\end{figure}

egin{figure}[htbp]
\centering
\includegraphics[width=\linewidth]{figures/cordes_ensemble__violon__p22_18.png}
\caption{Illustration (page 22)}
\end{figure}

egin{figure}[htbp]
\centering
\includegraphics[width=\linewidth]{figures/cordes_ensemble__violon__p22_19.png}
\caption{Illustration (page 22)}
\end{figure}

egin{figure}[htbp]
\centering
\includegraphics[width=\linewidth]{figures/cordes_ensemble__violon__p22_20.jpg}
\caption{Illustration (page 22)}
\end{figure}

egin{figure}[htbp]
\centering
\includegraphics[width=\linewidth]{figures/cordes_ensemble__violon__p22_21.jpg}
\caption{Illustration (page 22)}
\end{figure}

egin{figure}[htbp]
\centering
\includegraphics[width=\linewidth]{figures/cordes_ensemble__violon__p22_22.jpg}
\caption{Illustration (page 22)}
\end{figure}

egin{figure}[htbp]
\centering
\includegraphics[width=\linewidth]{figures/cordes_ensemble__violon__p22_23.jpg}
\caption{Illustration (page 22)}
\end{figure}

egin{figure}[htbp]
\centering
\includegraphics[width=\linewidth]{figures/cordes_ensemble__violon__p22_24.jpg}
\caption{Illustration (page 22)}
\end{figure}

egin{figure}[htbp]
\centering
\includegraphics[width=\linewidth]{figures/cordes_ensemble__violon__p22_25.jpg}
\caption{Illustration (page 22)}
\end{figure}

egin{figure}[htbp]
\centering
\includegraphics[width=\linewidth]{figures/cordes_ensemble__violon__p22_26.jpg}
\caption{Illustration (page 22)}
\end{figure}

egin{figure}[htbp]
\centering
\includegraphics[width=\linewidth]{figures/cordes_ensemble__violon__p22_27.png}
\caption{Illustration (page 22)}
\end{figure}

egin{figure}[htbp]
\centering
\includegraphics[width=\linewidth]{figures/cordes_ensemble__violon__p22_28.png}
\caption{Illustration (page 22)}
\end{figure}

egin{figure}[htbp]
\centering
\includegraphics[width=\linewidth]{figures/cordes_ensemble__violon__p22_29.png}
\caption{Illustration (page 22)}
\end{figure}

egin{figure}[htbp]
\centering
\includegraphics[width=\linewidth]{figures/cordes_ensemble__violon__p22_30.png}
\caption{Illustration (page 22)}
\end{figure}

egin{figure}[htbp]
\centering
\includegraphics[width=\linewidth]{figures/cordes_ensemble__violon__p22_31.png}
\caption{Illustration (page 22)}
\end{figure}

egin{figure}[htbp]
\centering
\includegraphics[width=\linewidth]{figures/cordes_ensemble__violon__p22_32.jpg}
\caption{Illustration (page 22)}
\end{figure}

egin{figure}[htbp]
\centering
\includegraphics[width=\linewidth]{figures/cordes_ensemble__violon__p22_33.png}
\caption{Illustration (page 22)}
\end{figure}
Ensemble d'Altos SourceF1 (Hz)F2 (Hz)F3 (Hz)F4 (Hz)VoyelleNAccor d Yan\_Adds1 190 ±8872 550 ±12503 869 ±13405 021 ±1616a (ah)148— F1=1 190 Hz, frontière 'a/e'. Variation: pp=934, ff=1 371 Hz. Ensemble d'Altos — sourdine SourceF1 (Hz)F2 (Hz)F3 (Hz)F4 (Hz)VoyelleNAccor d Yan\_Adds683 ±3251 697 ±9543 7235 711å (aw)38— Sourdine: 1 190 → 683 Hz. Zone 'å (aw)', timbre très voilé. Ensemble de Violoncelles SourceF1 (Hz)F2 (Hz)F3 (Hz)F4 (Hz)VoyelleNAccor d Yan\_Adds587 ±5221 563 ±10652 717 ±14273 850 ±1679o/å147— F1=587 Hz, frontière 'o/å', légèrement plus haut que le soliste (499 Hz). Variation: p=464, mf=531, ff=743 Hz.
\par\medskip

egin{figure}[htbp]
\centering
\includegraphics[width=\linewidth]{figures/cordes_ensemble__alto__p23_01.jpg}
\caption{Illustration (page 23)}
\end{figure}

egin{figure}[htbp]
\centering
\includegraphics[width=\linewidth]{figures/cordes_ensemble__alto__p23_02.jpg}
\caption{Illustration (page 23)}
\end{figure}

egin{figure}[htbp]
\centering
\includegraphics[width=\linewidth]{figures/cordes_ensemble__alto__p23_03.png}
\caption{Illustration (page 23)}
\end{figure}

egin{figure}[htbp]
\centering
\includegraphics[width=\linewidth]{figures/cordes_ensemble__alto__p23_04.jpg}
\caption{Illustration (page 23)}
\end{figure}

egin{figure}[htbp]
\centering
\includegraphics[width=\linewidth]{figures/cordes_ensemble__alto__p23_05.jpg}
\caption{Illustration (page 23)}
\end{figure}

egin{figure}[htbp]
\centering
\includegraphics[width=\linewidth]{figures/cordes_ensemble__alto__p23_06.jpg}
\caption{Illustration (page 23)}
\end{figure}

egin{figure}[htbp]
\centering
\includegraphics[width=\linewidth]{figures/cordes_ensemble__alto__p23_07.jpg}
\caption{Illustration (page 23)}
\end{figure}

egin{figure}[htbp]
\centering
\includegraphics[width=\linewidth]{figures/cordes_ensemble__alto__p23_08.png}
\caption{Illustration (page 23)}
\end{figure}

egin{figure}[htbp]
\centering
\includegraphics[width=\linewidth]{figures/cordes_ensemble__alto__p23_09.jpg}
\caption{Illustration (page 23)}
\end{figure}

egin{figure}[htbp]
\centering
\includegraphics[width=\linewidth]{figures/cordes_ensemble__alto__p23_10.jpg}
\caption{Illustration (page 23)}
\end{figure}

egin{figure}[htbp]
\centering
\includegraphics[width=\linewidth]{figures/cordes_ensemble__alto__p23_11.png}
\caption{Illustration (page 23)}
\end{figure}

egin{figure}[htbp]
\centering
\includegraphics[width=\linewidth]{figures/cordes_ensemble__alto__p23_12.png}
\caption{Illustration (page 23)}
\end{figure}

egin{figure}[htbp]
\centering
\includegraphics[width=\linewidth]{figures/cordes_ensemble__alto__p23_13.png}
\caption{Illustration (page 23)}
\end{figure}

egin{figure}[htbp]
\centering
\includegraphics[width=\linewidth]{figures/cordes_ensemble__alto__p23_14.png}
\caption{Illustration (page 23)}
\end{figure}

egin{figure}[htbp]
\centering
\includegraphics[width=\linewidth]{figures/cordes_ensemble__alto__p23_15.png}
\caption{Illustration (page 23)}
\end{figure}

egin{figure}[htbp]
\centering
\includegraphics[width=\linewidth]{figures/cordes_ensemble__alto__p23_16.png}
\caption{Illustration (page 23)}
\end{figure}

egin{figure}[htbp]
\centering
\includegraphics[width=\linewidth]{figures/cordes_ensemble__alto__p23_17.png}
\caption{Illustration (page 23)}
\end{figure}

egin{figure}[htbp]
\centering
\includegraphics[width=\linewidth]{figures/cordes_ensemble__alto__p23_18.png}
\caption{Illustration (page 23)}
\end{figure}

egin{figure}[htbp]
\centering
\includegraphics[width=\linewidth]{figures/cordes_ensemble__alto__p23_19.png}
\caption{Illustration (page 23)}
\end{figure}

egin{figure}[htbp]
\centering
\includegraphics[width=\linewidth]{figures/cordes_ensemble__alto__p23_20.jpg}
\caption{Illustration (page 23)}
\end{figure}

egin{figure}[htbp]
\centering
\includegraphics[width=\linewidth]{figures/cordes_ensemble__alto__p23_21.jpg}
\caption{Illustration (page 23)}
\end{figure}

egin{figure}[htbp]
\centering
\includegraphics[width=\linewidth]{figures/cordes_ensemble__alto__p23_22.jpg}
\caption{Illustration (page 23)}
\end{figure}

egin{figure}[htbp]
\centering
\includegraphics[width=\linewidth]{figures/cordes_ensemble__alto__p23_23.jpg}
\caption{Illustration (page 23)}
\end{figure}

egin{figure}[htbp]
\centering
\includegraphics[width=\linewidth]{figures/cordes_ensemble__alto__p23_24.jpg}
\caption{Illustration (page 23)}
\end{figure}

egin{figure}[htbp]
\centering
\includegraphics[width=\linewidth]{figures/cordes_ensemble__alto__p23_25.jpg}
\caption{Illustration (page 23)}
\end{figure}

egin{figure}[htbp]
\centering
\includegraphics[width=\linewidth]{figures/cordes_ensemble__alto__p23_26.jpg}
\caption{Illustration (page 23)}
\end{figure}

egin{figure}[htbp]
\centering
\includegraphics[width=\linewidth]{figures/cordes_ensemble__alto__p23_27.png}
\caption{Illustration (page 23)}
\end{figure}

egin{figure}[htbp]
\centering
\includegraphics[width=\linewidth]{figures/cordes_ensemble__alto__p23_28.png}
\caption{Illustration (page 23)}
\end{figure}

egin{figure}[htbp]
\centering
\includegraphics[width=\linewidth]{figures/cordes_ensemble__alto__p23_29.png}
\caption{Illustration (page 23)}
\end{figure}

egin{figure}[htbp]
\centering
\includegraphics[width=\linewidth]{figures/cordes_ensemble__alto__p23_30.png}
\caption{Illustration (page 23)}
\end{figure}

egin{figure}[htbp]
\centering
\includegraphics[width=\linewidth]{figures/cordes_ensemble__alto__p23_31.png}
\caption{Illustration (page 23)}
\end{figure}

egin{figure}[htbp]
\centering
\includegraphics[width=\linewidth]{figures/cordes_ensemble__alto__p23_32.jpg}
\caption{Illustration (page 23)}
\end{figure}

egin{figure}[htbp]
\centering
\includegraphics[width=\linewidth]{figures/cordes_ensemble__alto__p23_33.png}
\caption{Illustration (page 23)}
\end{figure}
Doublures •Basson: F1 ensemble 587 vs basson 502 (Δ=85) — bonne proximité. •Flûte basse: F1 587 vs 737 (Δ=150) — modérée. Ensemble de Violoncelles — sourdine SourceF1 (Hz)F2 (Hz)F3 (Hz)F4 (Hz)VoyelleNAccor d Yan\_Adds423 ±481 050 ±2522 5133 975o (oh)38— Sourdine: 587 → 423 Hz. Zone 'o (oh)' de plénitude, très faible variabilité (±48 Hz). Convergence avec cor (457 Hz, Δ=34). Ensemble de Contrebasses (non-vibrato) SourceF1 (Hz)F2 (Hz)F3 (Hz)F4 (Hz)VoyelleNAccor d Yan\_Adds350 ±30962 ±2611 830 ±3342 439 ±387u/o76— F1=350 Hz, zone 'u/o', stabilité extrême (±30 Hz). Plus haut que le soliste (200 Hz). L'ensemble ajoute de la présence tout en conservant la profondeur.
\par\medskip

egin{figure}[htbp]
\centering
\includegraphics[width=\linewidth]{figures/cordes_ensemble__basson__p24_01.jpg}
\caption{Illustration (page 24)}
\end{figure}

egin{figure}[htbp]
\centering
\includegraphics[width=\linewidth]{figures/cordes_ensemble__basson__p24_02.jpg}
\caption{Illustration (page 24)}
\end{figure}

egin{figure}[htbp]
\centering
\includegraphics[width=\linewidth]{figures/cordes_ensemble__basson__p24_03.png}
\caption{Illustration (page 24)}
\end{figure}

egin{figure}[htbp]
\centering
\includegraphics[width=\linewidth]{figures/cordes_ensemble__basson__p24_04.jpg}
\caption{Illustration (page 24)}
\end{figure}

egin{figure}[htbp]
\centering
\includegraphics[width=\linewidth]{figures/cordes_ensemble__basson__p24_05.jpg}
\caption{Illustration (page 24)}
\end{figure}

egin{figure}[htbp]
\centering
\includegraphics[width=\linewidth]{figures/cordes_ensemble__basson__p24_06.jpg}
\caption{Illustration (page 24)}
\end{figure}

egin{figure}[htbp]
\centering
\includegraphics[width=\linewidth]{figures/cordes_ensemble__basson__p24_07.jpg}
\caption{Illustration (page 24)}
\end{figure}

egin{figure}[htbp]
\centering
\includegraphics[width=\linewidth]{figures/cordes_ensemble__basson__p24_08.png}
\caption{Illustration (page 24)}
\end{figure}

egin{figure}[htbp]
\centering
\includegraphics[width=\linewidth]{figures/cordes_ensemble__basson__p24_09.jpg}
\caption{Illustration (page 24)}
\end{figure}

egin{figure}[htbp]
\centering
\includegraphics[width=\linewidth]{figures/cordes_ensemble__basson__p24_10.jpg}
\caption{Illustration (page 24)}
\end{figure}

egin{figure}[htbp]
\centering
\includegraphics[width=\linewidth]{figures/cordes_ensemble__basson__p24_11.png}
\caption{Illustration (page 24)}
\end{figure}

egin{figure}[htbp]
\centering
\includegraphics[width=\linewidth]{figures/cordes_ensemble__basson__p24_12.png}
\caption{Illustration (page 24)}
\end{figure}

egin{figure}[htbp]
\centering
\includegraphics[width=\linewidth]{figures/cordes_ensemble__basson__p24_13.png}
\caption{Illustration (page 24)}
\end{figure}

egin{figure}[htbp]
\centering
\includegraphics[width=\linewidth]{figures/cordes_ensemble__basson__p24_14.png}
\caption{Illustration (page 24)}
\end{figure}

egin{figure}[htbp]
\centering
\includegraphics[width=\linewidth]{figures/cordes_ensemble__basson__p24_15.png}
\caption{Illustration (page 24)}
\end{figure}

egin{figure}[htbp]
\centering
\includegraphics[width=\linewidth]{figures/cordes_ensemble__basson__p24_16.png}
\caption{Illustration (page 24)}
\end{figure}

egin{figure}[htbp]
\centering
\includegraphics[width=\linewidth]{figures/cordes_ensemble__basson__p24_17.png}
\caption{Illustration (page 24)}
\end{figure}

egin{figure}[htbp]
\centering
\includegraphics[width=\linewidth]{figures/cordes_ensemble__basson__p24_18.png}
\caption{Illustration (page 24)}
\end{figure}

egin{figure}[htbp]
\centering
\includegraphics[width=\linewidth]{figures/cordes_ensemble__basson__p24_19.png}
\caption{Illustration (page 24)}
\end{figure}

egin{figure}[htbp]
\centering
\includegraphics[width=\linewidth]{figures/cordes_ensemble__basson__p24_20.jpg}
\caption{Illustration (page 24)}
\end{figure}

egin{figure}[htbp]
\centering
\includegraphics[width=\linewidth]{figures/cordes_ensemble__basson__p24_21.jpg}
\caption{Illustration (page 24)}
\end{figure}

egin{figure}[htbp]
\centering
\includegraphics[width=\linewidth]{figures/cordes_ensemble__basson__p24_22.jpg}
\caption{Illustration (page 24)}
\end{figure}

egin{figure}[htbp]
\centering
\includegraphics[width=\linewidth]{figures/cordes_ensemble__basson__p24_23.jpg}
\caption{Illustration (page 24)}
\end{figure}

egin{figure}[htbp]
\centering
\includegraphics[width=\linewidth]{figures/cordes_ensemble__basson__p24_24.jpg}
\caption{Illustration (page 24)}
\end{figure}

egin{figure}[htbp]
\centering
\includegraphics[width=\linewidth]{figures/cordes_ensemble__basson__p24_25.jpg}
\caption{Illustration (page 24)}
\end{figure}

egin{figure}[htbp]
\centering
\includegraphics[width=\linewidth]{figures/cordes_ensemble__basson__p24_26.jpg}
\caption{Illustration (page 24)}
\end{figure}

egin{figure}[htbp]
\centering
\includegraphics[width=\linewidth]{figures/cordes_ensemble__basson__p24_27.png}
\caption{Illustration (page 24)}
\end{figure}

egin{figure}[htbp]
\centering
\includegraphics[width=\linewidth]{figures/cordes_ensemble__basson__p24_28.png}
\caption{Illustration (page 24)}
\end{figure}

egin{figure}[htbp]
\centering
\includegraphics[width=\linewidth]{figures/cordes_ensemble__basson__p24_29.png}
\caption{Illustration (page 24)}
\end{figure}

egin{figure}[htbp]
\centering
\includegraphics[width=\linewidth]{figures/cordes_ensemble__basson__p24_30.png}
\caption{Illustration (page 24)}
\end{figure}

egin{figure}[htbp]
\centering
\includegraphics[width=\linewidth]{figures/cordes_ensemble__basson__p24_31.png}
\caption{Illustration (page 24)}
\end{figure}

egin{figure}[htbp]
\centering
\includegraphics[width=\linewidth]{figures/cordes_ensemble__basson__p24_32.jpg}
\caption{Illustration (page 24)}
\end{figure}

egin{figure}[htbp]
\centering
\includegraphics[width=\linewidth]{figures/cordes_ensemble__basson__p24_33.png}
\caption{Illustration (page 24)}
\end{figure}
Synthèse: Effet de la sourdine sur les formants SectionOrdinario F1Sourdine F1RatioDéplacement vocalique Violons1 556 Hz777 Hz×0.50 e (eh) → å (aw) Altos1 190 Hz683 Hz×0.57 a (ah) → å (aw) Violoncelles 587 Hz423 Hz×0.72 o/å → o (oh) La sourdine abaisse systématiquement F1 de 30 à 50 \%, déplaçant le timbre d'une à deux catégories vocaliques vers le grave. L'effet est plus prononcé sur les aigus (violons ×0.50) que sur les graves (violoncelles ×0.72).
\par\medskip

egin{figure}[htbp]
\centering
\includegraphics[width=\linewidth]{figures/cordes_ensemble__violon__p25_01.jpg}
\caption{Illustration (page 25)}
\end{figure}

egin{figure}[htbp]
\centering
\includegraphics[width=\linewidth]{figures/cordes_ensemble__violon__p25_02.jpg}
\caption{Illustration (page 25)}
\end{figure}

egin{figure}[htbp]
\centering
\includegraphics[width=\linewidth]{figures/cordes_ensemble__violon__p25_03.png}
\caption{Illustration (page 25)}
\end{figure}

egin{figure}[htbp]
\centering
\includegraphics[width=\linewidth]{figures/cordes_ensemble__violon__p25_04.jpg}
\caption{Illustration (page 25)}
\end{figure}

egin{figure}[htbp]
\centering
\includegraphics[width=\linewidth]{figures/cordes_ensemble__violon__p25_05.jpg}
\caption{Illustration (page 25)}
\end{figure}

egin{figure}[htbp]
\centering
\includegraphics[width=\linewidth]{figures/cordes_ensemble__violon__p25_06.jpg}
\caption{Illustration (page 25)}
\end{figure}

egin{figure}[htbp]
\centering
\includegraphics[width=\linewidth]{figures/cordes_ensemble__violon__p25_07.jpg}
\caption{Illustration (page 25)}
\end{figure}

egin{figure}[htbp]
\centering
\includegraphics[width=\linewidth]{figures/cordes_ensemble__violon__p25_08.png}
\caption{Illustration (page 25)}
\end{figure}

egin{figure}[htbp]
\centering
\includegraphics[width=\linewidth]{figures/cordes_ensemble__violon__p25_09.jpg}
\caption{Illustration (page 25)}
\end{figure}

egin{figure}[htbp]
\centering
\includegraphics[width=\linewidth]{figures/cordes_ensemble__violon__p25_10.jpg}
\caption{Illustration (page 25)}
\end{figure}

egin{figure}[htbp]
\centering
\includegraphics[width=\linewidth]{figures/cordes_ensemble__violon__p25_11.png}
\caption{Illustration (page 25)}
\end{figure}

egin{figure}[htbp]
\centering
\includegraphics[width=\linewidth]{figures/cordes_ensemble__violon__p25_12.png}
\caption{Illustration (page 25)}
\end{figure}

egin{figure}[htbp]
\centering
\includegraphics[width=\linewidth]{figures/cordes_ensemble__violon__p25_13.png}
\caption{Illustration (page 25)}
\end{figure}

egin{figure}[htbp]
\centering
\includegraphics[width=\linewidth]{figures/cordes_ensemble__violon__p25_14.png}
\caption{Illustration (page 25)}
\end{figure}

egin{figure}[htbp]
\centering
\includegraphics[width=\linewidth]{figures/cordes_ensemble__violon__p25_15.png}
\caption{Illustration (page 25)}
\end{figure}

egin{figure}[htbp]
\centering
\includegraphics[width=\linewidth]{figures/cordes_ensemble__violon__p25_16.png}
\caption{Illustration (page 25)}
\end{figure}

egin{figure}[htbp]
\centering
\includegraphics[width=\linewidth]{figures/cordes_ensemble__violon__p25_17.png}
\caption{Illustration (page 25)}
\end{figure}

egin{figure}[htbp]
\centering
\includegraphics[width=\linewidth]{figures/cordes_ensemble__violon__p25_18.png}
\caption{Illustration (page 25)}
\end{figure}

egin{figure}[htbp]
\centering
\includegraphics[width=\linewidth]{figures/cordes_ensemble__violon__p25_19.png}
\caption{Illustration (page 25)}
\end{figure}

egin{figure}[htbp]
\centering
\includegraphics[width=\linewidth]{figures/cordes_ensemble__violon__p25_20.jpg}
\caption{Illustration (page 25)}
\end{figure}

egin{figure}[htbp]
\centering
\includegraphics[width=\linewidth]{figures/cordes_ensemble__violon__p25_21.jpg}
\caption{Illustration (page 25)}
\end{figure}

egin{figure}[htbp]
\centering
\includegraphics[width=\linewidth]{figures/cordes_ensemble__violon__p25_22.jpg}
\caption{Illustration (page 25)}
\end{figure}

egin{figure}[htbp]
\centering
\includegraphics[width=\linewidth]{figures/cordes_ensemble__violon__p25_23.jpg}
\caption{Illustration (page 25)}
\end{figure}

egin{figure}[htbp]
\centering
\includegraphics[width=\linewidth]{figures/cordes_ensemble__violon__p25_24.jpg}
\caption{Illustration (page 25)}
\end{figure}

egin{figure}[htbp]
\centering
\includegraphics[width=\linewidth]{figures/cordes_ensemble__violon__p25_25.jpg}
\caption{Illustration (page 25)}
\end{figure}

egin{figure}[htbp]
\centering
\includegraphics[width=\linewidth]{figures/cordes_ensemble__violon__p25_26.jpg}
\caption{Illustration (page 25)}
\end{figure}

egin{figure}[htbp]
\centering
\includegraphics[width=\linewidth]{figures/cordes_ensemble__violon__p25_27.png}
\caption{Illustration (page 25)}
\end{figure}

egin{figure}[htbp]
\centering
\includegraphics[width=\linewidth]{figures/cordes_ensemble__violon__p25_28.png}
\caption{Illustration (page 25)}
\end{figure}

egin{figure}[htbp]
\centering
\includegraphics[width=\linewidth]{figures/cordes_ensemble__violon__p25_29.png}
\caption{Illustration (page 25)}
\end{figure}

egin{figure}[htbp]
\centering
\includegraphics[width=\linewidth]{figures/cordes_ensemble__violon__p25_30.png}
\caption{Illustration (page 25)}
\end{figure}

egin{figure}[htbp]
\centering
\includegraphics[width=\linewidth]{figures/cordes_ensemble__violon__p25_31.png}
\caption{Illustration (page 25)}
\end{figure}

egin{figure}[htbp]
\centering
\includegraphics[width=\linewidth]{figures/cordes_ensemble__violon__p25_32.jpg}
\caption{Illustration (page 25)}
\end{figure}

egin{figure}[htbp]
\centering
\includegraphics[width=\linewidth]{figures/cordes_ensemble__violon__p25_33.png}
\caption{Illustration (page 25)}
\end{figure}

\section{Synthèse — Convergences formantiques}
\subsection{Tuba contrebasse F1 = 471 Hz}
VII. Synthèse — Convergences Formantiques Le Cluster 450–502 Hz Découverte majeure — 6 instruments convergent dans 45 Hz seulement Cor F1 = 457 Hz Tuba contrebasse F1 = 471 Hz Contrebasson F1 = 478 Hz Trombone F1 = 491 Hz Violoncelle F1 = 499 Hz Basson F1 = 502 Hz Moyenne: 483 Hz | Médiane: 485 Hz | Écart-type: 18 Hz | Voyelle commune: /o/ Ce cluster explique les doublures de Brahms, Bruckner, Wagner, Ravel, Mahler. Quand six instruments partagent la même résonance (Δ < 50 Hz), leurs timbres fusionnent quasi- totalement. L'énergie spectrale s'additionne constructivement dans cette zone, créant homogénéité timbrale maximale, puissance sans rugosité, et la sensation de 'plénitude' caractéristique de la voyelle /o/. Les Fondations \textasciitilde{}200–350 Hz Tuba (249) + Contrebasse solo (200) + Contrebasson sol. (281) + Contrebasse ensemble (350) = fondations orchestrales. Le Tuba crée un pont acoustique double: F1 (249) dans la zone /u/, F2 (494) ≈ Basson F1 (502) — Δ=9 Hz, pont vers le cluster. La Zone 'a (ah)': 890–1 050 Hz Nouvelle découverte: le trombone basse (894), la clarinette basse (909), la clarinette contrebasse (963) et le cor anglais (1 045) convergent dans la zone de puissance 'a (ah)'. Ceci explique l'efficacité des doublures graves dans les tutti de Mahler et Strauss. Figure 1 — Positions formantiques des 29 instruments
\par\medskip

egin{figure}[htbp]
\centering
\includegraphics[width=\linewidth]{figures/synthese__tuba_contrebasse__p26_01.jpg}
\caption{Illustration (page 26)}
\end{figure}

egin{figure}[htbp]
\centering
\includegraphics[width=\linewidth]{figures/synthese__tuba_contrebasse__p26_02.jpg}
\caption{Illustration (page 26)}
\end{figure}

egin{figure}[htbp]
\centering
\includegraphics[width=\linewidth]{figures/synthese__tuba_contrebasse__p26_03.png}
\caption{Illustration (page 26)}
\end{figure}

egin{figure}[htbp]
\centering
\includegraphics[width=\linewidth]{figures/synthese__tuba_contrebasse__p26_04.jpg}
\caption{Illustration (page 26)}
\end{figure}

egin{figure}[htbp]
\centering
\includegraphics[width=\linewidth]{figures/synthese__tuba_contrebasse__p26_05.jpg}
\caption{Illustration (page 26)}
\end{figure}

egin{figure}[htbp]
\centering
\includegraphics[width=\linewidth]{figures/synthese__tuba_contrebasse__p26_06.jpg}
\caption{Illustration (page 26)}
\end{figure}

egin{figure}[htbp]
\centering
\includegraphics[width=\linewidth]{figures/synthese__tuba_contrebasse__p26_07.jpg}
\caption{Illustration (page 26)}
\end{figure}

egin{figure}[htbp]
\centering
\includegraphics[width=\linewidth]{figures/synthese__tuba_contrebasse__p26_08.png}
\caption{Illustration (page 26)}
\end{figure}

egin{figure}[htbp]
\centering
\includegraphics[width=\linewidth]{figures/synthese__tuba_contrebasse__p26_09.jpg}
\caption{Illustration (page 26)}
\end{figure}

egin{figure}[htbp]
\centering
\includegraphics[width=\linewidth]{figures/synthese__tuba_contrebasse__p26_10.jpg}
\caption{Illustration (page 26)}
\end{figure}

egin{figure}[htbp]
\centering
\includegraphics[width=\linewidth]{figures/synthese__tuba_contrebasse__p26_11.png}
\caption{Illustration (page 26)}
\end{figure}

egin{figure}[htbp]
\centering
\includegraphics[width=\linewidth]{figures/synthese__tuba_contrebasse__p26_12.png}
\caption{Illustration (page 26)}
\end{figure}

egin{figure}[htbp]
\centering
\includegraphics[width=\linewidth]{figures/synthese__tuba_contrebasse__p26_13.png}
\caption{Illustration (page 26)}
\end{figure}

egin{figure}[htbp]
\centering
\includegraphics[width=\linewidth]{figures/synthese__tuba_contrebasse__p26_14.png}
\caption{Illustration (page 26)}
\end{figure}

egin{figure}[htbp]
\centering
\includegraphics[width=\linewidth]{figures/synthese__tuba_contrebasse__p26_15.png}
\caption{Illustration (page 26)}
\end{figure}

egin{figure}[htbp]
\centering
\includegraphics[width=\linewidth]{figures/synthese__tuba_contrebasse__p26_16.png}
\caption{Illustration (page 26)}
\end{figure}

egin{figure}[htbp]
\centering
\includegraphics[width=\linewidth]{figures/synthese__tuba_contrebasse__p26_17.png}
\caption{Illustration (page 26)}
\end{figure}

egin{figure}[htbp]
\centering
\includegraphics[width=\linewidth]{figures/synthese__tuba_contrebasse__p26_18.png}
\caption{Illustration (page 26)}
\end{figure}

egin{figure}[htbp]
\centering
\includegraphics[width=\linewidth]{figures/synthese__tuba_contrebasse__p26_19.png}
\caption{Illustration (page 26)}
\end{figure}

egin{figure}[htbp]
\centering
\includegraphics[width=\linewidth]{figures/synthese__tuba_contrebasse__p26_20.jpg}
\caption{Illustration (page 26)}
\end{figure}

egin{figure}[htbp]
\centering
\includegraphics[width=\linewidth]{figures/synthese__tuba_contrebasse__p26_21.jpg}
\caption{Illustration (page 26)}
\end{figure}

egin{figure}[htbp]
\centering
\includegraphics[width=\linewidth]{figures/synthese__tuba_contrebasse__p26_22.jpg}
\caption{Illustration (page 26)}
\end{figure}

egin{figure}[htbp]
\centering
\includegraphics[width=\linewidth]{figures/synthese__tuba_contrebasse__p26_23.jpg}
\caption{Illustration (page 26)}
\end{figure}

egin{figure}[htbp]
\centering
\includegraphics[width=\linewidth]{figures/synthese__tuba_contrebasse__p26_24.jpg}
\caption{Illustration (page 26)}
\end{figure}

egin{figure}[htbp]
\centering
\includegraphics[width=\linewidth]{figures/synthese__tuba_contrebasse__p26_25.jpg}
\caption{Illustration (page 26)}
\end{figure}

egin{figure}[htbp]
\centering
\includegraphics[width=\linewidth]{figures/synthese__tuba_contrebasse__p26_26.jpg}
\caption{Illustration (page 26)}
\end{figure}

egin{figure}[htbp]
\centering
\includegraphics[width=\linewidth]{figures/synthese__tuba_contrebasse__p26_27.png}
\caption{Illustration (page 26)}
\end{figure}

egin{figure}[htbp]
\centering
\includegraphics[width=\linewidth]{figures/synthese__tuba_contrebasse__p26_28.png}
\caption{Illustration (page 26)}
\end{figure}

egin{figure}[htbp]
\centering
\includegraphics[width=\linewidth]{figures/synthese__tuba_contrebasse__p26_29.png}
\caption{Illustration (page 26)}
\end{figure}

egin{figure}[htbp]
\centering
\includegraphics[width=\linewidth]{figures/synthese__tuba_contrebasse__p26_30.png}
\caption{Illustration (page 26)}
\end{figure}

egin{figure}[htbp]
\centering
\includegraphics[width=\linewidth]{figures/synthese__tuba_contrebasse__p26_31.png}
\caption{Illustration (page 26)}
\end{figure}

egin{figure}[htbp]
\centering
\includegraphics[width=\linewidth]{figures/synthese__tuba_contrebasse__p26_32.jpg}
\caption{Illustration (page 26)}
\end{figure}

egin{figure}[htbp]
\centering
\includegraphics[width=\linewidth]{figures/synthese__tuba_contrebasse__p26_33.png}
\caption{Illustration (page 26)}
\end{figure}
Figure 1 — Vue d'ensemble formantique, 29 instruments Ce diagramme représente les quatre premiers formants (F1 à F4) de chaque instrument sur un axe de fréquence logarithmique (150–8 000 Hz). Chaque instrument occupe une ligne horizontale, triés du plus grave (Contrebasse, F1=200 Hz) au plus aigu spectralement (Piccolo). Les cercles pleins proviennent de la base SOL2020 IRCAM (instruments solistes), les losanges des données Yan\_Adds (ensembles et instruments supplémentaires). La taille et l'opacité des marqueurs décroissent de F1 (le plus gros) à F4 (le plus petit), reflétant l'importance perceptive décroissante des formants successifs. Les bandes colorées en arrière-plan correspondent aux zones vocales de Meyer (2009): u (oo) en bleu pâle (200–400 Hz), o (oh) en vert pâle (400– 600 Hz), å (aw) en orange pâle (600–800 Hz), a (ah) en rose (800–1 250 Hz), e (eh) en violet pâle (1 250–2 600 Hz), i (ee) en jaune pâle (> 2 600 Hz). Deux zones de convergence sont surlignées: le cluster 450–502 Hz (rouge pâle) et la zone de fondation \textasciitilde{}200 Hz. Figure 2 — Espace vocalique F1/F2
\par\medskip

egin{figure}[htbp]
\centering
\includegraphics[width=\linewidth]{figures/synthese__contrebasse__p27_01.jpg}
\caption{Illustration (page 27)}
\end{figure}

egin{figure}[htbp]
\centering
\includegraphics[width=\linewidth]{figures/synthese__contrebasse__p27_02.jpg}
\caption{Illustration (page 27)}
\end{figure}

egin{figure}[htbp]
\centering
\includegraphics[width=\linewidth]{figures/synthese__contrebasse__p27_03.png}
\caption{Illustration (page 27)}
\end{figure}

egin{figure}[htbp]
\centering
\includegraphics[width=\linewidth]{figures/synthese__contrebasse__p27_04.jpg}
\caption{Illustration (page 27)}
\end{figure}

egin{figure}[htbp]
\centering
\includegraphics[width=\linewidth]{figures/synthese__contrebasse__p27_05.jpg}
\caption{Illustration (page 27)}
\end{figure}

egin{figure}[htbp]
\centering
\includegraphics[width=\linewidth]{figures/synthese__contrebasse__p27_06.jpg}
\caption{Illustration (page 27)}
\end{figure}

egin{figure}[htbp]
\centering
\includegraphics[width=\linewidth]{figures/synthese__contrebasse__p27_07.jpg}
\caption{Illustration (page 27)}
\end{figure}

egin{figure}[htbp]
\centering
\includegraphics[width=\linewidth]{figures/synthese__contrebasse__p27_08.png}
\caption{Illustration (page 27)}
\end{figure}

egin{figure}[htbp]
\centering
\includegraphics[width=\linewidth]{figures/synthese__contrebasse__p27_09.jpg}
\caption{Illustration (page 27)}
\end{figure}

egin{figure}[htbp]
\centering
\includegraphics[width=\linewidth]{figures/synthese__contrebasse__p27_10.jpg}
\caption{Illustration (page 27)}
\end{figure}

egin{figure}[htbp]
\centering
\includegraphics[width=\linewidth]{figures/synthese__contrebasse__p27_11.png}
\caption{Illustration (page 27)}
\end{figure}

egin{figure}[htbp]
\centering
\includegraphics[width=\linewidth]{figures/synthese__contrebasse__p27_12.png}
\caption{Illustration (page 27)}
\end{figure}

egin{figure}[htbp]
\centering
\includegraphics[width=\linewidth]{figures/synthese__contrebasse__p27_13.png}
\caption{Illustration (page 27)}
\end{figure}

egin{figure}[htbp]
\centering
\includegraphics[width=\linewidth]{figures/synthese__contrebasse__p27_14.png}
\caption{Illustration (page 27)}
\end{figure}

egin{figure}[htbp]
\centering
\includegraphics[width=\linewidth]{figures/synthese__contrebasse__p27_15.png}
\caption{Illustration (page 27)}
\end{figure}

egin{figure}[htbp]
\centering
\includegraphics[width=\linewidth]{figures/synthese__contrebasse__p27_16.png}
\caption{Illustration (page 27)}
\end{figure}

egin{figure}[htbp]
\centering
\includegraphics[width=\linewidth]{figures/synthese__contrebasse__p27_17.png}
\caption{Illustration (page 27)}
\end{figure}

egin{figure}[htbp]
\centering
\includegraphics[width=\linewidth]{figures/synthese__contrebasse__p27_18.png}
\caption{Illustration (page 27)}
\end{figure}

egin{figure}[htbp]
\centering
\includegraphics[width=\linewidth]{figures/synthese__contrebasse__p27_19.png}
\caption{Illustration (page 27)}
\end{figure}

egin{figure}[htbp]
\centering
\includegraphics[width=\linewidth]{figures/synthese__contrebasse__p27_20.jpg}
\caption{Illustration (page 27)}
\end{figure}

egin{figure}[htbp]
\centering
\includegraphics[width=\linewidth]{figures/synthese__contrebasse__p27_21.jpg}
\caption{Illustration (page 27)}
\end{figure}

egin{figure}[htbp]
\centering
\includegraphics[width=\linewidth]{figures/synthese__contrebasse__p27_22.jpg}
\caption{Illustration (page 27)}
\end{figure}

egin{figure}[htbp]
\centering
\includegraphics[width=\linewidth]{figures/synthese__contrebasse__p27_23.jpg}
\caption{Illustration (page 27)}
\end{figure}

egin{figure}[htbp]
\centering
\includegraphics[width=\linewidth]{figures/synthese__contrebasse__p27_24.jpg}
\caption{Illustration (page 27)}
\end{figure}

egin{figure}[htbp]
\centering
\includegraphics[width=\linewidth]{figures/synthese__contrebasse__p27_25.jpg}
\caption{Illustration (page 27)}
\end{figure}

egin{figure}[htbp]
\centering
\includegraphics[width=\linewidth]{figures/synthese__contrebasse__p27_26.jpg}
\caption{Illustration (page 27)}
\end{figure}

egin{figure}[htbp]
\centering
\includegraphics[width=\linewidth]{figures/synthese__contrebasse__p27_27.png}
\caption{Illustration (page 27)}
\end{figure}

egin{figure}[htbp]
\centering
\includegraphics[width=\linewidth]{figures/synthese__contrebasse__p27_28.png}
\caption{Illustration (page 27)}
\end{figure}

egin{figure}[htbp]
\centering
\includegraphics[width=\linewidth]{figures/synthese__contrebasse__p27_29.png}
\caption{Illustration (page 27)}
\end{figure}

egin{figure}[htbp]
\centering
\includegraphics[width=\linewidth]{figures/synthese__contrebasse__p27_30.png}
\caption{Illustration (page 27)}
\end{figure}

egin{figure}[htbp]
\centering
\includegraphics[width=\linewidth]{figures/synthese__contrebasse__p27_31.png}
\caption{Illustration (page 27)}
\end{figure}

egin{figure}[htbp]
\centering
\includegraphics[width=\linewidth]{figures/synthese__contrebasse__p27_32.jpg}
\caption{Illustration (page 27)}
\end{figure}

egin{figure}[htbp]
\centering
\includegraphics[width=\linewidth]{figures/synthese__contrebasse__p27_33.png}
\caption{Illustration (page 27)}
\end{figure}
Figure 2 — Espace vocalique F1/F2, 29 instruments Ce diagramme place chaque instrument dans un espace bidimensionnel défini par son premier formant (F1, axe horizontal) et son second formant (F2, axe vertical), selon la même convention que les diagrammes vocaliques en phonétique. Les ellipses en tirets gris représentent les zones des voyelles cardinales d'après Goldstein (Haskins Laboratories) et McCarthy. La forme des marqueurs distingue les familles: carrés pour les cuivres, losanges pour les bois, cercles pour les cordes solistes, triangles pour les cordes d'ensemble. La zone jaune en tirets signale le cluster de convergence 450–502 Hz, où six instruments se regroupent dans la zone vocalique 'o (oh)'. Ce diagramme confirme que les instruments qui fusionnent le mieux en orchestre partagent littéralement la même 'voyelle' acoustique. Figure 3 — Le Cluster de Convergence Figure 3 — Enveloppes spectrales des 6 instruments du cluster
\par\medskip

egin{figure}[htbp]
\centering
\includegraphics[width=\linewidth]{figures/synthese__contrebasse__p28_01.jpg}
\caption{Illustration (page 28)}
\end{figure}

egin{figure}[htbp]
\centering
\includegraphics[width=\linewidth]{figures/synthese__contrebasse__p28_02.jpg}
\caption{Illustration (page 28)}
\end{figure}

egin{figure}[htbp]
\centering
\includegraphics[width=\linewidth]{figures/synthese__contrebasse__p28_03.png}
\caption{Illustration (page 28)}
\end{figure}

egin{figure}[htbp]
\centering
\includegraphics[width=\linewidth]{figures/synthese__contrebasse__p28_04.jpg}
\caption{Illustration (page 28)}
\end{figure}

egin{figure}[htbp]
\centering
\includegraphics[width=\linewidth]{figures/synthese__contrebasse__p28_05.jpg}
\caption{Illustration (page 28)}
\end{figure}

egin{figure}[htbp]
\centering
\includegraphics[width=\linewidth]{figures/synthese__contrebasse__p28_06.jpg}
\caption{Illustration (page 28)}
\end{figure}

egin{figure}[htbp]
\centering
\includegraphics[width=\linewidth]{figures/synthese__contrebasse__p28_07.jpg}
\caption{Illustration (page 28)}
\end{figure}

egin{figure}[htbp]
\centering
\includegraphics[width=\linewidth]{figures/synthese__contrebasse__p28_08.png}
\caption{Illustration (page 28)}
\end{figure}

egin{figure}[htbp]
\centering
\includegraphics[width=\linewidth]{figures/synthese__contrebasse__p28_09.jpg}
\caption{Illustration (page 28)}
\end{figure}

egin{figure}[htbp]
\centering
\includegraphics[width=\linewidth]{figures/synthese__contrebasse__p28_10.jpg}
\caption{Illustration (page 28)}
\end{figure}

egin{figure}[htbp]
\centering
\includegraphics[width=\linewidth]{figures/synthese__contrebasse__p28_11.png}
\caption{Illustration (page 28)}
\end{figure}

egin{figure}[htbp]
\centering
\includegraphics[width=\linewidth]{figures/synthese__contrebasse__p28_12.png}
\caption{Illustration (page 28)}
\end{figure}

egin{figure}[htbp]
\centering
\includegraphics[width=\linewidth]{figures/synthese__contrebasse__p28_13.png}
\caption{Illustration (page 28)}
\end{figure}

egin{figure}[htbp]
\centering
\includegraphics[width=\linewidth]{figures/synthese__contrebasse__p28_14.png}
\caption{Illustration (page 28)}
\end{figure}

egin{figure}[htbp]
\centering
\includegraphics[width=\linewidth]{figures/synthese__contrebasse__p28_15.png}
\caption{Illustration (page 28)}
\end{figure}

egin{figure}[htbp]
\centering
\includegraphics[width=\linewidth]{figures/synthese__contrebasse__p28_16.png}
\caption{Illustration (page 28)}
\end{figure}

egin{figure}[htbp]
\centering
\includegraphics[width=\linewidth]{figures/synthese__contrebasse__p28_17.png}
\caption{Illustration (page 28)}
\end{figure}

egin{figure}[htbp]
\centering
\includegraphics[width=\linewidth]{figures/synthese__contrebasse__p28_18.png}
\caption{Illustration (page 28)}
\end{figure}

egin{figure}[htbp]
\centering
\includegraphics[width=\linewidth]{figures/synthese__contrebasse__p28_19.png}
\caption{Illustration (page 28)}
\end{figure}

egin{figure}[htbp]
\centering
\includegraphics[width=\linewidth]{figures/synthese__contrebasse__p28_20.jpg}
\caption{Illustration (page 28)}
\end{figure}

egin{figure}[htbp]
\centering
\includegraphics[width=\linewidth]{figures/synthese__contrebasse__p28_21.jpg}
\caption{Illustration (page 28)}
\end{figure}

egin{figure}[htbp]
\centering
\includegraphics[width=\linewidth]{figures/synthese__contrebasse__p28_22.jpg}
\caption{Illustration (page 28)}
\end{figure}

egin{figure}[htbp]
\centering
\includegraphics[width=\linewidth]{figures/synthese__contrebasse__p28_23.jpg}
\caption{Illustration (page 28)}
\end{figure}

egin{figure}[htbp]
\centering
\includegraphics[width=\linewidth]{figures/synthese__contrebasse__p28_24.jpg}
\caption{Illustration (page 28)}
\end{figure}

egin{figure}[htbp]
\centering
\includegraphics[width=\linewidth]{figures/synthese__contrebasse__p28_25.jpg}
\caption{Illustration (page 28)}
\end{figure}

egin{figure}[htbp]
\centering
\includegraphics[width=\linewidth]{figures/synthese__contrebasse__p28_26.jpg}
\caption{Illustration (page 28)}
\end{figure}

egin{figure}[htbp]
\centering
\includegraphics[width=\linewidth]{figures/synthese__contrebasse__p28_27.png}
\caption{Illustration (page 28)}
\end{figure}

egin{figure}[htbp]
\centering
\includegraphics[width=\linewidth]{figures/synthese__contrebasse__p28_28.png}
\caption{Illustration (page 28)}
\end{figure}

egin{figure}[htbp]
\centering
\includegraphics[width=\linewidth]{figures/synthese__contrebasse__p28_29.png}
\caption{Illustration (page 28)}
\end{figure}

egin{figure}[htbp]
\centering
\includegraphics[width=\linewidth]{figures/synthese__contrebasse__p28_30.png}
\caption{Illustration (page 28)}
\end{figure}

egin{figure}[htbp]
\centering
\includegraphics[width=\linewidth]{figures/synthese__contrebasse__p28_31.png}
\caption{Illustration (page 28)}
\end{figure}

egin{figure}[htbp]
\centering
\includegraphics[width=\linewidth]{figures/synthese__contrebasse__p28_32.jpg}
\caption{Illustration (page 28)}
\end{figure}

egin{figure}[htbp]
\centering
\includegraphics[width=\linewidth]{figures/synthese__contrebasse__p28_33.png}
\caption{Illustration (page 28)}
\end{figure}
\subsection{Contrebasson en vert foncé) et le Violoncelle en violet. La bande jaune verticale surligne la zone de}
Ce diagramme superpose les enveloppes spectrales schématiques de sept instruments clés, représentées sous forme de courbes gaussiennes centrées sur leurs formants mesurés. Les traits pleins épais représentent les trois instruments du 'triangle' original (Cor en bleu, Trombone en rouge, Basson en vert), tandis que les traits tiretés montrent les instruments qui élargissent ce triangle en un cluster de six (CB Tuba en orange, Contrebasson en vert foncé) et le Violoncelle en violet. La bande jaune verticale surligne la zone de convergence 450–502 Hz où les six pics F1 se superposent quasi-parfaitement. C'est le mécanisme acoustique de la fusion orchestrale: quand les résonances principales de plusieurs instruments coïncident dans une bande étroite, l'oreille perçoit un timbre unique et homogène plutôt que des sources séparées. Figure 4 — Matrice de convergence F1 (27 instruments) Figure 4 — Matrice de distance formantique (heatmap 27×27) Ce diagramme est une matrice symétrique de 27×27 instruments, où chaque cellule indique l'écart absolu en Hz entre les F1 de deux instruments. Code couleur: vert foncé (écart ≈ 0, convergence quasi-parfaite) au rouge foncé (écart > 2 000 Hz). Les instruments sont triés par F1 croissant, faisant apparaître un grand carré vert dans le coin supérieur gauche (instruments graves, F1 < 600 Hz). Ce diagramme met en évidence: (1) les 'blocs' d'instruments naturellement compatibles (le bloc vert du cluster, le bloc Tuba/Contrebasse); (2) les 'ponts' entre familles — Basson/Violoncelle (3 Hz), Clarinette basse/Trombone basse (15 Hz); (3) les combinaisons à éviter (rouge vif) — Piccolo + Contrebasse (2 136 Hz), Trompette + Tuba (2 352 Hz). C'est la 'carte de compatibilité timbrale' de l'orchestre complet.
\par\medskip

egin{figure}[htbp]
\centering
\includegraphics[width=\linewidth]{figures/synthese__trombone__p29_01.jpg}
\caption{Illustration (page 29)}
\end{figure}

egin{figure}[htbp]
\centering
\includegraphics[width=\linewidth]{figures/synthese__trombone__p29_02.jpg}
\caption{Illustration (page 29)}
\end{figure}

egin{figure}[htbp]
\centering
\includegraphics[width=\linewidth]{figures/synthese__trombone__p29_03.png}
\caption{Illustration (page 29)}
\end{figure}

egin{figure}[htbp]
\centering
\includegraphics[width=\linewidth]{figures/synthese__trombone__p29_04.jpg}
\caption{Illustration (page 29)}
\end{figure}

egin{figure}[htbp]
\centering
\includegraphics[width=\linewidth]{figures/synthese__trombone__p29_05.jpg}
\caption{Illustration (page 29)}
\end{figure}

egin{figure}[htbp]
\centering
\includegraphics[width=\linewidth]{figures/synthese__trombone__p29_06.jpg}
\caption{Illustration (page 29)}
\end{figure}

egin{figure}[htbp]
\centering
\includegraphics[width=\linewidth]{figures/synthese__trombone__p29_07.jpg}
\caption{Illustration (page 29)}
\end{figure}

egin{figure}[htbp]
\centering
\includegraphics[width=\linewidth]{figures/synthese__trombone__p29_08.png}
\caption{Illustration (page 29)}
\end{figure}

egin{figure}[htbp]
\centering
\includegraphics[width=\linewidth]{figures/synthese__trombone__p29_09.jpg}
\caption{Illustration (page 29)}
\end{figure}

egin{figure}[htbp]
\centering
\includegraphics[width=\linewidth]{figures/synthese__trombone__p29_10.jpg}
\caption{Illustration (page 29)}
\end{figure}

egin{figure}[htbp]
\centering
\includegraphics[width=\linewidth]{figures/synthese__trombone__p29_11.png}
\caption{Illustration (page 29)}
\end{figure}

egin{figure}[htbp]
\centering
\includegraphics[width=\linewidth]{figures/synthese__trombone__p29_12.png}
\caption{Illustration (page 29)}
\end{figure}

egin{figure}[htbp]
\centering
\includegraphics[width=\linewidth]{figures/synthese__trombone__p29_13.png}
\caption{Illustration (page 29)}
\end{figure}

egin{figure}[htbp]
\centering
\includegraphics[width=\linewidth]{figures/synthese__trombone__p29_14.png}
\caption{Illustration (page 29)}
\end{figure}

egin{figure}[htbp]
\centering
\includegraphics[width=\linewidth]{figures/synthese__trombone__p29_15.png}
\caption{Illustration (page 29)}
\end{figure}

egin{figure}[htbp]
\centering
\includegraphics[width=\linewidth]{figures/synthese__trombone__p29_16.png}
\caption{Illustration (page 29)}
\end{figure}

egin{figure}[htbp]
\centering
\includegraphics[width=\linewidth]{figures/synthese__trombone__p29_17.png}
\caption{Illustration (page 29)}
\end{figure}

egin{figure}[htbp]
\centering
\includegraphics[width=\linewidth]{figures/synthese__trombone__p29_18.png}
\caption{Illustration (page 29)}
\end{figure}

egin{figure}[htbp]
\centering
\includegraphics[width=\linewidth]{figures/synthese__trombone__p29_19.png}
\caption{Illustration (page 29)}
\end{figure}

egin{figure}[htbp]
\centering
\includegraphics[width=\linewidth]{figures/synthese__trombone__p29_20.jpg}
\caption{Illustration (page 29)}
\end{figure}

egin{figure}[htbp]
\centering
\includegraphics[width=\linewidth]{figures/synthese__trombone__p29_21.jpg}
\caption{Illustration (page 29)}
\end{figure}

egin{figure}[htbp]
\centering
\includegraphics[width=\linewidth]{figures/synthese__trombone__p29_22.jpg}
\caption{Illustration (page 29)}
\end{figure}

egin{figure}[htbp]
\centering
\includegraphics[width=\linewidth]{figures/synthese__trombone__p29_23.jpg}
\caption{Illustration (page 29)}
\end{figure}

egin{figure}[htbp]
\centering
\includegraphics[width=\linewidth]{figures/synthese__trombone__p29_24.jpg}
\caption{Illustration (page 29)}
\end{figure}

egin{figure}[htbp]
\centering
\includegraphics[width=\linewidth]{figures/synthese__trombone__p29_25.jpg}
\caption{Illustration (page 29)}
\end{figure}

egin{figure}[htbp]
\centering
\includegraphics[width=\linewidth]{figures/synthese__trombone__p29_26.jpg}
\caption{Illustration (page 29)}
\end{figure}

egin{figure}[htbp]
\centering
\includegraphics[width=\linewidth]{figures/synthese__trombone__p29_27.png}
\caption{Illustration (page 29)}
\end{figure}

egin{figure}[htbp]
\centering
\includegraphics[width=\linewidth]{figures/synthese__trombone__p29_28.png}
\caption{Illustration (page 29)}
\end{figure}

egin{figure}[htbp]
\centering
\includegraphics[width=\linewidth]{figures/synthese__trombone__p29_29.png}
\caption{Illustration (page 29)}
\end{figure}

egin{figure}[htbp]
\centering
\includegraphics[width=\linewidth]{figures/synthese__trombone__p29_30.png}
\caption{Illustration (page 29)}
\end{figure}

egin{figure}[htbp]
\centering
\includegraphics[width=\linewidth]{figures/synthese__trombone__p29_31.png}
\caption{Illustration (page 29)}
\end{figure}

egin{figure}[htbp]
\centering
\includegraphics[width=\linewidth]{figures/synthese__trombone__p29_32.jpg}
\caption{Illustration (page 29)}
\end{figure}

egin{figure}[htbp]
\centering
\includegraphics[width=\linewidth]{figures/synthese__trombone__p29_33.png}
\caption{Illustration (page 29)}
\end{figure}
Figure 5 — Carte Fp globale: F2 vs Fp (30 instruments) Comparaison F2 (peak-picking, gris) vs Fp (centroïde de bande, couleur) pour les 30 instruments audités. Les barres horizontales représentent ±1σ. Fp est systématiquement plus stable que F2, en particulier pour les instruments à large tessiture (cordes, clarinettes, trompette). Figure 6 — Enveloppes spectrales Yan\_Adds (14 instruments) Enveloppes spectrales moyennes des 14 instruments Yan\_Adds (technique ordinario/non-vibrato). Les lignes rouges indiquent F1-F4 (peak-picking v22), la ligne verte indique Fp (centroïde de bande).
\par\medskip

egin{figure}[htbp]
\centering
\includegraphics[width=\linewidth]{figures/synthese__trompette__p30_01.jpg}
\caption{Illustration (page 30)}
\end{figure}

egin{figure}[htbp]
\centering
\includegraphics[width=\linewidth]{figures/synthese__trompette__p30_02.jpg}
\caption{Illustration (page 30)}
\end{figure}

egin{figure}[htbp]
\centering
\includegraphics[width=\linewidth]{figures/synthese__trompette__p30_03.png}
\caption{Illustration (page 30)}
\end{figure}

egin{figure}[htbp]
\centering
\includegraphics[width=\linewidth]{figures/synthese__trompette__p30_04.jpg}
\caption{Illustration (page 30)}
\end{figure}

egin{figure}[htbp]
\centering
\includegraphics[width=\linewidth]{figures/synthese__trompette__p30_05.jpg}
\caption{Illustration (page 30)}
\end{figure}

egin{figure}[htbp]
\centering
\includegraphics[width=\linewidth]{figures/synthese__trompette__p30_06.jpg}
\caption{Illustration (page 30)}
\end{figure}

egin{figure}[htbp]
\centering
\includegraphics[width=\linewidth]{figures/synthese__trompette__p30_07.jpg}
\caption{Illustration (page 30)}
\end{figure}

egin{figure}[htbp]
\centering
\includegraphics[width=\linewidth]{figures/synthese__trompette__p30_08.png}
\caption{Illustration (page 30)}
\end{figure}

egin{figure}[htbp]
\centering
\includegraphics[width=\linewidth]{figures/synthese__trompette__p30_09.jpg}
\caption{Illustration (page 30)}
\end{figure}

egin{figure}[htbp]
\centering
\includegraphics[width=\linewidth]{figures/synthese__trompette__p30_10.jpg}
\caption{Illustration (page 30)}
\end{figure}

egin{figure}[htbp]
\centering
\includegraphics[width=\linewidth]{figures/synthese__trompette__p30_11.png}
\caption{Illustration (page 30)}
\end{figure}

egin{figure}[htbp]
\centering
\includegraphics[width=\linewidth]{figures/synthese__trompette__p30_12.png}
\caption{Illustration (page 30)}
\end{figure}

egin{figure}[htbp]
\centering
\includegraphics[width=\linewidth]{figures/synthese__trompette__p30_13.png}
\caption{Illustration (page 30)}
\end{figure}

egin{figure}[htbp]
\centering
\includegraphics[width=\linewidth]{figures/synthese__trompette__p30_14.png}
\caption{Illustration (page 30)}
\end{figure}

egin{figure}[htbp]
\centering
\includegraphics[width=\linewidth]{figures/synthese__trompette__p30_15.png}
\caption{Illustration (page 30)}
\end{figure}

egin{figure}[htbp]
\centering
\includegraphics[width=\linewidth]{figures/synthese__trompette__p30_16.png}
\caption{Illustration (page 30)}
\end{figure}

egin{figure}[htbp]
\centering
\includegraphics[width=\linewidth]{figures/synthese__trompette__p30_17.png}
\caption{Illustration (page 30)}
\end{figure}

egin{figure}[htbp]
\centering
\includegraphics[width=\linewidth]{figures/synthese__trompette__p30_18.png}
\caption{Illustration (page 30)}
\end{figure}

egin{figure}[htbp]
\centering
\includegraphics[width=\linewidth]{figures/synthese__trompette__p30_19.png}
\caption{Illustration (page 30)}
\end{figure}

egin{figure}[htbp]
\centering
\includegraphics[width=\linewidth]{figures/synthese__trompette__p30_20.jpg}
\caption{Illustration (page 30)}
\end{figure}

egin{figure}[htbp]
\centering
\includegraphics[width=\linewidth]{figures/synthese__trompette__p30_21.jpg}
\caption{Illustration (page 30)}
\end{figure}

egin{figure}[htbp]
\centering
\includegraphics[width=\linewidth]{figures/synthese__trompette__p30_22.jpg}
\caption{Illustration (page 30)}
\end{figure}

egin{figure}[htbp]
\centering
\includegraphics[width=\linewidth]{figures/synthese__trompette__p30_23.jpg}
\caption{Illustration (page 30)}
\end{figure}

egin{figure}[htbp]
\centering
\includegraphics[width=\linewidth]{figures/synthese__trompette__p30_24.jpg}
\caption{Illustration (page 30)}
\end{figure}

egin{figure}[htbp]
\centering
\includegraphics[width=\linewidth]{figures/synthese__trompette__p30_25.jpg}
\caption{Illustration (page 30)}
\end{figure}

egin{figure}[htbp]
\centering
\includegraphics[width=\linewidth]{figures/synthese__trompette__p30_26.jpg}
\caption{Illustration (page 30)}
\end{figure}

egin{figure}[htbp]
\centering
\includegraphics[width=\linewidth]{figures/synthese__trompette__p30_27.png}
\caption{Illustration (page 30)}
\end{figure}

egin{figure}[htbp]
\centering
\includegraphics[width=\linewidth]{figures/synthese__trompette__p30_28.png}
\caption{Illustration (page 30)}
\end{figure}

egin{figure}[htbp]
\centering
\includegraphics[width=\linewidth]{figures/synthese__trompette__p30_29.png}
\caption{Illustration (page 30)}
\end{figure}

egin{figure}[htbp]
\centering
\includegraphics[width=\linewidth]{figures/synthese__trompette__p30_30.png}
\caption{Illustration (page 30)}
\end{figure}

egin{figure}[htbp]
\centering
\includegraphics[width=\linewidth]{figures/synthese__trompette__p30_31.png}
\caption{Illustration (page 30)}
\end{figure}

egin{figure}[htbp]
\centering
\includegraphics[width=\linewidth]{figures/synthese__trompette__p30_32.jpg}
\caption{Illustration (page 30)}
\end{figure}

egin{figure}[htbp]
\centering
\includegraphics[width=\linewidth]{figures/synthese__trompette__p30_33.png}
\caption{Illustration (page 30)}
\end{figure}

\section{Espace vocalique complet}
VIII. Espace Vocalique Complet Les formants instrumentaux se positionnent dans l'espace des voyelles cardinales. Cette distribution n'est pas aléatoire: elle reflète les contraintes acoustiques fondamentales des résonateurs (tubes, colonnes d'air, corps vibrants) qui sont les mêmes pour la voix humaine et les instruments d'orchestre. VoyellePlage Hz Instruments (F1) /u/ (ou)200–400Contrebasse solo (200) · Tuba (249) · Contrebasson (281) · Contrebasse ensemble (350) /o/ (oh) ★ CLUSTER400–600Baryton Sax (374) · Ténor Sax (438) · Cor (457) · Tuba CB (471) · Cbn (478) · Trombone (491) · Vcl (499) · Basson (502) · Cor anglais (507) · Cl. basse (588) · Vcl- Ens (587) /å/ (aw)600–800 Vcl-Ens sord. (423) · Fl. basse (737) · Fl. CB (746) /a/ (ah)800–1 250Trb. basse (894) · Cl. basse (909) · Vln-Ens sord. (777) · Cl. Sib (1 016) · Violon (1 034) · Cor anglais (1 045) · Cl. Cb (963) · Alto-Ens (1 190) /e/ (eh)1 250–2 600Trompette (1 324) · Flûte (1 354) · Hautbois (1 460) · Vln-Ens (1 556) · Cl. Mib (1 747) · Piccolo (2 336) — [Alto déplacé: F1 corrigé = 369 Hz → zone /o/] /i/ (ee)> 2 600 Piccolo (2 658) L'orchestration peut être comprise comme une 'polyphonie de voyelles', où chaque instrument apporte une couleur vocalique spécifique. Les convergences formantiques correspondent aux 'accords de voyelles' — des combinaisons naturellement harmonieuses. Tous les 29 instruments se positionnent à l'intérieur de l'espace vocalique standard (F1: 200–3 000 Hz).
\par\medskip

egin{figure}[htbp]
\centering
\includegraphics[width=\linewidth]{figures/espace_vocalique__contrebasse__p31_01.jpg}
\caption{Illustration (page 31)}
\end{figure}

egin{figure}[htbp]
\centering
\includegraphics[width=\linewidth]{figures/espace_vocalique__contrebasse__p31_02.jpg}
\caption{Illustration (page 31)}
\end{figure}

egin{figure}[htbp]
\centering
\includegraphics[width=\linewidth]{figures/espace_vocalique__contrebasse__p31_03.png}
\caption{Illustration (page 31)}
\end{figure}

egin{figure}[htbp]
\centering
\includegraphics[width=\linewidth]{figures/espace_vocalique__contrebasse__p31_04.jpg}
\caption{Illustration (page 31)}
\end{figure}

egin{figure}[htbp]
\centering
\includegraphics[width=\linewidth]{figures/espace_vocalique__contrebasse__p31_05.jpg}
\caption{Illustration (page 31)}
\end{figure}

egin{figure}[htbp]
\centering
\includegraphics[width=\linewidth]{figures/espace_vocalique__contrebasse__p31_06.jpg}
\caption{Illustration (page 31)}
\end{figure}

egin{figure}[htbp]
\centering
\includegraphics[width=\linewidth]{figures/espace_vocalique__contrebasse__p31_07.jpg}
\caption{Illustration (page 31)}
\end{figure}

egin{figure}[htbp]
\centering
\includegraphics[width=\linewidth]{figures/espace_vocalique__contrebasse__p31_08.png}
\caption{Illustration (page 31)}
\end{figure}

egin{figure}[htbp]
\centering
\includegraphics[width=\linewidth]{figures/espace_vocalique__contrebasse__p31_09.jpg}
\caption{Illustration (page 31)}
\end{figure}

egin{figure}[htbp]
\centering
\includegraphics[width=\linewidth]{figures/espace_vocalique__contrebasse__p31_10.jpg}
\caption{Illustration (page 31)}
\end{figure}

egin{figure}[htbp]
\centering
\includegraphics[width=\linewidth]{figures/espace_vocalique__contrebasse__p31_11.png}
\caption{Illustration (page 31)}
\end{figure}

egin{figure}[htbp]
\centering
\includegraphics[width=\linewidth]{figures/espace_vocalique__contrebasse__p31_12.png}
\caption{Illustration (page 31)}
\end{figure}

egin{figure}[htbp]
\centering
\includegraphics[width=\linewidth]{figures/espace_vocalique__contrebasse__p31_13.png}
\caption{Illustration (page 31)}
\end{figure}

egin{figure}[htbp]
\centering
\includegraphics[width=\linewidth]{figures/espace_vocalique__contrebasse__p31_14.png}
\caption{Illustration (page 31)}
\end{figure}

egin{figure}[htbp]
\centering
\includegraphics[width=\linewidth]{figures/espace_vocalique__contrebasse__p31_15.png}
\caption{Illustration (page 31)}
\end{figure}

egin{figure}[htbp]
\centering
\includegraphics[width=\linewidth]{figures/espace_vocalique__contrebasse__p31_16.png}
\caption{Illustration (page 31)}
\end{figure}

egin{figure}[htbp]
\centering
\includegraphics[width=\linewidth]{figures/espace_vocalique__contrebasse__p31_17.png}
\caption{Illustration (page 31)}
\end{figure}

egin{figure}[htbp]
\centering
\includegraphics[width=\linewidth]{figures/espace_vocalique__contrebasse__p31_18.png}
\caption{Illustration (page 31)}
\end{figure}

egin{figure}[htbp]
\centering
\includegraphics[width=\linewidth]{figures/espace_vocalique__contrebasse__p31_19.png}
\caption{Illustration (page 31)}
\end{figure}

egin{figure}[htbp]
\centering
\includegraphics[width=\linewidth]{figures/espace_vocalique__contrebasse__p31_20.jpg}
\caption{Illustration (page 31)}
\end{figure}

egin{figure}[htbp]
\centering
\includegraphics[width=\linewidth]{figures/espace_vocalique__contrebasse__p31_21.jpg}
\caption{Illustration (page 31)}
\end{figure}

egin{figure}[htbp]
\centering
\includegraphics[width=\linewidth]{figures/espace_vocalique__contrebasse__p31_22.jpg}
\caption{Illustration (page 31)}
\end{figure}

egin{figure}[htbp]
\centering
\includegraphics[width=\linewidth]{figures/espace_vocalique__contrebasse__p31_23.jpg}
\caption{Illustration (page 31)}
\end{figure}

egin{figure}[htbp]
\centering
\includegraphics[width=\linewidth]{figures/espace_vocalique__contrebasse__p31_24.jpg}
\caption{Illustration (page 31)}
\end{figure}

egin{figure}[htbp]
\centering
\includegraphics[width=\linewidth]{figures/espace_vocalique__contrebasse__p31_25.jpg}
\caption{Illustration (page 31)}
\end{figure}

egin{figure}[htbp]
\centering
\includegraphics[width=\linewidth]{figures/espace_vocalique__contrebasse__p31_26.jpg}
\caption{Illustration (page 31)}
\end{figure}

egin{figure}[htbp]
\centering
\includegraphics[width=\linewidth]{figures/espace_vocalique__contrebasse__p31_27.png}
\caption{Illustration (page 31)}
\end{figure}

egin{figure}[htbp]
\centering
\includegraphics[width=\linewidth]{figures/espace_vocalique__contrebasse__p31_28.png}
\caption{Illustration (page 31)}
\end{figure}

egin{figure}[htbp]
\centering
\includegraphics[width=\linewidth]{figures/espace_vocalique__contrebasse__p31_29.png}
\caption{Illustration (page 31)}
\end{figure}

egin{figure}[htbp]
\centering
\includegraphics[width=\linewidth]{figures/espace_vocalique__contrebasse__p31_30.png}
\caption{Illustration (page 31)}
\end{figure}

egin{figure}[htbp]
\centering
\includegraphics[width=\linewidth]{figures/espace_vocalique__contrebasse__p31_31.png}
\caption{Illustration (page 31)}
\end{figure}

egin{figure}[htbp]
\centering
\includegraphics[width=\linewidth]{figures/espace_vocalique__contrebasse__p31_32.jpg}
\caption{Illustration (page 31)}
\end{figure}

egin{figure}[htbp]
\centering
\includegraphics[width=\linewidth]{figures/espace_vocalique__contrebasse__p31_33.png}
\caption{Illustration (page 31)}
\end{figure}

\section{Validation et comparaison avec la littérature}
IX. Validation et Comparaison avec la Littérature Concordance multi-sources — résultats globaux Niveau de concordanceNombre d'instrumentsPourcentage Critère ✓✓ Excellente 2379 \% < 80 Hz d'écart ✓ Bonne 414 \% 80–200 Hz d'écart \textasciitilde{} Modérée 27 \%Registre-dépendant TOTAL 2993 \%Taux de validation global Exemples de concordance parfaite InstrumentSOL2020 Littérature Écart BassonSOL 502 HzGiesler 500 · Backus 500 · Meyer 500 Δ = 0–2 Hz TromboneSOL 491 HzGiesler 520–600 · Backus 500 · Meyer 480– 600Δ = 9 Hz ViolonSOL 1 034 Hz Giesler 1 000–1 200 · Meyer \textasciitilde{}1 100 Δ = 66 Hz CorSOL 457 HzGiesler 250–500 · Backus 400–500 · Meyer 340Δ < 50 Hz Comparaison avec McCarty/CCRMA (2003) John McCarty (CCRMA, Stanford, 2003) a publié une analyse formantique de 11 instruments utilisant le toolkit COLEA (matlab). L'auteur lui-même note: "The accuracy of colea, or my ability to use it seems to be questionable" et "the accuracy of this chart is probably questionable". Critère McCarty (2003) SOL2020 (2026) Échantillons \textasciitilde{}11 fichiers uniques 2 298 (200–800 par instrument) Méthode COLEA (matlab)scipy.signal.find\_peaks + enveloppes spectrales 4096 FFT Validation Aucune source croisée Multi-sources (Giesler, Backus, Meyer) Reconnaissance Auteur reconnaît les limites Concordance 93 \% Date 2003 (technologie limitée) 2026 (méthodes modernes) Nos données SOL2020 sont EN ACCORD avec McCarty pour tous les instruments comparables, mais avec un échantillonnage 100–1 000× plus grand et une validation multi-sources indépendantes. Notre analyse confirme et étend les observations de McCarty avec une robustesse statistique significativement supérieure.
\par\medskip

egin{figure}[htbp]
\centering
\includegraphics[width=\linewidth]{figures/validation__basson__p32_01.jpg}
\caption{Illustration (page 32)}
\end{figure}

egin{figure}[htbp]
\centering
\includegraphics[width=\linewidth]{figures/validation__basson__p32_02.jpg}
\caption{Illustration (page 32)}
\end{figure}

egin{figure}[htbp]
\centering
\includegraphics[width=\linewidth]{figures/validation__basson__p32_03.png}
\caption{Illustration (page 32)}
\end{figure}

egin{figure}[htbp]
\centering
\includegraphics[width=\linewidth]{figures/validation__basson__p32_04.jpg}
\caption{Illustration (page 32)}
\end{figure}

egin{figure}[htbp]
\centering
\includegraphics[width=\linewidth]{figures/validation__basson__p32_05.jpg}
\caption{Illustration (page 32)}
\end{figure}

egin{figure}[htbp]
\centering
\includegraphics[width=\linewidth]{figures/validation__basson__p32_06.jpg}
\caption{Illustration (page 32)}
\end{figure}

egin{figure}[htbp]
\centering
\includegraphics[width=\linewidth]{figures/validation__basson__p32_07.jpg}
\caption{Illustration (page 32)}
\end{figure}

egin{figure}[htbp]
\centering
\includegraphics[width=\linewidth]{figures/validation__basson__p32_08.png}
\caption{Illustration (page 32)}
\end{figure}

egin{figure}[htbp]
\centering
\includegraphics[width=\linewidth]{figures/validation__basson__p32_09.jpg}
\caption{Illustration (page 32)}
\end{figure}

egin{figure}[htbp]
\centering
\includegraphics[width=\linewidth]{figures/validation__basson__p32_10.jpg}
\caption{Illustration (page 32)}
\end{figure}

egin{figure}[htbp]
\centering
\includegraphics[width=\linewidth]{figures/validation__basson__p32_11.png}
\caption{Illustration (page 32)}
\end{figure}

egin{figure}[htbp]
\centering
\includegraphics[width=\linewidth]{figures/validation__basson__p32_12.png}
\caption{Illustration (page 32)}
\end{figure}

egin{figure}[htbp]
\centering
\includegraphics[width=\linewidth]{figures/validation__basson__p32_13.png}
\caption{Illustration (page 32)}
\end{figure}

egin{figure}[htbp]
\centering
\includegraphics[width=\linewidth]{figures/validation__basson__p32_14.png}
\caption{Illustration (page 32)}
\end{figure}

egin{figure}[htbp]
\centering
\includegraphics[width=\linewidth]{figures/validation__basson__p32_15.png}
\caption{Illustration (page 32)}
\end{figure}

egin{figure}[htbp]
\centering
\includegraphics[width=\linewidth]{figures/validation__basson__p32_16.png}
\caption{Illustration (page 32)}
\end{figure}

egin{figure}[htbp]
\centering
\includegraphics[width=\linewidth]{figures/validation__basson__p32_17.png}
\caption{Illustration (page 32)}
\end{figure}

egin{figure}[htbp]
\centering
\includegraphics[width=\linewidth]{figures/validation__basson__p32_18.png}
\caption{Illustration (page 32)}
\end{figure}

egin{figure}[htbp]
\centering
\includegraphics[width=\linewidth]{figures/validation__basson__p32_19.png}
\caption{Illustration (page 32)}
\end{figure}

egin{figure}[htbp]
\centering
\includegraphics[width=\linewidth]{figures/validation__basson__p32_20.jpg}
\caption{Illustration (page 32)}
\end{figure}

egin{figure}[htbp]
\centering
\includegraphics[width=\linewidth]{figures/validation__basson__p32_21.jpg}
\caption{Illustration (page 32)}
\end{figure}

egin{figure}[htbp]
\centering
\includegraphics[width=\linewidth]{figures/validation__basson__p32_22.jpg}
\caption{Illustration (page 32)}
\end{figure}

egin{figure}[htbp]
\centering
\includegraphics[width=\linewidth]{figures/validation__basson__p32_23.jpg}
\caption{Illustration (page 32)}
\end{figure}

egin{figure}[htbp]
\centering
\includegraphics[width=\linewidth]{figures/validation__basson__p32_24.jpg}
\caption{Illustration (page 32)}
\end{figure}

egin{figure}[htbp]
\centering
\includegraphics[width=\linewidth]{figures/validation__basson__p32_25.jpg}
\caption{Illustration (page 32)}
\end{figure}

egin{figure}[htbp]
\centering
\includegraphics[width=\linewidth]{figures/validation__basson__p32_26.jpg}
\caption{Illustration (page 32)}
\end{figure}

egin{figure}[htbp]
\centering
\includegraphics[width=\linewidth]{figures/validation__basson__p32_27.png}
\caption{Illustration (page 32)}
\end{figure}

egin{figure}[htbp]
\centering
\includegraphics[width=\linewidth]{figures/validation__basson__p32_28.png}
\caption{Illustration (page 32)}
\end{figure}

egin{figure}[htbp]
\centering
\includegraphics[width=\linewidth]{figures/validation__basson__p32_29.png}
\caption{Illustration (page 32)}
\end{figure}

egin{figure}[htbp]
\centering
\includegraphics[width=\linewidth]{figures/validation__basson__p32_30.png}
\caption{Illustration (page 32)}
\end{figure}

egin{figure}[htbp]
\centering
\includegraphics[width=\linewidth]{figures/validation__basson__p32_31.png}
\caption{Illustration (page 32)}
\end{figure}

egin{figure}[htbp]
\centering
\includegraphics[width=\linewidth]{figures/validation__basson__p32_32.jpg}
\caption{Illustration (page 32)}
\end{figure}

egin{figure}[htbp]
\centering
\includegraphics[width=\linewidth]{figures/validation__basson__p32_33.png}
\caption{Illustration (page 32)}
\end{figure}

\section{Applications à l'orchestration}
X. Applications à l'Orchestration Prédiction de la fusion timbrale La distance formantique (Δ en Hz) entre deux instruments prédit le degré de fusion timbrale: Δ F1 (Hz)Type de fusion Effet acoustique Exemples orchestraux < 50 HzFusion quasi- totaleRésonances identiques, timbres indissociablesBasson + Vcl (3 Hz), Trb + Vcl (8 Hz) 50–100 HzMélange homogène Légère distinction perceptibleCor + Basson (45 Hz), Tuba + CB (49 Hz) 100–300 HzMélange avec distinctionDeux couleurs complémentaires Trp + Violon (290 Hz) 300–500 HzSéparation partielle Couleurs distinctes Htb + Violon (426 Hz) > 500 HzContraste intentionnelEffets dramatiques, couleurs séparéesTuba + Piccolo (2 409 Hz) Doublures archétypales par zone Zone formantique Doublure Compositeurs types Caractère Cluster /o/ 450–502 Hz Cor + Basson (Δ=45)Mozart, Brahms, Beethoven, MendelssohnTimbre pastoral, plénitude Cluster /o/ 450–502 HzTrombone + Violoncelle (Δ=8)Wagner (Tristan), Bruckner symphoniesPuissance grave, homogénéité Cluster /o/ 450–502 HzBasson + Violoncelle (Δ=3)Tous compositeurs classiques/romantiquesFusion légendaire Zone /a/ 800–1 250 HzTrompette + Violon (Δ=290)Bach (Brandenburg), fanfares baroquesProjection maximale Zone /a/ 800–1 250 HzHautbois + Violon (Δ=426)Vivaldi, Mozart, baroque Clarté, médiums Fondations /u/ 200–350 HzTuba + Contrebasse (Δ=49)Tous orchestres depuis WagnerAncrage harmonique stable Complémentarité spectrale Au-delà de la fusion par convergence, certaines doublures exploitent la complémentarité spectrale: •Clarinette (harmoniques impairs) + Basson (harmoniques pairs) = spectre complet. •Flûte (pas de formant fixe, timbre pur) + Hautbois (formant fixe) = clarté + présence. •Sourdine comme transposition timbrale: abaisse F1 de 30–50 \%, déplace d'une à deux catégories vocaliques vers le grave.
\par\medskip

egin{figure}[htbp]
\centering
\includegraphics[width=\linewidth]{figures/applications__basson__p33_01.jpg}
\caption{Illustration (page 33)}
\end{figure}

egin{figure}[htbp]
\centering
\includegraphics[width=\linewidth]{figures/applications__basson__p33_02.jpg}
\caption{Illustration (page 33)}
\end{figure}

egin{figure}[htbp]
\centering
\includegraphics[width=\linewidth]{figures/applications__basson__p33_03.png}
\caption{Illustration (page 33)}
\end{figure}

egin{figure}[htbp]
\centering
\includegraphics[width=\linewidth]{figures/applications__basson__p33_04.jpg}
\caption{Illustration (page 33)}
\end{figure}

egin{figure}[htbp]
\centering
\includegraphics[width=\linewidth]{figures/applications__basson__p33_05.jpg}
\caption{Illustration (page 33)}
\end{figure}

egin{figure}[htbp]
\centering
\includegraphics[width=\linewidth]{figures/applications__basson__p33_06.jpg}
\caption{Illustration (page 33)}
\end{figure}

egin{figure}[htbp]
\centering
\includegraphics[width=\linewidth]{figures/applications__basson__p33_07.jpg}
\caption{Illustration (page 33)}
\end{figure}

egin{figure}[htbp]
\centering
\includegraphics[width=\linewidth]{figures/applications__basson__p33_08.png}
\caption{Illustration (page 33)}
\end{figure}

egin{figure}[htbp]
\centering
\includegraphics[width=\linewidth]{figures/applications__basson__p33_09.jpg}
\caption{Illustration (page 33)}
\end{figure}

egin{figure}[htbp]
\centering
\includegraphics[width=\linewidth]{figures/applications__basson__p33_10.jpg}
\caption{Illustration (page 33)}
\end{figure}

egin{figure}[htbp]
\centering
\includegraphics[width=\linewidth]{figures/applications__basson__p33_11.png}
\caption{Illustration (page 33)}
\end{figure}

egin{figure}[htbp]
\centering
\includegraphics[width=\linewidth]{figures/applications__basson__p33_12.png}
\caption{Illustration (page 33)}
\end{figure}

egin{figure}[htbp]
\centering
\includegraphics[width=\linewidth]{figures/applications__basson__p33_13.png}
\caption{Illustration (page 33)}
\end{figure}

egin{figure}[htbp]
\centering
\includegraphics[width=\linewidth]{figures/applications__basson__p33_14.png}
\caption{Illustration (page 33)}
\end{figure}

egin{figure}[htbp]
\centering
\includegraphics[width=\linewidth]{figures/applications__basson__p33_15.png}
\caption{Illustration (page 33)}
\end{figure}

egin{figure}[htbp]
\centering
\includegraphics[width=\linewidth]{figures/applications__basson__p33_16.png}
\caption{Illustration (page 33)}
\end{figure}

egin{figure}[htbp]
\centering
\includegraphics[width=\linewidth]{figures/applications__basson__p33_17.png}
\caption{Illustration (page 33)}
\end{figure}

egin{figure}[htbp]
\centering
\includegraphics[width=\linewidth]{figures/applications__basson__p33_18.png}
\caption{Illustration (page 33)}
\end{figure}

egin{figure}[htbp]
\centering
\includegraphics[width=\linewidth]{figures/applications__basson__p33_19.png}
\caption{Illustration (page 33)}
\end{figure}

egin{figure}[htbp]
\centering
\includegraphics[width=\linewidth]{figures/applications__basson__p33_20.jpg}
\caption{Illustration (page 33)}
\end{figure}

egin{figure}[htbp]
\centering
\includegraphics[width=\linewidth]{figures/applications__basson__p33_21.jpg}
\caption{Illustration (page 33)}
\end{figure}

egin{figure}[htbp]
\centering
\includegraphics[width=\linewidth]{figures/applications__basson__p33_22.jpg}
\caption{Illustration (page 33)}
\end{figure}

egin{figure}[htbp]
\centering
\includegraphics[width=\linewidth]{figures/applications__basson__p33_23.jpg}
\caption{Illustration (page 33)}
\end{figure}

egin{figure}[htbp]
\centering
\includegraphics[width=\linewidth]{figures/applications__basson__p33_24.jpg}
\caption{Illustration (page 33)}
\end{figure}

egin{figure}[htbp]
\centering
\includegraphics[width=\linewidth]{figures/applications__basson__p33_25.jpg}
\caption{Illustration (page 33)}
\end{figure}

egin{figure}[htbp]
\centering
\includegraphics[width=\linewidth]{figures/applications__basson__p33_26.jpg}
\caption{Illustration (page 33)}
\end{figure}

egin{figure}[htbp]
\centering
\includegraphics[width=\linewidth]{figures/applications__basson__p33_27.png}
\caption{Illustration (page 33)}
\end{figure}

egin{figure}[htbp]
\centering
\includegraphics[width=\linewidth]{figures/applications__basson__p33_28.png}
\caption{Illustration (page 33)}
\end{figure}

egin{figure}[htbp]
\centering
\includegraphics[width=\linewidth]{figures/applications__basson__p33_29.png}
\caption{Illustration (page 33)}
\end{figure}

egin{figure}[htbp]
\centering
\includegraphics[width=\linewidth]{figures/applications__basson__p33_30.png}
\caption{Illustration (page 33)}
\end{figure}

egin{figure}[htbp]
\centering
\includegraphics[width=\linewidth]{figures/applications__basson__p33_31.png}
\caption{Illustration (page 33)}
\end{figure}

egin{figure}[htbp]
\centering
\includegraphics[width=\linewidth]{figures/applications__basson__p33_32.jpg}
\caption{Illustration (page 33)}
\end{figure}

egin{figure}[htbp]
\centering
\includegraphics[width=\linewidth]{figures/applications__basson__p33_33.png}
\caption{Illustration (page 33)}
\end{figure}

\section{Tableau des doublures vérifiées}
\subsection{Violoncelle + Basson (F1–F1) 4995023✓✓Timbres frères — convergence la plus}
XI. Tableau des Doublures Vérifiées Trié par écart croissant. Lignes bleues = doublures nouvelles (Yan\_Adds). ✓✓ < 60 Hz. Doublure HzHzΔQ.Signification orchestrale Violoncelle + Basson (F1–F1) 4995023✓✓Timbres frères — convergence la plus serrée Violoncelle + Trombone (F1–F1) 4994918✓✓Fusion graves nobles Tuba F2 + Basson F1 4945029✓✓Pont fondations bois-cuivres Trombone + Basson (F1–F1) 49150212✓✓Pont bois-cuivres graves Contrebasson + Trombone (F1) — N47849113✓✓Cluster graves renforcé CB Tuba + Cor (F1–F1) — N 47145714✓✓Fusion cuivres profonds Cl. basse + Trb. basse (F1) — N 90989415✓✓Convergence zone 'a (ah)' CB Tuba + Trombone (F1) — N 47149120✓✓Cuivres graves Contrebasson + Violoncelle — N 47849921✓✓Fusion graves bois-cordes Contrebasson + Basson (F1) — N 47850224✓✓Famille bassons Cor anglais + Cl. Sib (F1) — N 1 0451 01629✓✓Médiums bois chaleureux Cor F2 + Basson F2 68571429✓✓Médiums bois-cuivres CB Tuba + Basson (F1) — N 47150231✓✓Cuivres-bois graves Cor + Trombone (F1–F1) 45749133✓✓Homogénéité cuivres Vcl-Ens sord. + Cor (F1) — N 42345734✓✓Cordes sord. + cuivres Cor + Violoncelle (F1–F1) 45749942✓✓Chant lyrique Cor + Basson (F1–F1) 45750245✓✓Transition bois-cuivres classique Tuba + Contrebasse (F1) 24920049✓✓Fondations orchestrales Cl. contrebasse + Cl. basse — N 96390954✓✓Famille clarinettes graves
\par\medskip

egin{figure}[htbp]
\centering
\includegraphics[width=\linewidth]{figures/doublures__violon__p34_01.jpg}
\caption{Illustration (page 34)}
\end{figure}

egin{figure}[htbp]
\centering
\includegraphics[width=\linewidth]{figures/doublures__violon__p34_02.jpg}
\caption{Illustration (page 34)}
\end{figure}

egin{figure}[htbp]
\centering
\includegraphics[width=\linewidth]{figures/doublures__violon__p34_03.png}
\caption{Illustration (page 34)}
\end{figure}

egin{figure}[htbp]
\centering
\includegraphics[width=\linewidth]{figures/doublures__violon__p34_04.jpg}
\caption{Illustration (page 34)}
\end{figure}

egin{figure}[htbp]
\centering
\includegraphics[width=\linewidth]{figures/doublures__violon__p34_05.jpg}
\caption{Illustration (page 34)}
\end{figure}

egin{figure}[htbp]
\centering
\includegraphics[width=\linewidth]{figures/doublures__violon__p34_06.jpg}
\caption{Illustration (page 34)}
\end{figure}

egin{figure}[htbp]
\centering
\includegraphics[width=\linewidth]{figures/doublures__violon__p34_07.jpg}
\caption{Illustration (page 34)}
\end{figure}

egin{figure}[htbp]
\centering
\includegraphics[width=\linewidth]{figures/doublures__violon__p34_08.png}
\caption{Illustration (page 34)}
\end{figure}

egin{figure}[htbp]
\centering
\includegraphics[width=\linewidth]{figures/doublures__violon__p34_09.jpg}
\caption{Illustration (page 34)}
\end{figure}

egin{figure}[htbp]
\centering
\includegraphics[width=\linewidth]{figures/doublures__violon__p34_10.jpg}
\caption{Illustration (page 34)}
\end{figure}

egin{figure}[htbp]
\centering
\includegraphics[width=\linewidth]{figures/doublures__violon__p34_11.png}
\caption{Illustration (page 34)}
\end{figure}

egin{figure}[htbp]
\centering
\includegraphics[width=\linewidth]{figures/doublures__violon__p34_12.png}
\caption{Illustration (page 34)}
\end{figure}

egin{figure}[htbp]
\centering
\includegraphics[width=\linewidth]{figures/doublures__violon__p34_13.png}
\caption{Illustration (page 34)}
\end{figure}

egin{figure}[htbp]
\centering
\includegraphics[width=\linewidth]{figures/doublures__violon__p34_14.png}
\caption{Illustration (page 34)}
\end{figure}

egin{figure}[htbp]
\centering
\includegraphics[width=\linewidth]{figures/doublures__violon__p34_15.png}
\caption{Illustration (page 34)}
\end{figure}

egin{figure}[htbp]
\centering
\includegraphics[width=\linewidth]{figures/doublures__violon__p34_16.png}
\caption{Illustration (page 34)}
\end{figure}

egin{figure}[htbp]
\centering
\includegraphics[width=\linewidth]{figures/doublures__violon__p34_17.png}
\caption{Illustration (page 34)}
\end{figure}

egin{figure}[htbp]
\centering
\includegraphics[width=\linewidth]{figures/doublures__violon__p34_18.png}
\caption{Illustration (page 34)}
\end{figure}

egin{figure}[htbp]
\centering
\includegraphics[width=\linewidth]{figures/doublures__violon__p34_19.png}
\caption{Illustration (page 34)}
\end{figure}

egin{figure}[htbp]
\centering
\includegraphics[width=\linewidth]{figures/doublures__violon__p34_20.jpg}
\caption{Illustration (page 34)}
\end{figure}

egin{figure}[htbp]
\centering
\includegraphics[width=\linewidth]{figures/doublures__violon__p34_21.jpg}
\caption{Illustration (page 34)}
\end{figure}

egin{figure}[htbp]
\centering
\includegraphics[width=\linewidth]{figures/doublures__violon__p34_22.jpg}
\caption{Illustration (page 34)}
\end{figure}

egin{figure}[htbp]
\centering
\includegraphics[width=\linewidth]{figures/doublures__violon__p34_23.jpg}
\caption{Illustration (page 34)}
\end{figure}

egin{figure}[htbp]
\centering
\includegraphics[width=\linewidth]{figures/doublures__violon__p34_24.jpg}
\caption{Illustration (page 34)}
\end{figure}

egin{figure}[htbp]
\centering
\includegraphics[width=\linewidth]{figures/doublures__violon__p34_25.jpg}
\caption{Illustration (page 34)}
\end{figure}

egin{figure}[htbp]
\centering
\includegraphics[width=\linewidth]{figures/doublures__violon__p34_26.jpg}
\caption{Illustration (page 34)}
\end{figure}

egin{figure}[htbp]
\centering
\includegraphics[width=\linewidth]{figures/doublures__violon__p34_27.png}
\caption{Illustration (page 34)}
\end{figure}

egin{figure}[htbp]
\centering
\includegraphics[width=\linewidth]{figures/doublures__violon__p34_28.png}
\caption{Illustration (page 34)}
\end{figure}

egin{figure}[htbp]
\centering
\includegraphics[width=\linewidth]{figures/doublures__violon__p34_29.png}
\caption{Illustration (page 34)}
\end{figure}

egin{figure}[htbp]
\centering
\includegraphics[width=\linewidth]{figures/doublures__violon__p34_30.png}
\caption{Illustration (page 34)}
\end{figure}

egin{figure}[htbp]
\centering
\includegraphics[width=\linewidth]{figures/doublures__violon__p34_31.png}
\caption{Illustration (page 34)}
\end{figure}

egin{figure}[htbp]
\centering
\includegraphics[width=\linewidth]{figures/doublures__violon__p34_32.jpg}
\caption{Illustration (page 34)}
\end{figure}

egin{figure}[htbp]
\centering
\includegraphics[width=\linewidth]{figures/doublures__violon__p34_33.png}
\caption{Illustration (page 34)}
\end{figure}

\section{Principes d'orchestration acoustique}
XII. Principes d'Orchestration Acoustique 1. Convergence formantique (fusion) Les doublures les plus efficaces exploitent des F1 convergents: le cluster 450–502 Hz (6 instruments), les fondations 200–350 Hz, la zone 'a' 890–1 050 Hz. Fusion quasi-totale quand Δ < 50 Hz. 2. Complémentarité spectrale (enrichissement) Clarinette (harmoniques impairs) + Basson (harmoniques pairs) = spectre complet. Flûte (timbre pur) + Hautbois (formant fixe) = clarté + présence. 3. Effet de section (compression formantique) Les ensembles de cordes 'compriment' les formants vers la zone médiane (violon section 1 556 Hz vs soliste 1 034 Hz). L'orchestre est naturellement plus homogène que la somme de ses parties. 4. La sourdine comme transposition timbrale La sourdine abaisse F1 de 30 à 50 \%, déplaçant systématiquement le timbre d'une à deux catégories vocaliques vers le grave. Outil d'orchestration puissant pour modifier la couleur sans changer les notes. 5. Familles formantiques transversales Les données révèlent des 'familles formantiques' transversales: clarinettes graves (basse 909, contrebasse 963, Δ=54 Hz); flûtes graves (basse 737, contrebasse 746, Δ=9 Hz); bassons (basson 502, contrebasson 478, Δ=24 Hz). Les instruments à l'octave inférieure conservent souvent des formants très proches. 6. Masquage évité Les grandes doublures évitent le masquage fréquentiel où un instrument couvrirait les formants essentiels de l'autre. La carte de compatibilité (Figure 4) identifie les combinaisons à éviter. Conclusion L'analyse formantique de 29 instruments (4 037 échantillons, quatre sources, taux de validation 93 \%) révèle que les doublures orchestrales classiques reposent sur des convergences mesurables et robustes. Le cluster 450–502 Hz (Cor–CB Tuba–Contrebasson–Trombone–Violoncelle–Basson) et la zone de fondation 200–350 Hz (Tuba–Contrebasses) constituent les piliers acoustiques de l'orchestre. Ces convergences expliquent acoustiquement des pratiques empiriques développées sur trois siècles. Deux découvertes majeures émergent des nouvelles données d'ensemble: (1) l'effet de section comprime les formants vers la zone médiane, créant naturellement l'homogénéité orchestrale; (2) la sourdine opère une transposition timbrale systématique de une à deux voyelles vers le grave. L'intuition des grands orchestrateurs se trouve validée par la mesure spectrale directe. L'art de l'orchestration repose sur des fondements acoustiques profonds et universels — les mêmes contraintes physiques qui gouvernent la voix humaine régissent l'ensemble des instruments d'orchestre.
\par\medskip

egin{figure}[htbp]
\centering
\includegraphics[width=\linewidth]{figures/principes__basson__p35_01.jpg}
\caption{Illustration (page 35)}
\end{figure}

egin{figure}[htbp]
\centering
\includegraphics[width=\linewidth]{figures/principes__basson__p35_02.jpg}
\caption{Illustration (page 35)}
\end{figure}

egin{figure}[htbp]
\centering
\includegraphics[width=\linewidth]{figures/principes__basson__p35_03.png}
\caption{Illustration (page 35)}
\end{figure}

egin{figure}[htbp]
\centering
\includegraphics[width=\linewidth]{figures/principes__basson__p35_04.jpg}
\caption{Illustration (page 35)}
\end{figure}

egin{figure}[htbp]
\centering
\includegraphics[width=\linewidth]{figures/principes__basson__p35_05.jpg}
\caption{Illustration (page 35)}
\end{figure}

egin{figure}[htbp]
\centering
\includegraphics[width=\linewidth]{figures/principes__basson__p35_06.jpg}
\caption{Illustration (page 35)}
\end{figure}

egin{figure}[htbp]
\centering
\includegraphics[width=\linewidth]{figures/principes__basson__p35_07.jpg}
\caption{Illustration (page 35)}
\end{figure}

egin{figure}[htbp]
\centering
\includegraphics[width=\linewidth]{figures/principes__basson__p35_08.png}
\caption{Illustration (page 35)}
\end{figure}

egin{figure}[htbp]
\centering
\includegraphics[width=\linewidth]{figures/principes__basson__p35_09.jpg}
\caption{Illustration (page 35)}
\end{figure}

egin{figure}[htbp]
\centering
\includegraphics[width=\linewidth]{figures/principes__basson__p35_10.jpg}
\caption{Illustration (page 35)}
\end{figure}

egin{figure}[htbp]
\centering
\includegraphics[width=\linewidth]{figures/principes__basson__p35_11.png}
\caption{Illustration (page 35)}
\end{figure}

egin{figure}[htbp]
\centering
\includegraphics[width=\linewidth]{figures/principes__basson__p35_12.png}
\caption{Illustration (page 35)}
\end{figure}

egin{figure}[htbp]
\centering
\includegraphics[width=\linewidth]{figures/principes__basson__p35_13.png}
\caption{Illustration (page 35)}
\end{figure}

egin{figure}[htbp]
\centering
\includegraphics[width=\linewidth]{figures/principes__basson__p35_14.png}
\caption{Illustration (page 35)}
\end{figure}

egin{figure}[htbp]
\centering
\includegraphics[width=\linewidth]{figures/principes__basson__p35_15.png}
\caption{Illustration (page 35)}
\end{figure}

egin{figure}[htbp]
\centering
\includegraphics[width=\linewidth]{figures/principes__basson__p35_16.png}
\caption{Illustration (page 35)}
\end{figure}

egin{figure}[htbp]
\centering
\includegraphics[width=\linewidth]{figures/principes__basson__p35_17.png}
\caption{Illustration (page 35)}
\end{figure}

egin{figure}[htbp]
\centering
\includegraphics[width=\linewidth]{figures/principes__basson__p35_18.png}
\caption{Illustration (page 35)}
\end{figure}

egin{figure}[htbp]
\centering
\includegraphics[width=\linewidth]{figures/principes__basson__p35_19.png}
\caption{Illustration (page 35)}
\end{figure}

egin{figure}[htbp]
\centering
\includegraphics[width=\linewidth]{figures/principes__basson__p35_20.jpg}
\caption{Illustration (page 35)}
\end{figure}

egin{figure}[htbp]
\centering
\includegraphics[width=\linewidth]{figures/principes__basson__p35_21.jpg}
\caption{Illustration (page 35)}
\end{figure}

egin{figure}[htbp]
\centering
\includegraphics[width=\linewidth]{figures/principes__basson__p35_22.jpg}
\caption{Illustration (page 35)}
\end{figure}

egin{figure}[htbp]
\centering
\includegraphics[width=\linewidth]{figures/principes__basson__p35_23.jpg}
\caption{Illustration (page 35)}
\end{figure}

egin{figure}[htbp]
\centering
\includegraphics[width=\linewidth]{figures/principes__basson__p35_24.jpg}
\caption{Illustration (page 35)}
\end{figure}

egin{figure}[htbp]
\centering
\includegraphics[width=\linewidth]{figures/principes__basson__p35_25.jpg}
\caption{Illustration (page 35)}
\end{figure}

egin{figure}[htbp]
\centering
\includegraphics[width=\linewidth]{figures/principes__basson__p35_26.jpg}
\caption{Illustration (page 35)}
\end{figure}

egin{figure}[htbp]
\centering
\includegraphics[width=\linewidth]{figures/principes__basson__p35_27.png}
\caption{Illustration (page 35)}
\end{figure}

egin{figure}[htbp]
\centering
\includegraphics[width=\linewidth]{figures/principes__basson__p35_28.png}
\caption{Illustration (page 35)}
\end{figure}

egin{figure}[htbp]
\centering
\includegraphics[width=\linewidth]{figures/principes__basson__p35_29.png}
\caption{Illustration (page 35)}
\end{figure}

egin{figure}[htbp]
\centering
\includegraphics[width=\linewidth]{figures/principes__basson__p35_30.png}
\caption{Illustration (page 35)}
\end{figure}

egin{figure}[htbp]
\centering
\includegraphics[width=\linewidth]{figures/principes__basson__p35_31.png}
\caption{Illustration (page 35)}
\end{figure}

egin{figure}[htbp]
\centering
\includegraphics[width=\linewidth]{figures/principes__basson__p35_32.jpg}
\caption{Illustration (page 35)}
\end{figure}

egin{figure}[htbp]
\centering
\includegraphics[width=\linewidth]{figures/principes__basson__p35_33.png}
\caption{Illustration (page 35)}
\end{figure}
Pour l'analyse musicale: expliquer les choix de doublures des compositeurs, identifier les 'accords de timbres', comprendre l'évolution des pratiques orchestrales. Pour la composition: optimiser les doublures selon l'effet désiré, prédire les résultats timbraux, créer des progressions de couleurs cohérentes.
\par\medskip

egin{figure}[htbp]
\centering
\includegraphics[width=\linewidth]{figures/principes__basson__p36_01.jpg}
\caption{Illustration (page 36)}
\end{figure}

egin{figure}[htbp]
\centering
\includegraphics[width=\linewidth]{figures/principes__basson__p36_02.jpg}
\caption{Illustration (page 36)}
\end{figure}

egin{figure}[htbp]
\centering
\includegraphics[width=\linewidth]{figures/principes__basson__p36_03.png}
\caption{Illustration (page 36)}
\end{figure}

egin{figure}[htbp]
\centering
\includegraphics[width=\linewidth]{figures/principes__basson__p36_04.jpg}
\caption{Illustration (page 36)}
\end{figure}

egin{figure}[htbp]
\centering
\includegraphics[width=\linewidth]{figures/principes__basson__p36_05.jpg}
\caption{Illustration (page 36)}
\end{figure}

egin{figure}[htbp]
\centering
\includegraphics[width=\linewidth]{figures/principes__basson__p36_06.jpg}
\caption{Illustration (page 36)}
\end{figure}

egin{figure}[htbp]
\centering
\includegraphics[width=\linewidth]{figures/principes__basson__p36_07.jpg}
\caption{Illustration (page 36)}
\end{figure}

egin{figure}[htbp]
\centering
\includegraphics[width=\linewidth]{figures/principes__basson__p36_08.png}
\caption{Illustration (page 36)}
\end{figure}

egin{figure}[htbp]
\centering
\includegraphics[width=\linewidth]{figures/principes__basson__p36_09.jpg}
\caption{Illustration (page 36)}
\end{figure}

egin{figure}[htbp]
\centering
\includegraphics[width=\linewidth]{figures/principes__basson__p36_10.jpg}
\caption{Illustration (page 36)}
\end{figure}

egin{figure}[htbp]
\centering
\includegraphics[width=\linewidth]{figures/principes__basson__p36_11.png}
\caption{Illustration (page 36)}
\end{figure}

egin{figure}[htbp]
\centering
\includegraphics[width=\linewidth]{figures/principes__basson__p36_12.png}
\caption{Illustration (page 36)}
\end{figure}

egin{figure}[htbp]
\centering
\includegraphics[width=\linewidth]{figures/principes__basson__p36_13.png}
\caption{Illustration (page 36)}
\end{figure}

egin{figure}[htbp]
\centering
\includegraphics[width=\linewidth]{figures/principes__basson__p36_14.png}
\caption{Illustration (page 36)}
\end{figure}

egin{figure}[htbp]
\centering
\includegraphics[width=\linewidth]{figures/principes__basson__p36_15.png}
\caption{Illustration (page 36)}
\end{figure}

egin{figure}[htbp]
\centering
\includegraphics[width=\linewidth]{figures/principes__basson__p36_16.png}
\caption{Illustration (page 36)}
\end{figure}

egin{figure}[htbp]
\centering
\includegraphics[width=\linewidth]{figures/principes__basson__p36_17.png}
\caption{Illustration (page 36)}
\end{figure}

egin{figure}[htbp]
\centering
\includegraphics[width=\linewidth]{figures/principes__basson__p36_18.png}
\caption{Illustration (page 36)}
\end{figure}

egin{figure}[htbp]
\centering
\includegraphics[width=\linewidth]{figures/principes__basson__p36_19.png}
\caption{Illustration (page 36)}
\end{figure}

egin{figure}[htbp]
\centering
\includegraphics[width=\linewidth]{figures/principes__basson__p36_20.jpg}
\caption{Illustration (page 36)}
\end{figure}

egin{figure}[htbp]
\centering
\includegraphics[width=\linewidth]{figures/principes__basson__p36_21.jpg}
\caption{Illustration (page 36)}
\end{figure}

egin{figure}[htbp]
\centering
\includegraphics[width=\linewidth]{figures/principes__basson__p36_22.jpg}
\caption{Illustration (page 36)}
\end{figure}

egin{figure}[htbp]
\centering
\includegraphics[width=\linewidth]{figures/principes__basson__p36_23.jpg}
\caption{Illustration (page 36)}
\end{figure}

egin{figure}[htbp]
\centering
\includegraphics[width=\linewidth]{figures/principes__basson__p36_24.jpg}
\caption{Illustration (page 36)}
\end{figure}

egin{figure}[htbp]
\centering
\includegraphics[width=\linewidth]{figures/principes__basson__p36_25.jpg}
\caption{Illustration (page 36)}
\end{figure}

egin{figure}[htbp]
\centering
\includegraphics[width=\linewidth]{figures/principes__basson__p36_26.jpg}
\caption{Illustration (page 36)}
\end{figure}

egin{figure}[htbp]
\centering
\includegraphics[width=\linewidth]{figures/principes__basson__p36_27.png}
\caption{Illustration (page 36)}
\end{figure}

egin{figure}[htbp]
\centering
\includegraphics[width=\linewidth]{figures/principes__basson__p36_28.png}
\caption{Illustration (page 36)}
\end{figure}

egin{figure}[htbp]
\centering
\includegraphics[width=\linewidth]{figures/principes__basson__p36_29.png}
\caption{Illustration (page 36)}
\end{figure}

egin{figure}[htbp]
\centering
\includegraphics[width=\linewidth]{figures/principes__basson__p36_30.png}
\caption{Illustration (page 36)}
\end{figure}

egin{figure}[htbp]
\centering
\includegraphics[width=\linewidth]{figures/principes__basson__p36_31.png}
\caption{Illustration (page 36)}
\end{figure}

egin{figure}[htbp]
\centering
\includegraphics[width=\linewidth]{figures/principes__basson__p36_32.jpg}
\caption{Illustration (page 36)}
\end{figure}

egin{figure}[htbp]
\centering
\includegraphics[width=\linewidth]{figures/principes__basson__p36_33.png}
\caption{Illustration (page 36)}
\end{figure}
Annexe: Sources et Références Bases de données Studio On Line (SOL2020), IRCAM 2 391 échantillons analysés · 12 instruments solistes · Enregistrements professionnels multi-dynamiques · Spectral envelopes (4 096 FFT, 44 100 Hz) Yan\_Adds (extension personnelle) 1 646 échantillons · 17 instruments supplémentaires: ensembles de cordes, bois auxiliaires, cuivres graves · Même méthode d'extraction specenv, modes ordinario Références académiques Auteur Titre Éditeur Description Backus, J. (1969)The Acoustical Foundations of MusicW.W. Norton \& Company, New YorkRéférence standard en acoustique musicale. Cuivres et bois principaux. Giesler, M. (1985)Instrumentation: Ein Hand- und LehrbuchBreitkopf \& Härtel, WiesbadenMesures détaillées, associations vocaliques explicites. 13 instruments couverts. Meyer, J. (2009)Acoustics and the Performance of Music (5e éd.)Springer-Verlag, New YorkApproche scientifique moderne. Chapitre 5: spectres et formants. Zones vocaliques définies. McCarty, J. (2003) Instrumental Vowel SpaceCCRMA, Stanford Univ.Graphique vowel space F1/F2. Analyse COLEA de 11 instruments. Disponible en ligne. Logiciels •Python 3.x + scipy.signal.find\_peaks •numpy (analyse numérique) •python-docx (génération documents) © 2026 — Référence Formantique des Instruments d'Orchestre Base: SOL2020 (IRCAM)  |  Validation: Giesler (1985), Backus (1969), Meyer (2009)
\par\medskip

egin{figure}[htbp]
\centering
\includegraphics[width=\linewidth]{figures/annexe__basson__p37_01.jpg}
\caption{Illustration (page 37)}
\end{figure}

egin{figure}[htbp]
\centering
\includegraphics[width=\linewidth]{figures/annexe__basson__p37_02.jpg}
\caption{Illustration (page 37)}
\end{figure}

egin{figure}[htbp]
\centering
\includegraphics[width=\linewidth]{figures/annexe__basson__p37_03.png}
\caption{Illustration (page 37)}
\end{figure}

egin{figure}[htbp]
\centering
\includegraphics[width=\linewidth]{figures/annexe__basson__p37_04.jpg}
\caption{Illustration (page 37)}
\end{figure}

egin{figure}[htbp]
\centering
\includegraphics[width=\linewidth]{figures/annexe__basson__p37_05.jpg}
\caption{Illustration (page 37)}
\end{figure}

egin{figure}[htbp]
\centering
\includegraphics[width=\linewidth]{figures/annexe__basson__p37_06.jpg}
\caption{Illustration (page 37)}
\end{figure}

egin{figure}[htbp]
\centering
\includegraphics[width=\linewidth]{figures/annexe__basson__p37_07.jpg}
\caption{Illustration (page 37)}
\end{figure}

egin{figure}[htbp]
\centering
\includegraphics[width=\linewidth]{figures/annexe__basson__p37_08.png}
\caption{Illustration (page 37)}
\end{figure}

egin{figure}[htbp]
\centering
\includegraphics[width=\linewidth]{figures/annexe__basson__p37_09.jpg}
\caption{Illustration (page 37)}
\end{figure}

egin{figure}[htbp]
\centering
\includegraphics[width=\linewidth]{figures/annexe__basson__p37_10.jpg}
\caption{Illustration (page 37)}
\end{figure}

egin{figure}[htbp]
\centering
\includegraphics[width=\linewidth]{figures/annexe__basson__p37_11.png}
\caption{Illustration (page 37)}
\end{figure}

egin{figure}[htbp]
\centering
\includegraphics[width=\linewidth]{figures/annexe__basson__p37_12.png}
\caption{Illustration (page 37)}
\end{figure}

egin{figure}[htbp]
\centering
\includegraphics[width=\linewidth]{figures/annexe__basson__p37_13.png}
\caption{Illustration (page 37)}
\end{figure}

egin{figure}[htbp]
\centering
\includegraphics[width=\linewidth]{figures/annexe__basson__p37_14.png}
\caption{Illustration (page 37)}
\end{figure}

egin{figure}[htbp]
\centering
\includegraphics[width=\linewidth]{figures/annexe__basson__p37_15.png}
\caption{Illustration (page 37)}
\end{figure}

egin{figure}[htbp]
\centering
\includegraphics[width=\linewidth]{figures/annexe__basson__p37_16.png}
\caption{Illustration (page 37)}
\end{figure}

egin{figure}[htbp]
\centering
\includegraphics[width=\linewidth]{figures/annexe__basson__p37_17.png}
\caption{Illustration (page 37)}
\end{figure}

egin{figure}[htbp]
\centering
\includegraphics[width=\linewidth]{figures/annexe__basson__p37_18.png}
\caption{Illustration (page 37)}
\end{figure}

egin{figure}[htbp]
\centering
\includegraphics[width=\linewidth]{figures/annexe__basson__p37_19.png}
\caption{Illustration (page 37)}
\end{figure}

egin{figure}[htbp]
\centering
\includegraphics[width=\linewidth]{figures/annexe__basson__p37_20.jpg}
\caption{Illustration (page 37)}
\end{figure}

egin{figure}[htbp]
\centering
\includegraphics[width=\linewidth]{figures/annexe__basson__p37_21.jpg}
\caption{Illustration (page 37)}
\end{figure}

egin{figure}[htbp]
\centering
\includegraphics[width=\linewidth]{figures/annexe__basson__p37_22.jpg}
\caption{Illustration (page 37)}
\end{figure}

egin{figure}[htbp]
\centering
\includegraphics[width=\linewidth]{figures/annexe__basson__p37_23.jpg}
\caption{Illustration (page 37)}
\end{figure}

egin{figure}[htbp]
\centering
\includegraphics[width=\linewidth]{figures/annexe__basson__p37_24.jpg}
\caption{Illustration (page 37)}
\end{figure}

egin{figure}[htbp]
\centering
\includegraphics[width=\linewidth]{figures/annexe__basson__p37_25.jpg}
\caption{Illustration (page 37)}
\end{figure}

egin{figure}[htbp]
\centering
\includegraphics[width=\linewidth]{figures/annexe__basson__p37_26.jpg}
\caption{Illustration (page 37)}
\end{figure}

egin{figure}[htbp]
\centering
\includegraphics[width=\linewidth]{figures/annexe__basson__p37_27.png}
\caption{Illustration (page 37)}
\end{figure}

egin{figure}[htbp]
\centering
\includegraphics[width=\linewidth]{figures/annexe__basson__p37_28.png}
\caption{Illustration (page 37)}
\end{figure}

egin{figure}[htbp]
\centering
\includegraphics[width=\linewidth]{figures/annexe__basson__p37_29.png}
\caption{Illustration (page 37)}
\end{figure}

egin{figure}[htbp]
\centering
\includegraphics[width=\linewidth]{figures/annexe__basson__p37_30.png}
\caption{Illustration (page 37)}
\end{figure}

egin{figure}[htbp]
\centering
\includegraphics[width=\linewidth]{figures/annexe__basson__p37_31.png}
\caption{Illustration (page 37)}
\end{figure}

egin{figure}[htbp]
\centering
\includegraphics[width=\linewidth]{figures/annexe__basson__p37_32.jpg}
\caption{Illustration (page 37)}
\end{figure}

egin{figure}[htbp]
\centering
\includegraphics[width=\linewidth]{figures/annexe__basson__p37_33.png}
\caption{Illustration (page 37)}
\end{figure}

\section*{SECTION II — Analyse spectrale complète}
SECTION II — ANALYSE SPECTRALE COMPLÈTE F1–F6 spectraux stricts · Toutes techniques · SOL2020 IRCAM + Yan\_Adds-Divers · Script v3 corrigé II.1 — Contexte méthodologique La Section I de ce document mesure le formant perceptuel dominant (Fp) — le pic spectral le plus proéminent énergétiquement — en jeu sustained uniquement. Cette Section II adopte une approche complémentaire: elle cartographie les formants spectraux stricts F1 à F6 (premier pic, deuxième pic, etc. par ordre de fréquence croissante) sur l'intégralité du corpus SOL2020, toutes techniques de jeu confondues. Pour 8 instruments (Trombone, Violoncelle, Saxophone alto, Trompette, Clarinette Sib, Hautbois, Flûte, Violon), le Fp de la Section I correspond au F2 spectral de cette section — le premier mode (F1 spectral) étant acoustiquement moins rayonné. Chaque fiche signale ce décalage par une ligne de rappel (→). II.2 — Table de correspondance Fp (Section I) / F1 spectral (Section II) InstrumentAncien Fp (Sec.I)F1 spectral (Sec.II)Fp centroïd eBande / Note Tuba basse 249 Hz226 Hz1239 Hz1000-2000 Fp centroïde capture la résonance du pavillon Cor en Fa 457 Hz388 Hz738 Hz600-1400 Fp centroïde dans zone Meyer (/ɔ/) Trombone ténor 491 Hz237 Hz1218 Hz1000-2000 Fp centroïde zone résonance pavillon Trompette en Do1324 Hz786 Hz1046 Hz600-1400 Fp centroïde capture brillance (σ÷10 vs F2) Basson 502 Hz495 Hz1079 Hz800-1600 Fp centroïde zone anche double Hautbois1460 Hz743 Hz1485 Hz1000-2000 Fp centroïde confirme
\par\medskip

egin{figure}[htbp]
\centering
\includegraphics[width=\linewidth]{figures/section_ii__trombone__p38_01.jpg}
\caption{Illustration (page 38)}
\end{figure}

egin{figure}[htbp]
\centering
\includegraphics[width=\linewidth]{figures/section_ii__trombone__p38_02.jpg}
\caption{Illustration (page 38)}
\end{figure}

egin{figure}[htbp]
\centering
\includegraphics[width=\linewidth]{figures/section_ii__trombone__p38_03.png}
\caption{Illustration (page 38)}
\end{figure}

egin{figure}[htbp]
\centering
\includegraphics[width=\linewidth]{figures/section_ii__trombone__p38_04.jpg}
\caption{Illustration (page 38)}
\end{figure}

egin{figure}[htbp]
\centering
\includegraphics[width=\linewidth]{figures/section_ii__trombone__p38_05.jpg}
\caption{Illustration (page 38)}
\end{figure}

egin{figure}[htbp]
\centering
\includegraphics[width=\linewidth]{figures/section_ii__trombone__p38_06.jpg}
\caption{Illustration (page 38)}
\end{figure}

egin{figure}[htbp]
\centering
\includegraphics[width=\linewidth]{figures/section_ii__trombone__p38_07.jpg}
\caption{Illustration (page 38)}
\end{figure}

egin{figure}[htbp]
\centering
\includegraphics[width=\linewidth]{figures/section_ii__trombone__p38_08.png}
\caption{Illustration (page 38)}
\end{figure}

egin{figure}[htbp]
\centering
\includegraphics[width=\linewidth]{figures/section_ii__trombone__p38_09.jpg}
\caption{Illustration (page 38)}
\end{figure}

egin{figure}[htbp]
\centering
\includegraphics[width=\linewidth]{figures/section_ii__trombone__p38_10.jpg}
\caption{Illustration (page 38)}
\end{figure}

egin{figure}[htbp]
\centering
\includegraphics[width=\linewidth]{figures/section_ii__trombone__p38_11.png}
\caption{Illustration (page 38)}
\end{figure}

egin{figure}[htbp]
\centering
\includegraphics[width=\linewidth]{figures/section_ii__trombone__p38_12.png}
\caption{Illustration (page 38)}
\end{figure}

egin{figure}[htbp]
\centering
\includegraphics[width=\linewidth]{figures/section_ii__trombone__p38_13.png}
\caption{Illustration (page 38)}
\end{figure}

egin{figure}[htbp]
\centering
\includegraphics[width=\linewidth]{figures/section_ii__trombone__p38_14.png}
\caption{Illustration (page 38)}
\end{figure}

egin{figure}[htbp]
\centering
\includegraphics[width=\linewidth]{figures/section_ii__trombone__p38_15.png}
\caption{Illustration (page 38)}
\end{figure}

egin{figure}[htbp]
\centering
\includegraphics[width=\linewidth]{figures/section_ii__trombone__p38_16.png}
\caption{Illustration (page 38)}
\end{figure}

egin{figure}[htbp]
\centering
\includegraphics[width=\linewidth]{figures/section_ii__trombone__p38_17.png}
\caption{Illustration (page 38)}
\end{figure}

egin{figure}[htbp]
\centering
\includegraphics[width=\linewidth]{figures/section_ii__trombone__p38_18.png}
\caption{Illustration (page 38)}
\end{figure}

egin{figure}[htbp]
\centering
\includegraphics[width=\linewidth]{figures/section_ii__trombone__p38_19.png}
\caption{Illustration (page 38)}
\end{figure}

egin{figure}[htbp]
\centering
\includegraphics[width=\linewidth]{figures/section_ii__trombone__p38_20.jpg}
\caption{Illustration (page 38)}
\end{figure}

egin{figure}[htbp]
\centering
\includegraphics[width=\linewidth]{figures/section_ii__trombone__p38_21.jpg}
\caption{Illustration (page 38)}
\end{figure}

egin{figure}[htbp]
\centering
\includegraphics[width=\linewidth]{figures/section_ii__trombone__p38_22.jpg}
\caption{Illustration (page 38)}
\end{figure}

egin{figure}[htbp]
\centering
\includegraphics[width=\linewidth]{figures/section_ii__trombone__p38_23.jpg}
\caption{Illustration (page 38)}
\end{figure}

egin{figure}[htbp]
\centering
\includegraphics[width=\linewidth]{figures/section_ii__trombone__p38_24.jpg}
\caption{Illustration (page 38)}
\end{figure}

egin{figure}[htbp]
\centering
\includegraphics[width=\linewidth]{figures/section_ii__trombone__p38_25.jpg}
\caption{Illustration (page 38)}
\end{figure}

egin{figure}[htbp]
\centering
\includegraphics[width=\linewidth]{figures/section_ii__trombone__p38_26.jpg}
\caption{Illustration (page 38)}
\end{figure}

egin{figure}[htbp]
\centering
\includegraphics[width=\linewidth]{figures/section_ii__trombone__p38_27.png}
\caption{Illustration (page 38)}
\end{figure}

egin{figure}[htbp]
\centering
\includegraphics[width=\linewidth]{figures/section_ii__trombone__p38_28.png}
\caption{Illustration (page 38)}
\end{figure}

egin{figure}[htbp]
\centering
\includegraphics[width=\linewidth]{figures/section_ii__trombone__p38_29.png}
\caption{Illustration (page 38)}
\end{figure}

egin{figure}[htbp]
\centering
\includegraphics[width=\linewidth]{figures/section_ii__trombone__p38_30.png}
\caption{Illustration (page 38)}
\end{figure}

egin{figure}[htbp]
\centering
\includegraphics[width=\linewidth]{figures/section_ii__trombone__p38_31.png}
\caption{Illustration (page 38)}
\end{figure}

egin{figure}[htbp]
\centering
\includegraphics[width=\linewidth]{figures/section_ii__trombone__p38_32.jpg}
\caption{Illustration (page 38)}
\end{figure}

egin{figure}[htbp]
\centering
\includegraphics[width=\linewidth]{figures/section_ii__trombone__p38_33.png}
\caption{Illustration (page 38)}
\end{figure}
\subsection{Clarinette Sib1016}
zone anche double Clarinette Sib1016 Hz463 Hz1412 Hz1000-2000 Fp centroïde stable (σ÷4.7 vs F2) Flûte1354 Hz743 Hz1535 Hz1000-2000 Fp centroïde zone embouchur e Saxophone alto 829 Hz398 Hz1440 Hz1000-2000 Fp centroïde zone corps/pavill on Violon1034 Hz506 Hz893 Hz600-1400 Fp centroïde = Meyer Hauptforma nt  (σ=92)✓ Alto 377 Hz369 Hz1221 Hz800-1600 Fp centroïde résonance caisse alto Violoncelle 499 Hz205 Hz1025 Hz600-1400 Fp centroïde résonance caisse (σ=70) Contrebasse 200 Hz172 Hz1324 Hz1000-2000 Fp centroïde résonance caisse haute Fond vert: Fp = F1 spectral (cohérence totale) · Fond jaune pâle: divergence modérée · Fond rose: divergence forte (Fp = F2 ou F3 spectral) II.3 — Fiches par instrument — Toutes techniques
\par\medskip

egin{figure}[htbp]
\centering
\includegraphics[width=\linewidth]{figures/section_ii__saxophone_alto__p39_01.jpg}
\caption{Illustration (page 39)}
\end{figure}

egin{figure}[htbp]
\centering
\includegraphics[width=\linewidth]{figures/section_ii__saxophone_alto__p39_02.jpg}
\caption{Illustration (page 39)}
\end{figure}

egin{figure}[htbp]
\centering
\includegraphics[width=\linewidth]{figures/section_ii__saxophone_alto__p39_03.png}
\caption{Illustration (page 39)}
\end{figure}

egin{figure}[htbp]
\centering
\includegraphics[width=\linewidth]{figures/section_ii__saxophone_alto__p39_04.jpg}
\caption{Illustration (page 39)}
\end{figure}

egin{figure}[htbp]
\centering
\includegraphics[width=\linewidth]{figures/section_ii__saxophone_alto__p39_05.jpg}
\caption{Illustration (page 39)}
\end{figure}

egin{figure}[htbp]
\centering
\includegraphics[width=\linewidth]{figures/section_ii__saxophone_alto__p39_06.jpg}
\caption{Illustration (page 39)}
\end{figure}

egin{figure}[htbp]
\centering
\includegraphics[width=\linewidth]{figures/section_ii__saxophone_alto__p39_07.jpg}
\caption{Illustration (page 39)}
\end{figure}

egin{figure}[htbp]
\centering
\includegraphics[width=\linewidth]{figures/section_ii__saxophone_alto__p39_08.png}
\caption{Illustration (page 39)}
\end{figure}

egin{figure}[htbp]
\centering
\includegraphics[width=\linewidth]{figures/section_ii__saxophone_alto__p39_09.jpg}
\caption{Illustration (page 39)}
\end{figure}

egin{figure}[htbp]
\centering
\includegraphics[width=\linewidth]{figures/section_ii__saxophone_alto__p39_10.jpg}
\caption{Illustration (page 39)}
\end{figure}

egin{figure}[htbp]
\centering
\includegraphics[width=\linewidth]{figures/section_ii__saxophone_alto__p39_11.png}
\caption{Illustration (page 39)}
\end{figure}

egin{figure}[htbp]
\centering
\includegraphics[width=\linewidth]{figures/section_ii__saxophone_alto__p39_12.png}
\caption{Illustration (page 39)}
\end{figure}

egin{figure}[htbp]
\centering
\includegraphics[width=\linewidth]{figures/section_ii__saxophone_alto__p39_13.png}
\caption{Illustration (page 39)}
\end{figure}

egin{figure}[htbp]
\centering
\includegraphics[width=\linewidth]{figures/section_ii__saxophone_alto__p39_14.png}
\caption{Illustration (page 39)}
\end{figure}

egin{figure}[htbp]
\centering
\includegraphics[width=\linewidth]{figures/section_ii__saxophone_alto__p39_15.png}
\caption{Illustration (page 39)}
\end{figure}

egin{figure}[htbp]
\centering
\includegraphics[width=\linewidth]{figures/section_ii__saxophone_alto__p39_16.png}
\caption{Illustration (page 39)}
\end{figure}

egin{figure}[htbp]
\centering
\includegraphics[width=\linewidth]{figures/section_ii__saxophone_alto__p39_17.png}
\caption{Illustration (page 39)}
\end{figure}

egin{figure}[htbp]
\centering
\includegraphics[width=\linewidth]{figures/section_ii__saxophone_alto__p39_18.png}
\caption{Illustration (page 39)}
\end{figure}

egin{figure}[htbp]
\centering
\includegraphics[width=\linewidth]{figures/section_ii__saxophone_alto__p39_19.png}
\caption{Illustration (page 39)}
\end{figure}

egin{figure}[htbp]
\centering
\includegraphics[width=\linewidth]{figures/section_ii__saxophone_alto__p39_20.jpg}
\caption{Illustration (page 39)}
\end{figure}

egin{figure}[htbp]
\centering
\includegraphics[width=\linewidth]{figures/section_ii__saxophone_alto__p39_21.jpg}
\caption{Illustration (page 39)}
\end{figure}

egin{figure}[htbp]
\centering
\includegraphics[width=\linewidth]{figures/section_ii__saxophone_alto__p39_22.jpg}
\caption{Illustration (page 39)}
\end{figure}

egin{figure}[htbp]
\centering
\includegraphics[width=\linewidth]{figures/section_ii__saxophone_alto__p39_23.jpg}
\caption{Illustration (page 39)}
\end{figure}

egin{figure}[htbp]
\centering
\includegraphics[width=\linewidth]{figures/section_ii__saxophone_alto__p39_24.jpg}
\caption{Illustration (page 39)}
\end{figure}

egin{figure}[htbp]
\centering
\includegraphics[width=\linewidth]{figures/section_ii__saxophone_alto__p39_25.jpg}
\caption{Illustration (page 39)}
\end{figure}

egin{figure}[htbp]
\centering
\includegraphics[width=\linewidth]{figures/section_ii__saxophone_alto__p39_26.jpg}
\caption{Illustration (page 39)}
\end{figure}

egin{figure}[htbp]
\centering
\includegraphics[width=\linewidth]{figures/section_ii__saxophone_alto__p39_27.png}
\caption{Illustration (page 39)}
\end{figure}

egin{figure}[htbp]
\centering
\includegraphics[width=\linewidth]{figures/section_ii__saxophone_alto__p39_28.png}
\caption{Illustration (page 39)}
\end{figure}

egin{figure}[htbp]
\centering
\includegraphics[width=\linewidth]{figures/section_ii__saxophone_alto__p39_29.png}
\caption{Illustration (page 39)}
\end{figure}

egin{figure}[htbp]
\centering
\includegraphics[width=\linewidth]{figures/section_ii__saxophone_alto__p39_30.png}
\caption{Illustration (page 39)}
\end{figure}

egin{figure}[htbp]
\centering
\includegraphics[width=\linewidth]{figures/section_ii__saxophone_alto__p39_31.png}
\caption{Illustration (page 39)}
\end{figure}

egin{figure}[htbp]
\centering
\includegraphics[width=\linewidth]{figures/section_ii__saxophone_alto__p39_32.jpg}
\caption{Illustration (page 39)}
\end{figure}

egin{figure}[htbp]
\centering
\includegraphics[width=\linewidth]{figures/section_ii__saxophone_alto__p39_33.png}
\caption{Illustration (page 39)}
\end{figure}
\subsection{Tuba basse   [872 échantillons]}
Famille: Cuivres Tuba basse   [872 échantillons] TechniqueNF1F2F3F4F5F6Fiab. Ordinario108226452980129215501863●●● Crescendo99226456955123115071726●●● Staccato662154521001128114751561●●● Trille72215431845———●●● Flutter / Flatterzunge1072194529911123——●●● Glissando312164317431066——●●● Multiphoniques3215452————●●○ Slap4420761673710161285—●●● Brassy / Cuivré3291592883118414751766●●○ Bisbigliando36215431958———●●● Son éolien11837971184148620244877●○○ Buzz43129581195126015611863●●○ Autres1822234941020129615901842●●● Fond vert = ordinario · Fond jaune = F1 dans cluster /o/ 420–550 Hz · Fond violet = technique étendue · Fiabilité: ●●● ≥0.55 / ●●○ 0.35–0.54 / ●○○ <0.35 Cor en Fa   [961 échantillons] TechniqueNF1F2F3F4F5F6Fiab. Ordinario1343886891012137818632369●●○ Crescendo1314447431678219827563374●●○ Staccato414637321152158318632218●●○ Trille56421738105813181389—●●○ Flutter / Flatterzunge1534157891232179921212470●●○ Slap444527431098180918522207●●○ Brassy / Cuivré434747541863221831443962●●○ Autres2234267361133188023312724●●○ Fond vert = ordinario · Fond jaune = F1 dans cluster /o/ 420–550 Hz · Fond violet = technique étendue · Fiabilité: ●●● ≥0.55 / ●●○ 0.35–0.54 / ●○○ <0.35 Trombone ténor   [671 échantillons] → F1 spectral = 237 Hz (/u/) · Fp (Section I) = 491 Hz = F2 spectral TechniqueNF1F2F3F4F5F6Fiab. Ordinario117237506775128117441970●●○ Crescendo103312588889138918012145●●○ Staccato37226495775129218632175●●○ Flutter / Flatterzunge109250541822132616522145●●○ Glissando3021548474312701281—●●● Slap12194506764131421322153●●○ Brassy / Cuivré314957861270156118632358●●○ Autres123367640934138217282146●●○ Fond vert = ordinario · Fond jaune = F1 dans cluster /o/ 420–550 Hz · Fond violet = technique étendue · Fiabilité: ●●● ≥0.55 / ●●○ 0.35–0.54 / ●○○ <0.35
\par\medskip

egin{figure}[htbp]
\centering
\includegraphics[width=\linewidth]{figures/section_ii__tuba__p40_01.jpg}
\caption{Illustration (page 40)}
\end{figure}

egin{figure}[htbp]
\centering
\includegraphics[width=\linewidth]{figures/section_ii__tuba__p40_02.jpg}
\caption{Illustration (page 40)}
\end{figure}

egin{figure}[htbp]
\centering
\includegraphics[width=\linewidth]{figures/section_ii__tuba__p40_03.png}
\caption{Illustration (page 40)}
\end{figure}

egin{figure}[htbp]
\centering
\includegraphics[width=\linewidth]{figures/section_ii__tuba__p40_04.jpg}
\caption{Illustration (page 40)}
\end{figure}

egin{figure}[htbp]
\centering
\includegraphics[width=\linewidth]{figures/section_ii__tuba__p40_05.jpg}
\caption{Illustration (page 40)}
\end{figure}

egin{figure}[htbp]
\centering
\includegraphics[width=\linewidth]{figures/section_ii__tuba__p40_06.jpg}
\caption{Illustration (page 40)}
\end{figure}

egin{figure}[htbp]
\centering
\includegraphics[width=\linewidth]{figures/section_ii__tuba__p40_07.jpg}
\caption{Illustration (page 40)}
\end{figure}

egin{figure}[htbp]
\centering
\includegraphics[width=\linewidth]{figures/section_ii__tuba__p40_08.png}
\caption{Illustration (page 40)}
\end{figure}

egin{figure}[htbp]
\centering
\includegraphics[width=\linewidth]{figures/section_ii__tuba__p40_09.jpg}
\caption{Illustration (page 40)}
\end{figure}

egin{figure}[htbp]
\centering
\includegraphics[width=\linewidth]{figures/section_ii__tuba__p40_10.jpg}
\caption{Illustration (page 40)}
\end{figure}

egin{figure}[htbp]
\centering
\includegraphics[width=\linewidth]{figures/section_ii__tuba__p40_11.png}
\caption{Illustration (page 40)}
\end{figure}

egin{figure}[htbp]
\centering
\includegraphics[width=\linewidth]{figures/section_ii__tuba__p40_12.png}
\caption{Illustration (page 40)}
\end{figure}

egin{figure}[htbp]
\centering
\includegraphics[width=\linewidth]{figures/section_ii__tuba__p40_13.png}
\caption{Illustration (page 40)}
\end{figure}

egin{figure}[htbp]
\centering
\includegraphics[width=\linewidth]{figures/section_ii__tuba__p40_14.png}
\caption{Illustration (page 40)}
\end{figure}

egin{figure}[htbp]
\centering
\includegraphics[width=\linewidth]{figures/section_ii__tuba__p40_15.png}
\caption{Illustration (page 40)}
\end{figure}

egin{figure}[htbp]
\centering
\includegraphics[width=\linewidth]{figures/section_ii__tuba__p40_16.png}
\caption{Illustration (page 40)}
\end{figure}

egin{figure}[htbp]
\centering
\includegraphics[width=\linewidth]{figures/section_ii__tuba__p40_17.png}
\caption{Illustration (page 40)}
\end{figure}

egin{figure}[htbp]
\centering
\includegraphics[width=\linewidth]{figures/section_ii__tuba__p40_18.png}
\caption{Illustration (page 40)}
\end{figure}

egin{figure}[htbp]
\centering
\includegraphics[width=\linewidth]{figures/section_ii__tuba__p40_19.png}
\caption{Illustration (page 40)}
\end{figure}

egin{figure}[htbp]
\centering
\includegraphics[width=\linewidth]{figures/section_ii__tuba__p40_20.jpg}
\caption{Illustration (page 40)}
\end{figure}

egin{figure}[htbp]
\centering
\includegraphics[width=\linewidth]{figures/section_ii__tuba__p40_21.jpg}
\caption{Illustration (page 40)}
\end{figure}

egin{figure}[htbp]
\centering
\includegraphics[width=\linewidth]{figures/section_ii__tuba__p40_22.jpg}
\caption{Illustration (page 40)}
\end{figure}

egin{figure}[htbp]
\centering
\includegraphics[width=\linewidth]{figures/section_ii__tuba__p40_23.jpg}
\caption{Illustration (page 40)}
\end{figure}

egin{figure}[htbp]
\centering
\includegraphics[width=\linewidth]{figures/section_ii__tuba__p40_24.jpg}
\caption{Illustration (page 40)}
\end{figure}

egin{figure}[htbp]
\centering
\includegraphics[width=\linewidth]{figures/section_ii__tuba__p40_25.jpg}
\caption{Illustration (page 40)}
\end{figure}

egin{figure}[htbp]
\centering
\includegraphics[width=\linewidth]{figures/section_ii__tuba__p40_26.jpg}
\caption{Illustration (page 40)}
\end{figure}

egin{figure}[htbp]
\centering
\includegraphics[width=\linewidth]{figures/section_ii__tuba__p40_27.png}
\caption{Illustration (page 40)}
\end{figure}

egin{figure}[htbp]
\centering
\includegraphics[width=\linewidth]{figures/section_ii__tuba__p40_28.png}
\caption{Illustration (page 40)}
\end{figure}

egin{figure}[htbp]
\centering
\includegraphics[width=\linewidth]{figures/section_ii__tuba__p40_29.png}
\caption{Illustration (page 40)}
\end{figure}

egin{figure}[htbp]
\centering
\includegraphics[width=\linewidth]{figures/section_ii__tuba__p40_30.png}
\caption{Illustration (page 40)}
\end{figure}

egin{figure}[htbp]
\centering
\includegraphics[width=\linewidth]{figures/section_ii__tuba__p40_31.png}
\caption{Illustration (page 40)}
\end{figure}

egin{figure}[htbp]
\centering
\includegraphics[width=\linewidth]{figures/section_ii__tuba__p40_32.jpg}
\caption{Illustration (page 40)}
\end{figure}

egin{figure}[htbp]
\centering
\includegraphics[width=\linewidth]{figures/section_ii__tuba__p40_33.png}
\caption{Illustration (page 40)}
\end{figure}
\subsection{Trompette en Do   [605 échantillons]}
Trompette en Do   [605 échantillons] → F1 spectral = 786 Hz (/a/) · Fp (Section I) = 1 324 Hz = F2 spectral TechniqueNF1F2F3F4F5F6Fiab. Ordinario9678613241766230429283714●○○ Crescendo81194026573149358441614695●○○ Staccato2984013891970248731123714●○○ Trille494878861253167621012322●●○ Flutter / Flatterzunge9079713141755233629283510●○○ Harmoniques nat.1182614201783225626543795●●○ Glissando2874314702156287835964337●●○ Slap183027641109145418842466●●○ Brassy / Cuivré31176629393768438249855599●○○ Autres9974512421709215227113693●○○ Fond vert = ordinario · Fond jaune = F1 dans cluster /o/ 420–550 Hz · Fond violet = technique étendue · Fiabilité: ●●● ≥0.55 / ●●○ 0.35–0.54 / ●○○ <0.35
\par\medskip

egin{figure}[htbp]
\centering
\includegraphics[width=\linewidth]{figures/section_ii__trompette__p41_01.jpg}
\caption{Illustration (page 41)}
\end{figure}

egin{figure}[htbp]
\centering
\includegraphics[width=\linewidth]{figures/section_ii__trompette__p41_02.jpg}
\caption{Illustration (page 41)}
\end{figure}

egin{figure}[htbp]
\centering
\includegraphics[width=\linewidth]{figures/section_ii__trompette__p41_03.png}
\caption{Illustration (page 41)}
\end{figure}

egin{figure}[htbp]
\centering
\includegraphics[width=\linewidth]{figures/section_ii__trompette__p41_04.jpg}
\caption{Illustration (page 41)}
\end{figure}

egin{figure}[htbp]
\centering
\includegraphics[width=\linewidth]{figures/section_ii__trompette__p41_05.jpg}
\caption{Illustration (page 41)}
\end{figure}

egin{figure}[htbp]
\centering
\includegraphics[width=\linewidth]{figures/section_ii__trompette__p41_06.jpg}
\caption{Illustration (page 41)}
\end{figure}

egin{figure}[htbp]
\centering
\includegraphics[width=\linewidth]{figures/section_ii__trompette__p41_07.jpg}
\caption{Illustration (page 41)}
\end{figure}

egin{figure}[htbp]
\centering
\includegraphics[width=\linewidth]{figures/section_ii__trompette__p41_08.png}
\caption{Illustration (page 41)}
\end{figure}

egin{figure}[htbp]
\centering
\includegraphics[width=\linewidth]{figures/section_ii__trompette__p41_09.jpg}
\caption{Illustration (page 41)}
\end{figure}

egin{figure}[htbp]
\centering
\includegraphics[width=\linewidth]{figures/section_ii__trompette__p41_10.jpg}
\caption{Illustration (page 41)}
\end{figure}

egin{figure}[htbp]
\centering
\includegraphics[width=\linewidth]{figures/section_ii__trompette__p41_11.png}
\caption{Illustration (page 41)}
\end{figure}

egin{figure}[htbp]
\centering
\includegraphics[width=\linewidth]{figures/section_ii__trompette__p41_12.png}
\caption{Illustration (page 41)}
\end{figure}

egin{figure}[htbp]
\centering
\includegraphics[width=\linewidth]{figures/section_ii__trompette__p41_13.png}
\caption{Illustration (page 41)}
\end{figure}

egin{figure}[htbp]
\centering
\includegraphics[width=\linewidth]{figures/section_ii__trompette__p41_14.png}
\caption{Illustration (page 41)}
\end{figure}

egin{figure}[htbp]
\centering
\includegraphics[width=\linewidth]{figures/section_ii__trompette__p41_15.png}
\caption{Illustration (page 41)}
\end{figure}

egin{figure}[htbp]
\centering
\includegraphics[width=\linewidth]{figures/section_ii__trompette__p41_16.png}
\caption{Illustration (page 41)}
\end{figure}

egin{figure}[htbp]
\centering
\includegraphics[width=\linewidth]{figures/section_ii__trompette__p41_17.png}
\caption{Illustration (page 41)}
\end{figure}

egin{figure}[htbp]
\centering
\includegraphics[width=\linewidth]{figures/section_ii__trompette__p41_18.png}
\caption{Illustration (page 41)}
\end{figure}

egin{figure}[htbp]
\centering
\includegraphics[width=\linewidth]{figures/section_ii__trompette__p41_19.png}
\caption{Illustration (page 41)}
\end{figure}

egin{figure}[htbp]
\centering
\includegraphics[width=\linewidth]{figures/section_ii__trompette__p41_20.jpg}
\caption{Illustration (page 41)}
\end{figure}

egin{figure}[htbp]
\centering
\includegraphics[width=\linewidth]{figures/section_ii__trompette__p41_21.jpg}
\caption{Illustration (page 41)}
\end{figure}

egin{figure}[htbp]
\centering
\includegraphics[width=\linewidth]{figures/section_ii__trompette__p41_22.jpg}
\caption{Illustration (page 41)}
\end{figure}

egin{figure}[htbp]
\centering
\includegraphics[width=\linewidth]{figures/section_ii__trompette__p41_23.jpg}
\caption{Illustration (page 41)}
\end{figure}

egin{figure}[htbp]
\centering
\includegraphics[width=\linewidth]{figures/section_ii__trompette__p41_24.jpg}
\caption{Illustration (page 41)}
\end{figure}

egin{figure}[htbp]
\centering
\includegraphics[width=\linewidth]{figures/section_ii__trompette__p41_25.jpg}
\caption{Illustration (page 41)}
\end{figure}

egin{figure}[htbp]
\centering
\includegraphics[width=\linewidth]{figures/section_ii__trompette__p41_26.jpg}
\caption{Illustration (page 41)}
\end{figure}

egin{figure}[htbp]
\centering
\includegraphics[width=\linewidth]{figures/section_ii__trompette__p41_27.png}
\caption{Illustration (page 41)}
\end{figure}

egin{figure}[htbp]
\centering
\includegraphics[width=\linewidth]{figures/section_ii__trompette__p41_28.png}
\caption{Illustration (page 41)}
\end{figure}

egin{figure}[htbp]
\centering
\includegraphics[width=\linewidth]{figures/section_ii__trompette__p41_29.png}
\caption{Illustration (page 41)}
\end{figure}

egin{figure}[htbp]
\centering
\includegraphics[width=\linewidth]{figures/section_ii__trompette__p41_30.png}
\caption{Illustration (page 41)}
\end{figure}

egin{figure}[htbp]
\centering
\includegraphics[width=\linewidth]{figures/section_ii__trompette__p41_31.png}
\caption{Illustration (page 41)}
\end{figure}

egin{figure}[htbp]
\centering
\includegraphics[width=\linewidth]{figures/section_ii__trompette__p41_32.jpg}
\caption{Illustration (page 41)}
\end{figure}

egin{figure}[htbp]
\centering
\includegraphics[width=\linewidth]{figures/section_ii__trompette__p41_33.png}
\caption{Illustration (page 41)}
\end{figure}
\subsection{Basson   [806 échantillons]}
Famille: Bois Basson   [806 échantillons] TechniqueNF1F2F3F4F5F6Fiab. Ordinario1264959151217165820032950●●○ Crescendo1154087411141155819062505●●○ Staccato414958611206165819702143●●○ Vibrato413776891195156119702606●●○ Trille765069111233157719242095●●○ Flutter / Flatterzunge324959151227162618952186●●○ Harmoniques nat.84097431184152920464554●●○ Glissando425289481217171220784317●●○ Multiphoniques564316891195151818522110●●○ Key click19194495926126016262250●●○ Autres1314908931246164219942579●●○ Fond vert = ordinario · Fond jaune = F1 dans cluster /o/ 420–550 Hz · Fond violet = technique étendue · Fiabilité: ●●● ≥0.55 / ●●○ 0.35–0.54 / ●○○ <0.35 Hautbois   [776 échantillons] → F1 spectral = 743 Hz (/a/) · Fp (Section I) = 1 460 Hz ≈ F2 spectral TechniqueNF1F2F3F4F5F6Fiab. Ordinario10774312601680230428103725●○○ Crescendo9974312851766236530073958●○○ Staccato3378612491636206730683768●○○ Vibrato3582913141744224034784393●○○ Trille8371613101853217426443098●○○ Flutter / Flatterzunge3679713461852262733054231●○○ Harmoniques nat.1388312701561212129933994●●○ Glissando3682915182078262734454264●○○ Multiphoniques77101214861981244431123941●○○ Key click973211951776215325523004●○○ Autres14679513041746225730414215●○○ Fond vert = ordinario · Fond jaune = F1 dans cluster /o/ 420–550 Hz · Fond violet = technique étendue · Fiabilité: ●●● ≥0.55 / ●●○ 0.35–0.54 / ●○○ <0.35 Clarinette Sib   [897 échantillons] → F1 spectral = 463 Hz (/o/) · Fp (Section I) = 1 016 Hz ≈ F2 spectral (harmoniques impairs) TechniqueNF1F2F3F4F5F6Fiab. Ordinario12646311301497185226813273●○○ Crescendo12753411651597219525823265●○○ Staccato4552813141766234727883510●○○ Trille784639431304162220572969●○○ Flutter / Flatterzunge13659313261920241229073295●○○ Glissando3773214642153284236073628●●● Multiphoniques353127751130159423903047●○○ Key click202696571023142118202961●●○ Autres13257112511845234330553471●○○
\par\medskip

egin{figure}[htbp]
\centering
\includegraphics[width=\linewidth]{figures/section_ii__basson__p42_01.jpg}
\caption{Illustration (page 42)}
\end{figure}

egin{figure}[htbp]
\centering
\includegraphics[width=\linewidth]{figures/section_ii__basson__p42_02.jpg}
\caption{Illustration (page 42)}
\end{figure}

egin{figure}[htbp]
\centering
\includegraphics[width=\linewidth]{figures/section_ii__basson__p42_03.png}
\caption{Illustration (page 42)}
\end{figure}

egin{figure}[htbp]
\centering
\includegraphics[width=\linewidth]{figures/section_ii__basson__p42_04.jpg}
\caption{Illustration (page 42)}
\end{figure}

egin{figure}[htbp]
\centering
\includegraphics[width=\linewidth]{figures/section_ii__basson__p42_05.jpg}
\caption{Illustration (page 42)}
\end{figure}

egin{figure}[htbp]
\centering
\includegraphics[width=\linewidth]{figures/section_ii__basson__p42_06.jpg}
\caption{Illustration (page 42)}
\end{figure}

egin{figure}[htbp]
\centering
\includegraphics[width=\linewidth]{figures/section_ii__basson__p42_07.jpg}
\caption{Illustration (page 42)}
\end{figure}

egin{figure}[htbp]
\centering
\includegraphics[width=\linewidth]{figures/section_ii__basson__p42_08.png}
\caption{Illustration (page 42)}
\end{figure}

egin{figure}[htbp]
\centering
\includegraphics[width=\linewidth]{figures/section_ii__basson__p42_09.jpg}
\caption{Illustration (page 42)}
\end{figure}

egin{figure}[htbp]
\centering
\includegraphics[width=\linewidth]{figures/section_ii__basson__p42_10.jpg}
\caption{Illustration (page 42)}
\end{figure}

egin{figure}[htbp]
\centering
\includegraphics[width=\linewidth]{figures/section_ii__basson__p42_11.png}
\caption{Illustration (page 42)}
\end{figure}

egin{figure}[htbp]
\centering
\includegraphics[width=\linewidth]{figures/section_ii__basson__p42_12.png}
\caption{Illustration (page 42)}
\end{figure}

egin{figure}[htbp]
\centering
\includegraphics[width=\linewidth]{figures/section_ii__basson__p42_13.png}
\caption{Illustration (page 42)}
\end{figure}

egin{figure}[htbp]
\centering
\includegraphics[width=\linewidth]{figures/section_ii__basson__p42_14.png}
\caption{Illustration (page 42)}
\end{figure}

egin{figure}[htbp]
\centering
\includegraphics[width=\linewidth]{figures/section_ii__basson__p42_15.png}
\caption{Illustration (page 42)}
\end{figure}

egin{figure}[htbp]
\centering
\includegraphics[width=\linewidth]{figures/section_ii__basson__p42_16.png}
\caption{Illustration (page 42)}
\end{figure}

egin{figure}[htbp]
\centering
\includegraphics[width=\linewidth]{figures/section_ii__basson__p42_17.png}
\caption{Illustration (page 42)}
\end{figure}

egin{figure}[htbp]
\centering
\includegraphics[width=\linewidth]{figures/section_ii__basson__p42_18.png}
\caption{Illustration (page 42)}
\end{figure}

egin{figure}[htbp]
\centering
\includegraphics[width=\linewidth]{figures/section_ii__basson__p42_19.png}
\caption{Illustration (page 42)}
\end{figure}

egin{figure}[htbp]
\centering
\includegraphics[width=\linewidth]{figures/section_ii__basson__p42_20.jpg}
\caption{Illustration (page 42)}
\end{figure}

egin{figure}[htbp]
\centering
\includegraphics[width=\linewidth]{figures/section_ii__basson__p42_21.jpg}
\caption{Illustration (page 42)}
\end{figure}

egin{figure}[htbp]
\centering
\includegraphics[width=\linewidth]{figures/section_ii__basson__p42_22.jpg}
\caption{Illustration (page 42)}
\end{figure}

egin{figure}[htbp]
\centering
\includegraphics[width=\linewidth]{figures/section_ii__basson__p42_23.jpg}
\caption{Illustration (page 42)}
\end{figure}

egin{figure}[htbp]
\centering
\includegraphics[width=\linewidth]{figures/section_ii__basson__p42_24.jpg}
\caption{Illustration (page 42)}
\end{figure}

egin{figure}[htbp]
\centering
\includegraphics[width=\linewidth]{figures/section_ii__basson__p42_25.jpg}
\caption{Illustration (page 42)}
\end{figure}

egin{figure}[htbp]
\centering
\includegraphics[width=\linewidth]{figures/section_ii__basson__p42_26.jpg}
\caption{Illustration (page 42)}
\end{figure}

egin{figure}[htbp]
\centering
\includegraphics[width=\linewidth]{figures/section_ii__basson__p42_27.png}
\caption{Illustration (page 42)}
\end{figure}

egin{figure}[htbp]
\centering
\includegraphics[width=\linewidth]{figures/section_ii__basson__p42_28.png}
\caption{Illustration (page 42)}
\end{figure}

egin{figure}[htbp]
\centering
\includegraphics[width=\linewidth]{figures/section_ii__basson__p42_29.png}
\caption{Illustration (page 42)}
\end{figure}

egin{figure}[htbp]
\centering
\includegraphics[width=\linewidth]{figures/section_ii__basson__p42_30.png}
\caption{Illustration (page 42)}
\end{figure}

egin{figure}[htbp]
\centering
\includegraphics[width=\linewidth]{figures/section_ii__basson__p42_31.png}
\caption{Illustration (page 42)}
\end{figure}

egin{figure}[htbp]
\centering
\includegraphics[width=\linewidth]{figures/section_ii__basson__p42_32.jpg}
\caption{Illustration (page 42)}
\end{figure}

egin{figure}[htbp]
\centering
\includegraphics[width=\linewidth]{figures/section_ii__basson__p42_33.png}
\caption{Illustration (page 42)}
\end{figure}
\subsection{Flûte traversière   [1135 échantillons]}
Fond vert = ordinario · Fond jaune = F1 dans cluster /o/ 420–550 Hz · Fond violet = technique étendue · Fiabilité: ●●● ≥0.55 / ●●○ 0.35–0.54 / ●○○ <0.35 Flûte traversière   [1135 échantillons] → F1 spectral = 743 Hz · Fp (Section I) = 1 354 Hz ≈ F2 spectral TechniqueNF1F2F3F4F5F6Fiab. Ordinario11874311631594198124443122●○○ Crescendo10367711441560200124663168●○○ Staccato3757111951615206726063510●○○ Trille7274914121807188420042092●○○ Flutter / Flatterzunge11657110661518191623152627●○○ Harmoniques nat.3886114101755227228324016●○○ Pizzicato203987111217160420562788●●○ Multiphoniques324318721270160419382476●●○ Son éolien133448081130149718842282●●○ Key click153987861120142118302186●●○ Tongue ram151945811217172330685566●●● Autres21865110611496191123683056●○○ Fond vert = ordinario · Fond jaune = F1 dans cluster /o/ 420–550 Hz · Fond violet = technique étendue · Fiabilité: ●●● ≥0.55 / ●●○ 0.35–0.54 / ●○○ <0.35
\par\medskip

egin{figure}[htbp]
\centering
\includegraphics[width=\linewidth]{figures/section_ii__basson__p43_01.jpg}
\caption{Illustration (page 43)}
\end{figure}

egin{figure}[htbp]
\centering
\includegraphics[width=\linewidth]{figures/section_ii__basson__p43_02.jpg}
\caption{Illustration (page 43)}
\end{figure}

egin{figure}[htbp]
\centering
\includegraphics[width=\linewidth]{figures/section_ii__basson__p43_03.png}
\caption{Illustration (page 43)}
\end{figure}

egin{figure}[htbp]
\centering
\includegraphics[width=\linewidth]{figures/section_ii__basson__p43_04.jpg}
\caption{Illustration (page 43)}
\end{figure}

egin{figure}[htbp]
\centering
\includegraphics[width=\linewidth]{figures/section_ii__basson__p43_05.jpg}
\caption{Illustration (page 43)}
\end{figure}

egin{figure}[htbp]
\centering
\includegraphics[width=\linewidth]{figures/section_ii__basson__p43_06.jpg}
\caption{Illustration (page 43)}
\end{figure}

egin{figure}[htbp]
\centering
\includegraphics[width=\linewidth]{figures/section_ii__basson__p43_07.jpg}
\caption{Illustration (page 43)}
\end{figure}

egin{figure}[htbp]
\centering
\includegraphics[width=\linewidth]{figures/section_ii__basson__p43_08.png}
\caption{Illustration (page 43)}
\end{figure}

egin{figure}[htbp]
\centering
\includegraphics[width=\linewidth]{figures/section_ii__basson__p43_09.jpg}
\caption{Illustration (page 43)}
\end{figure}

egin{figure}[htbp]
\centering
\includegraphics[width=\linewidth]{figures/section_ii__basson__p43_10.jpg}
\caption{Illustration (page 43)}
\end{figure}

egin{figure}[htbp]
\centering
\includegraphics[width=\linewidth]{figures/section_ii__basson__p43_11.png}
\caption{Illustration (page 43)}
\end{figure}

egin{figure}[htbp]
\centering
\includegraphics[width=\linewidth]{figures/section_ii__basson__p43_12.png}
\caption{Illustration (page 43)}
\end{figure}

egin{figure}[htbp]
\centering
\includegraphics[width=\linewidth]{figures/section_ii__basson__p43_13.png}
\caption{Illustration (page 43)}
\end{figure}

egin{figure}[htbp]
\centering
\includegraphics[width=\linewidth]{figures/section_ii__basson__p43_14.png}
\caption{Illustration (page 43)}
\end{figure}

egin{figure}[htbp]
\centering
\includegraphics[width=\linewidth]{figures/section_ii__basson__p43_15.png}
\caption{Illustration (page 43)}
\end{figure}

egin{figure}[htbp]
\centering
\includegraphics[width=\linewidth]{figures/section_ii__basson__p43_16.png}
\caption{Illustration (page 43)}
\end{figure}

egin{figure}[htbp]
\centering
\includegraphics[width=\linewidth]{figures/section_ii__basson__p43_17.png}
\caption{Illustration (page 43)}
\end{figure}

egin{figure}[htbp]
\centering
\includegraphics[width=\linewidth]{figures/section_ii__basson__p43_18.png}
\caption{Illustration (page 43)}
\end{figure}

egin{figure}[htbp]
\centering
\includegraphics[width=\linewidth]{figures/section_ii__basson__p43_19.png}
\caption{Illustration (page 43)}
\end{figure}

egin{figure}[htbp]
\centering
\includegraphics[width=\linewidth]{figures/section_ii__basson__p43_20.jpg}
\caption{Illustration (page 43)}
\end{figure}

egin{figure}[htbp]
\centering
\includegraphics[width=\linewidth]{figures/section_ii__basson__p43_21.jpg}
\caption{Illustration (page 43)}
\end{figure}

egin{figure}[htbp]
\centering
\includegraphics[width=\linewidth]{figures/section_ii__basson__p43_22.jpg}
\caption{Illustration (page 43)}
\end{figure}

egin{figure}[htbp]
\centering
\includegraphics[width=\linewidth]{figures/section_ii__basson__p43_23.jpg}
\caption{Illustration (page 43)}
\end{figure}

egin{figure}[htbp]
\centering
\includegraphics[width=\linewidth]{figures/section_ii__basson__p43_24.jpg}
\caption{Illustration (page 43)}
\end{figure}

egin{figure}[htbp]
\centering
\includegraphics[width=\linewidth]{figures/section_ii__basson__p43_25.jpg}
\caption{Illustration (page 43)}
\end{figure}

egin{figure}[htbp]
\centering
\includegraphics[width=\linewidth]{figures/section_ii__basson__p43_26.jpg}
\caption{Illustration (page 43)}
\end{figure}

egin{figure}[htbp]
\centering
\includegraphics[width=\linewidth]{figures/section_ii__basson__p43_27.png}
\caption{Illustration (page 43)}
\end{figure}

egin{figure}[htbp]
\centering
\includegraphics[width=\linewidth]{figures/section_ii__basson__p43_28.png}
\caption{Illustration (page 43)}
\end{figure}

egin{figure}[htbp]
\centering
\includegraphics[width=\linewidth]{figures/section_ii__basson__p43_29.png}
\caption{Illustration (page 43)}
\end{figure}

egin{figure}[htbp]
\centering
\includegraphics[width=\linewidth]{figures/section_ii__basson__p43_30.png}
\caption{Illustration (page 43)}
\end{figure}

egin{figure}[htbp]
\centering
\includegraphics[width=\linewidth]{figures/section_ii__basson__p43_31.png}
\caption{Illustration (page 43)}
\end{figure}

egin{figure}[htbp]
\centering
\includegraphics[width=\linewidth]{figures/section_ii__basson__p43_32.jpg}
\caption{Illustration (page 43)}
\end{figure}

egin{figure}[htbp]
\centering
\includegraphics[width=\linewidth]{figures/section_ii__basson__p43_33.png}
\caption{Illustration (page 43)}
\end{figure}
\subsection{Saxophone alto   [859 échantillons]}
Famille: Saxophones Saxophone alto   [859 échantillons] → F1 spectral = 398 Hz (/o/) · Fp (Section I) = 829 Hz = F2 spectral (tuyau conique) TechniqueNF1F2F3F4F5F6Fiab. Ordinario993988291303179823692756●○○ Crescendo933668291270171221142631●○○ Staccato3162412811680205623582928●○○ Trille654258491297188423632784●○○ Flutter / Flatterzunge333988291217188423903133●○○ Harmoniques nat.1047512181999226926433225●○○ Glissando323486981140171220462530●●● Multiphoniques474097541324182021532509●●○ Slap523448641243167022212732●●○ Bisbigliando294317971249157221752476●○○ Son éolien2690413141669193822502746●○○ Key click12258560915129216802282●●○ Double articulation333557111174164722932584●○○ Autres10441210031497194524682823●○○ Fond vert = ordinario · Fond jaune = F1 dans cluster /o/ 420–550 Hz · Fond violet = technique étendue · Fiabilité: ●●● ≥0.55 / ●●○ 0.35–0.54 / ●○○ <0.35
\par\medskip

egin{figure}[htbp]
\centering
\includegraphics[width=\linewidth]{figures/section_ii__saxophone_alto__p44_01.jpg}
\caption{Illustration (page 44)}
\end{figure}

egin{figure}[htbp]
\centering
\includegraphics[width=\linewidth]{figures/section_ii__saxophone_alto__p44_02.jpg}
\caption{Illustration (page 44)}
\end{figure}

egin{figure}[htbp]
\centering
\includegraphics[width=\linewidth]{figures/section_ii__saxophone_alto__p44_03.png}
\caption{Illustration (page 44)}
\end{figure}

egin{figure}[htbp]
\centering
\includegraphics[width=\linewidth]{figures/section_ii__saxophone_alto__p44_04.jpg}
\caption{Illustration (page 44)}
\end{figure}

egin{figure}[htbp]
\centering
\includegraphics[width=\linewidth]{figures/section_ii__saxophone_alto__p44_05.jpg}
\caption{Illustration (page 44)}
\end{figure}

egin{figure}[htbp]
\centering
\includegraphics[width=\linewidth]{figures/section_ii__saxophone_alto__p44_06.jpg}
\caption{Illustration (page 44)}
\end{figure}

egin{figure}[htbp]
\centering
\includegraphics[width=\linewidth]{figures/section_ii__saxophone_alto__p44_07.jpg}
\caption{Illustration (page 44)}
\end{figure}

egin{figure}[htbp]
\centering
\includegraphics[width=\linewidth]{figures/section_ii__saxophone_alto__p44_08.png}
\caption{Illustration (page 44)}
\end{figure}

egin{figure}[htbp]
\centering
\includegraphics[width=\linewidth]{figures/section_ii__saxophone_alto__p44_09.jpg}
\caption{Illustration (page 44)}
\end{figure}

egin{figure}[htbp]
\centering
\includegraphics[width=\linewidth]{figures/section_ii__saxophone_alto__p44_10.jpg}
\caption{Illustration (page 44)}
\end{figure}

egin{figure}[htbp]
\centering
\includegraphics[width=\linewidth]{figures/section_ii__saxophone_alto__p44_11.png}
\caption{Illustration (page 44)}
\end{figure}

egin{figure}[htbp]
\centering
\includegraphics[width=\linewidth]{figures/section_ii__saxophone_alto__p44_12.png}
\caption{Illustration (page 44)}
\end{figure}

egin{figure}[htbp]
\centering
\includegraphics[width=\linewidth]{figures/section_ii__saxophone_alto__p44_13.png}
\caption{Illustration (page 44)}
\end{figure}

egin{figure}[htbp]
\centering
\includegraphics[width=\linewidth]{figures/section_ii__saxophone_alto__p44_14.png}
\caption{Illustration (page 44)}
\end{figure}

egin{figure}[htbp]
\centering
\includegraphics[width=\linewidth]{figures/section_ii__saxophone_alto__p44_15.png}
\caption{Illustration (page 44)}
\end{figure}

egin{figure}[htbp]
\centering
\includegraphics[width=\linewidth]{figures/section_ii__saxophone_alto__p44_16.png}
\caption{Illustration (page 44)}
\end{figure}

egin{figure}[htbp]
\centering
\includegraphics[width=\linewidth]{figures/section_ii__saxophone_alto__p44_17.png}
\caption{Illustration (page 44)}
\end{figure}

egin{figure}[htbp]
\centering
\includegraphics[width=\linewidth]{figures/section_ii__saxophone_alto__p44_18.png}
\caption{Illustration (page 44)}
\end{figure}

egin{figure}[htbp]
\centering
\includegraphics[width=\linewidth]{figures/section_ii__saxophone_alto__p44_19.png}
\caption{Illustration (page 44)}
\end{figure}

egin{figure}[htbp]
\centering
\includegraphics[width=\linewidth]{figures/section_ii__saxophone_alto__p44_20.jpg}
\caption{Illustration (page 44)}
\end{figure}

egin{figure}[htbp]
\centering
\includegraphics[width=\linewidth]{figures/section_ii__saxophone_alto__p44_21.jpg}
\caption{Illustration (page 44)}
\end{figure}

egin{figure}[htbp]
\centering
\includegraphics[width=\linewidth]{figures/section_ii__saxophone_alto__p44_22.jpg}
\caption{Illustration (page 44)}
\end{figure}

egin{figure}[htbp]
\centering
\includegraphics[width=\linewidth]{figures/section_ii__saxophone_alto__p44_23.jpg}
\caption{Illustration (page 44)}
\end{figure}

egin{figure}[htbp]
\centering
\includegraphics[width=\linewidth]{figures/section_ii__saxophone_alto__p44_24.jpg}
\caption{Illustration (page 44)}
\end{figure}

egin{figure}[htbp]
\centering
\includegraphics[width=\linewidth]{figures/section_ii__saxophone_alto__p44_25.jpg}
\caption{Illustration (page 44)}
\end{figure}

egin{figure}[htbp]
\centering
\includegraphics[width=\linewidth]{figures/section_ii__saxophone_alto__p44_26.jpg}
\caption{Illustration (page 44)}
\end{figure}

egin{figure}[htbp]
\centering
\includegraphics[width=\linewidth]{figures/section_ii__saxophone_alto__p44_27.png}
\caption{Illustration (page 44)}
\end{figure}

egin{figure}[htbp]
\centering
\includegraphics[width=\linewidth]{figures/section_ii__saxophone_alto__p44_28.png}
\caption{Illustration (page 44)}
\end{figure}

egin{figure}[htbp]
\centering
\includegraphics[width=\linewidth]{figures/section_ii__saxophone_alto__p44_29.png}
\caption{Illustration (page 44)}
\end{figure}

egin{figure}[htbp]
\centering
\includegraphics[width=\linewidth]{figures/section_ii__saxophone_alto__p44_30.png}
\caption{Illustration (page 44)}
\end{figure}

egin{figure}[htbp]
\centering
\includegraphics[width=\linewidth]{figures/section_ii__saxophone_alto__p44_31.png}
\caption{Illustration (page 44)}
\end{figure}

egin{figure}[htbp]
\centering
\includegraphics[width=\linewidth]{figures/section_ii__saxophone_alto__p44_32.jpg}
\caption{Illustration (page 44)}
\end{figure}

egin{figure}[htbp]
\centering
\includegraphics[width=\linewidth]{figures/section_ii__saxophone_alto__p44_33.png}
\caption{Illustration (page 44)}
\end{figure}
\subsection{Violon   [2700 échantillons]}
Famille: Cordes Violon   [2700 échantillons] → F1 spectral = 506 Hz (/o/, caisse) · Fp (Section I) = 1 034 Hz = F2 (bridge hill) TechniqueNF1F2F3F4F5F6Fiab. Ordinario28450615182347328439084705●○○ Crescendo14449912792303307639254717●○○ Staccato6047410982207312239194565●○○ Vibrato2494099902229307938984436●○○ Trille12450611962164256936334285●○○ Trémolo28549511632100292837474404●○○ Harmoniques nat.11960911301948254934924254●○○ Sul ponticello20544911172225287837984404●○○ Sul tasto5051710122164261635104242●○○ Pizzicato4854038171255181725323586●●○ Col legno1243888511439212127804016●●○ Derrière chevalet753811522240306833815050●○○ Autres16947312022160284338344531●○○ Fond vert = ordinario · Fond jaune = F1 dans cluster /o/ 420–550 Hz · Fond violet = technique étendue · Fiabilité: ●●● ≥0.55 / ●●○ 0.35–0.54 / ●○○ <0.35 Alto   [3021 échantillons] → F1 spectral = 369 Hz (/o/) = Fp CORRIGÉ (Section I) — ancienne valeur 1 516 Hz était le bridge hill TechniqueNF1F2F3F4F5F6Fiab. Ordinario3093778291540220729393478●○○ Crescendo1603808141611223928923571●●○ Staccato643668611626233628213413●●○ Vibrato2973667541400211027883445●●○ Trille1284148451491213227943187●○○ Trémolo3093667751357201326493273●●○ Harmoniques nat.1313788751615219729053447●●○ Sul ponticello2913707961542221528833403●○○ Sul tasto1903667541276202425573149●●○ Pizzicato4052976611086153421702964●●○ Col legno1283667431373191623263117●●○ Derrière chevalet435510011960299343284942●○○ Autres1653829111566226229373660●●○ Fond vert = ordinario · Fond jaune = F1 dans cluster /o/ 420–550 Hz · Fond violet = technique étendue · Fiabilité: ●●● ≥0.55 / ●●○ 0.35–0.54 / ●○○ <0.35 Violoncelle   [2287 échantillons] → F1 spectral = 205 Hz (/u/) · Fp (Section I) = 499 Hz = F2 spectral TechniqueNF1F2F3F4F5F6Fiab. Ordinario2912055061130184120892509●○○ Crescendo1022285941124186122522630●●○ Staccato392055061044150719812498●●○ Vibrato2791944951087152920782476●●●
\par\medskip

egin{figure}[htbp]
\centering
\includegraphics[width=\linewidth]{figures/section_ii__violon__p45_01.jpg}
\caption{Illustration (page 45)}
\end{figure}

egin{figure}[htbp]
\centering
\includegraphics[width=\linewidth]{figures/section_ii__violon__p45_02.jpg}
\caption{Illustration (page 45)}
\end{figure}

egin{figure}[htbp]
\centering
\includegraphics[width=\linewidth]{figures/section_ii__violon__p45_03.png}
\caption{Illustration (page 45)}
\end{figure}

egin{figure}[htbp]
\centering
\includegraphics[width=\linewidth]{figures/section_ii__violon__p45_04.jpg}
\caption{Illustration (page 45)}
\end{figure}

egin{figure}[htbp]
\centering
\includegraphics[width=\linewidth]{figures/section_ii__violon__p45_05.jpg}
\caption{Illustration (page 45)}
\end{figure}

egin{figure}[htbp]
\centering
\includegraphics[width=\linewidth]{figures/section_ii__violon__p45_06.jpg}
\caption{Illustration (page 45)}
\end{figure}

egin{figure}[htbp]
\centering
\includegraphics[width=\linewidth]{figures/section_ii__violon__p45_07.jpg}
\caption{Illustration (page 45)}
\end{figure}

egin{figure}[htbp]
\centering
\includegraphics[width=\linewidth]{figures/section_ii__violon__p45_08.png}
\caption{Illustration (page 45)}
\end{figure}

egin{figure}[htbp]
\centering
\includegraphics[width=\linewidth]{figures/section_ii__violon__p45_09.jpg}
\caption{Illustration (page 45)}
\end{figure}

egin{figure}[htbp]
\centering
\includegraphics[width=\linewidth]{figures/section_ii__violon__p45_10.jpg}
\caption{Illustration (page 45)}
\end{figure}

egin{figure}[htbp]
\centering
\includegraphics[width=\linewidth]{figures/section_ii__violon__p45_11.png}
\caption{Illustration (page 45)}
\end{figure}

egin{figure}[htbp]
\centering
\includegraphics[width=\linewidth]{figures/section_ii__violon__p45_12.png}
\caption{Illustration (page 45)}
\end{figure}

egin{figure}[htbp]
\centering
\includegraphics[width=\linewidth]{figures/section_ii__violon__p45_13.png}
\caption{Illustration (page 45)}
\end{figure}

egin{figure}[htbp]
\centering
\includegraphics[width=\linewidth]{figures/section_ii__violon__p45_14.png}
\caption{Illustration (page 45)}
\end{figure}

egin{figure}[htbp]
\centering
\includegraphics[width=\linewidth]{figures/section_ii__violon__p45_15.png}
\caption{Illustration (page 45)}
\end{figure}

egin{figure}[htbp]
\centering
\includegraphics[width=\linewidth]{figures/section_ii__violon__p45_16.png}
\caption{Illustration (page 45)}
\end{figure}

egin{figure}[htbp]
\centering
\includegraphics[width=\linewidth]{figures/section_ii__violon__p45_17.png}
\caption{Illustration (page 45)}
\end{figure}

egin{figure}[htbp]
\centering
\includegraphics[width=\linewidth]{figures/section_ii__violon__p45_18.png}
\caption{Illustration (page 45)}
\end{figure}

egin{figure}[htbp]
\centering
\includegraphics[width=\linewidth]{figures/section_ii__violon__p45_19.png}
\caption{Illustration (page 45)}
\end{figure}

egin{figure}[htbp]
\centering
\includegraphics[width=\linewidth]{figures/section_ii__violon__p45_20.jpg}
\caption{Illustration (page 45)}
\end{figure}

egin{figure}[htbp]
\centering
\includegraphics[width=\linewidth]{figures/section_ii__violon__p45_21.jpg}
\caption{Illustration (page 45)}
\end{figure}

egin{figure}[htbp]
\centering
\includegraphics[width=\linewidth]{figures/section_ii__violon__p45_22.jpg}
\caption{Illustration (page 45)}
\end{figure}

egin{figure}[htbp]
\centering
\includegraphics[width=\linewidth]{figures/section_ii__violon__p45_23.jpg}
\caption{Illustration (page 45)}
\end{figure}

egin{figure}[htbp]
\centering
\includegraphics[width=\linewidth]{figures/section_ii__violon__p45_24.jpg}
\caption{Illustration (page 45)}
\end{figure}

egin{figure}[htbp]
\centering
\includegraphics[width=\linewidth]{figures/section_ii__violon__p45_25.jpg}
\caption{Illustration (page 45)}
\end{figure}

egin{figure}[htbp]
\centering
\includegraphics[width=\linewidth]{figures/section_ii__violon__p45_26.jpg}
\caption{Illustration (page 45)}
\end{figure}

egin{figure}[htbp]
\centering
\includegraphics[width=\linewidth]{figures/section_ii__violon__p45_27.png}
\caption{Illustration (page 45)}
\end{figure}

egin{figure}[htbp]
\centering
\includegraphics[width=\linewidth]{figures/section_ii__violon__p45_28.png}
\caption{Illustration (page 45)}
\end{figure}

egin{figure}[htbp]
\centering
\includegraphics[width=\linewidth]{figures/section_ii__violon__p45_29.png}
\caption{Illustration (page 45)}
\end{figure}

egin{figure}[htbp]
\centering
\includegraphics[width=\linewidth]{figures/section_ii__violon__p45_30.png}
\caption{Illustration (page 45)}
\end{figure}

egin{figure}[htbp]
\centering
\includegraphics[width=\linewidth]{figures/section_ii__violon__p45_31.png}
\caption{Illustration (page 45)}
\end{figure}

egin{figure}[htbp]
\centering
\includegraphics[width=\linewidth]{figures/section_ii__violon__p45_32.jpg}
\caption{Illustration (page 45)}
\end{figure}

egin{figure}[htbp]
\centering
\includegraphics[width=\linewidth]{figures/section_ii__violon__p45_33.png}
\caption{Illustration (page 45)}
\end{figure}
Trille902105921077150719602336●●○ Trémolo2192154841066148618842164●●○ Harmoniques nat.1142335921119169920452591●●● Sul ponticello2212015241132178421202495●●○ Sul tasto1461995221017134618952282●●○ Pizzicato309191488975128618342375●●● Col legno901886571028142719062379●●● Derrière chevalet4269581937126017982412●●○ Autres1082235901123172522672746●●○ Fond vert = ordinario · Fond jaune = F1 dans cluster /o/ 420–550 Hz · Fond violet = technique étendue · Fiabilité: ●●● ≥0.55 / ●●○ 0.35–0.54 / ●○○ <0.35 Contrebasse   [2576 échantillons] TechniqueNF1F2F3F4F5F6Fiab. Ordinario309172474808116314642347●●● Crescendo124179498962128415712365●●● Staccato52172463872121715401841●●● Vibrato297172474829116314752240●●● Trille108172474818114114592073●●● Trémolo228183474808114114641873●●● Harmoniques nat.119177521935132717342620●●● Sul ponticello240178481844116214802045●●● Sul tasto148172468802114714641977●●● Pizzicato384158486881123815592104●●● Col legno110151479915126516372094●●○ Derrière chevalet476411842347330540595792●○○ Autres125183505936124415201886●●● Fond vert = ordinario · Fond jaune = F1 dans cluster /o/ 420–550 Hz · Fond violet = technique étendue · Fiabilité: ●●● ≥0.55 / ●●○ 0.35–0.54 / ●○○ <0.35
\par\medskip

egin{figure}[htbp]
\centering
\includegraphics[width=\linewidth]{figures/section_ii__contrebasse__p46_01.jpg}
\caption{Illustration (page 46)}
\end{figure}

egin{figure}[htbp]
\centering
\includegraphics[width=\linewidth]{figures/section_ii__contrebasse__p46_02.jpg}
\caption{Illustration (page 46)}
\end{figure}

egin{figure}[htbp]
\centering
\includegraphics[width=\linewidth]{figures/section_ii__contrebasse__p46_03.png}
\caption{Illustration (page 46)}
\end{figure}

egin{figure}[htbp]
\centering
\includegraphics[width=\linewidth]{figures/section_ii__contrebasse__p46_04.jpg}
\caption{Illustration (page 46)}
\end{figure}

egin{figure}[htbp]
\centering
\includegraphics[width=\linewidth]{figures/section_ii__contrebasse__p46_05.jpg}
\caption{Illustration (page 46)}
\end{figure}

egin{figure}[htbp]
\centering
\includegraphics[width=\linewidth]{figures/section_ii__contrebasse__p46_06.jpg}
\caption{Illustration (page 46)}
\end{figure}

egin{figure}[htbp]
\centering
\includegraphics[width=\linewidth]{figures/section_ii__contrebasse__p46_07.jpg}
\caption{Illustration (page 46)}
\end{figure}

egin{figure}[htbp]
\centering
\includegraphics[width=\linewidth]{figures/section_ii__contrebasse__p46_08.png}
\caption{Illustration (page 46)}
\end{figure}

egin{figure}[htbp]
\centering
\includegraphics[width=\linewidth]{figures/section_ii__contrebasse__p46_09.jpg}
\caption{Illustration (page 46)}
\end{figure}

egin{figure}[htbp]
\centering
\includegraphics[width=\linewidth]{figures/section_ii__contrebasse__p46_10.jpg}
\caption{Illustration (page 46)}
\end{figure}

egin{figure}[htbp]
\centering
\includegraphics[width=\linewidth]{figures/section_ii__contrebasse__p46_11.png}
\caption{Illustration (page 46)}
\end{figure}

egin{figure}[htbp]
\centering
\includegraphics[width=\linewidth]{figures/section_ii__contrebasse__p46_12.png}
\caption{Illustration (page 46)}
\end{figure}

egin{figure}[htbp]
\centering
\includegraphics[width=\linewidth]{figures/section_ii__contrebasse__p46_13.png}
\caption{Illustration (page 46)}
\end{figure}

egin{figure}[htbp]
\centering
\includegraphics[width=\linewidth]{figures/section_ii__contrebasse__p46_14.png}
\caption{Illustration (page 46)}
\end{figure}

egin{figure}[htbp]
\centering
\includegraphics[width=\linewidth]{figures/section_ii__contrebasse__p46_15.png}
\caption{Illustration (page 46)}
\end{figure}

egin{figure}[htbp]
\centering
\includegraphics[width=\linewidth]{figures/section_ii__contrebasse__p46_16.png}
\caption{Illustration (page 46)}
\end{figure}

egin{figure}[htbp]
\centering
\includegraphics[width=\linewidth]{figures/section_ii__contrebasse__p46_17.png}
\caption{Illustration (page 46)}
\end{figure}

egin{figure}[htbp]
\centering
\includegraphics[width=\linewidth]{figures/section_ii__contrebasse__p46_18.png}
\caption{Illustration (page 46)}
\end{figure}

egin{figure}[htbp]
\centering
\includegraphics[width=\linewidth]{figures/section_ii__contrebasse__p46_19.png}
\caption{Illustration (page 46)}
\end{figure}

egin{figure}[htbp]
\centering
\includegraphics[width=\linewidth]{figures/section_ii__contrebasse__p46_20.jpg}
\caption{Illustration (page 46)}
\end{figure}

egin{figure}[htbp]
\centering
\includegraphics[width=\linewidth]{figures/section_ii__contrebasse__p46_21.jpg}
\caption{Illustration (page 46)}
\end{figure}

egin{figure}[htbp]
\centering
\includegraphics[width=\linewidth]{figures/section_ii__contrebasse__p46_22.jpg}
\caption{Illustration (page 46)}
\end{figure}

egin{figure}[htbp]
\centering
\includegraphics[width=\linewidth]{figures/section_ii__contrebasse__p46_23.jpg}
\caption{Illustration (page 46)}
\end{figure}

egin{figure}[htbp]
\centering
\includegraphics[width=\linewidth]{figures/section_ii__contrebasse__p46_24.jpg}
\caption{Illustration (page 46)}
\end{figure}

egin{figure}[htbp]
\centering
\includegraphics[width=\linewidth]{figures/section_ii__contrebasse__p46_25.jpg}
\caption{Illustration (page 46)}
\end{figure}

egin{figure}[htbp]
\centering
\includegraphics[width=\linewidth]{figures/section_ii__contrebasse__p46_26.jpg}
\caption{Illustration (page 46)}
\end{figure}

egin{figure}[htbp]
\centering
\includegraphics[width=\linewidth]{figures/section_ii__contrebasse__p46_27.png}
\caption{Illustration (page 46)}
\end{figure}

egin{figure}[htbp]
\centering
\includegraphics[width=\linewidth]{figures/section_ii__contrebasse__p46_28.png}
\caption{Illustration (page 46)}
\end{figure}

egin{figure}[htbp]
\centering
\includegraphics[width=\linewidth]{figures/section_ii__contrebasse__p46_29.png}
\caption{Illustration (page 46)}
\end{figure}

egin{figure}[htbp]
\centering
\includegraphics[width=\linewidth]{figures/section_ii__contrebasse__p46_30.png}
\caption{Illustration (page 46)}
\end{figure}

egin{figure}[htbp]
\centering
\includegraphics[width=\linewidth]{figures/section_ii__contrebasse__p46_31.png}
\caption{Illustration (page 46)}
\end{figure}

egin{figure}[htbp]
\centering
\includegraphics[width=\linewidth]{figures/section_ii__contrebasse__p46_32.jpg}
\caption{Illustration (page 46)}
\end{figure}

egin{figure}[htbp]
\centering
\includegraphics[width=\linewidth]{figures/section_ii__contrebasse__p46_33.png}
\caption{Illustration (page 46)}
\end{figure}
\subsection{Guitare   [559 échantillons]}
Famille: Cordes pincées et claviers Guitare   [559 échantillons] TechniqueNF1F2F3F4F5F6Fiab. Ordinario210194603980146422182939●●● Harmoniques nat.27194668969133520132035●●● Glissando151946031001141019813574●●○ Sul ponticello44205635969134616802756●●● Sul tasto37194592958131419812982●●● Pizzicato811946351007132521363224●●● Autres115205625980151320192371●●● Fond vert = ordinario · Fond jaune = F1 dans cluster /o/ 420–550 Hz · Fond violet = technique étendue · Fiabilité: ●●● ≥0.55 / ●●○ 0.35–0.54 / ●○○ <0.35 Harpe   [821 échantillons] TechniqueNF1F2F3F4F5F6Fiab. Ordinario2412696241034145423694156●●○ Trémolo1248614937189539944360●○○ Harmoniques nat.92282654998144320802967●●● Glissando652716981181159721112609●●○ Pizzicato79248581894132416582164●●○ Bisbigliando702856431171187626674144●●○ Buzz15215506872148617662401●●● Autres2582897081121163022003574●●○ Fond vert = ordinario · Fond jaune = F1 dans cluster /o/ 420–550 Hz · Fond violet = technique étendue · Fiabilité: ●●● ≥0.55 / ●●○ 0.35–0.54 / ●○○ <0.35 Accordéon   [1190 échantillons] TechniqueNF1F2F3F4F5F6Fiab. Ordinario6894209261486208926703962●○○ Crescendo1673758361472199524243582●○○ Staccato753346681281188423263854●○○ Son éolien23989151206158321862530●●○ Key click13777971303161522183348●○○ Autres2564198931376190926353611●○○ Fond vert = ordinario · Fond jaune = F1 dans cluster /o/ 420–550 Hz · Fond violet = technique étendue · Fiabilité: ●●● ≥0.55 / ●●○ 0.35–0.54 / ●○○ <0.35
\par\medskip

egin{figure}[htbp]
\centering
\includegraphics[width=\linewidth]{figures/section_ii__guitare__p47_01.jpg}
\caption{Illustration (page 47)}
\end{figure}

egin{figure}[htbp]
\centering
\includegraphics[width=\linewidth]{figures/section_ii__guitare__p47_02.jpg}
\caption{Illustration (page 47)}
\end{figure}

egin{figure}[htbp]
\centering
\includegraphics[width=\linewidth]{figures/section_ii__guitare__p47_03.png}
\caption{Illustration (page 47)}
\end{figure}

egin{figure}[htbp]
\centering
\includegraphics[width=\linewidth]{figures/section_ii__guitare__p47_04.jpg}
\caption{Illustration (page 47)}
\end{figure}

egin{figure}[htbp]
\centering
\includegraphics[width=\linewidth]{figures/section_ii__guitare__p47_05.jpg}
\caption{Illustration (page 47)}
\end{figure}

egin{figure}[htbp]
\centering
\includegraphics[width=\linewidth]{figures/section_ii__guitare__p47_06.jpg}
\caption{Illustration (page 47)}
\end{figure}

egin{figure}[htbp]
\centering
\includegraphics[width=\linewidth]{figures/section_ii__guitare__p47_07.jpg}
\caption{Illustration (page 47)}
\end{figure}

egin{figure}[htbp]
\centering
\includegraphics[width=\linewidth]{figures/section_ii__guitare__p47_08.png}
\caption{Illustration (page 47)}
\end{figure}

egin{figure}[htbp]
\centering
\includegraphics[width=\linewidth]{figures/section_ii__guitare__p47_09.jpg}
\caption{Illustration (page 47)}
\end{figure}

egin{figure}[htbp]
\centering
\includegraphics[width=\linewidth]{figures/section_ii__guitare__p47_10.jpg}
\caption{Illustration (page 47)}
\end{figure}

egin{figure}[htbp]
\centering
\includegraphics[width=\linewidth]{figures/section_ii__guitare__p47_11.png}
\caption{Illustration (page 47)}
\end{figure}

egin{figure}[htbp]
\centering
\includegraphics[width=\linewidth]{figures/section_ii__guitare__p47_12.png}
\caption{Illustration (page 47)}
\end{figure}

egin{figure}[htbp]
\centering
\includegraphics[width=\linewidth]{figures/section_ii__guitare__p47_13.png}
\caption{Illustration (page 47)}
\end{figure}

egin{figure}[htbp]
\centering
\includegraphics[width=\linewidth]{figures/section_ii__guitare__p47_14.png}
\caption{Illustration (page 47)}
\end{figure}

egin{figure}[htbp]
\centering
\includegraphics[width=\linewidth]{figures/section_ii__guitare__p47_15.png}
\caption{Illustration (page 47)}
\end{figure}

egin{figure}[htbp]
\centering
\includegraphics[width=\linewidth]{figures/section_ii__guitare__p47_16.png}
\caption{Illustration (page 47)}
\end{figure}

egin{figure}[htbp]
\centering
\includegraphics[width=\linewidth]{figures/section_ii__guitare__p47_17.png}
\caption{Illustration (page 47)}
\end{figure}

egin{figure}[htbp]
\centering
\includegraphics[width=\linewidth]{figures/section_ii__guitare__p47_18.png}
\caption{Illustration (page 47)}
\end{figure}

egin{figure}[htbp]
\centering
\includegraphics[width=\linewidth]{figures/section_ii__guitare__p47_19.png}
\caption{Illustration (page 47)}
\end{figure}

egin{figure}[htbp]
\centering
\includegraphics[width=\linewidth]{figures/section_ii__guitare__p47_20.jpg}
\caption{Illustration (page 47)}
\end{figure}

egin{figure}[htbp]
\centering
\includegraphics[width=\linewidth]{figures/section_ii__guitare__p47_21.jpg}
\caption{Illustration (page 47)}
\end{figure}

egin{figure}[htbp]
\centering
\includegraphics[width=\linewidth]{figures/section_ii__guitare__p47_22.jpg}
\caption{Illustration (page 47)}
\end{figure}

egin{figure}[htbp]
\centering
\includegraphics[width=\linewidth]{figures/section_ii__guitare__p47_23.jpg}
\caption{Illustration (page 47)}
\end{figure}

egin{figure}[htbp]
\centering
\includegraphics[width=\linewidth]{figures/section_ii__guitare__p47_24.jpg}
\caption{Illustration (page 47)}
\end{figure}

egin{figure}[htbp]
\centering
\includegraphics[width=\linewidth]{figures/section_ii__guitare__p47_25.jpg}
\caption{Illustration (page 47)}
\end{figure}

egin{figure}[htbp]
\centering
\includegraphics[width=\linewidth]{figures/section_ii__guitare__p47_26.jpg}
\caption{Illustration (page 47)}
\end{figure}

egin{figure}[htbp]
\centering
\includegraphics[width=\linewidth]{figures/section_ii__guitare__p47_27.png}
\caption{Illustration (page 47)}
\end{figure}

egin{figure}[htbp]
\centering
\includegraphics[width=\linewidth]{figures/section_ii__guitare__p47_28.png}
\caption{Illustration (page 47)}
\end{figure}

egin{figure}[htbp]
\centering
\includegraphics[width=\linewidth]{figures/section_ii__guitare__p47_29.png}
\caption{Illustration (page 47)}
\end{figure}

egin{figure}[htbp]
\centering
\includegraphics[width=\linewidth]{figures/section_ii__guitare__p47_30.png}
\caption{Illustration (page 47)}
\end{figure}

egin{figure}[htbp]
\centering
\includegraphics[width=\linewidth]{figures/section_ii__guitare__p47_31.png}
\caption{Illustration (page 47)}
\end{figure}

egin{figure}[htbp]
\centering
\includegraphics[width=\linewidth]{figures/section_ii__guitare__p47_32.jpg}
\caption{Illustration (page 47)}
\end{figure}

egin{figure}[htbp]
\centering
\includegraphics[width=\linewidth]{figures/section_ii__guitare__p47_33.png}
\caption{Illustration (page 47)}
\end{figure}
II.4 — Sourdines et variantes avec sourdine Les sourdines déplacent systématiquement les formants vers le grave (ratio ×0.5 à ×0.85 selon l'instrument et le type de sourdine). Ces fiches permettent d'analyser la transformation spectrale induite et de comparer aux valeurs ordinario correspondantes. Cuivres — sourdines Tuba basse + sourdine   [102 échantillons] TechniqueNF1F2F3F4F5F6Fiab. Ordinario1022264841012126018632132●●○ Fond vert = ordinario · Fond jaune = F1 dans cluster /o/ 420–550 Hz · Fond violet = technique étendue · Fiabilité: ●●● ≥0.55 / ●●○ 0.35–0.54 / ●○○ <0.35 Cor + sourdine   [87 échantillons] TechniqueNF1F2F3F4F5F6Fiab. Ordinario46344614948137817332143●●○ Flutter / Flatterzunge413346141001136717231938●●○ Fond vert = ordinario · Fond jaune = F1 dans cluster /o/ 420–550 Hz · Fond violet = technique étendue · Fiabilité: ●●● ≥0.55 / ●●○ 0.35–0.54 / ●○○ <0.35 Trombone + sourdine cup   [62 échantillons] TechniqueNF1F2F3F4F5F6Fiab. Ordinario313667641012133517442153●○○ Flutter / Flatterzunge314528941303168020132175●○○ Fond vert = ordinario · Fond jaune = F1 dans cluster /o/ 420–550 Hz · Fond violet = technique étendue · Fiabilité: ●●● ≥0.55 / ●●○ 0.35–0.54 / ●○○ <0.35 Trombone + sourdine harmon   [66 échantillons] TechniqueNF1F2F3F4F5F6Fiab. Ordinario351624411044185230473758●●○ Flutter / Flatterzunge311624411044130329393219●●○ Fond vert = ordinario · Fond jaune = F1 dans cluster /o/ 420–550 Hz · Fond violet = technique étendue · Fiabilité: ●●● ≥0.55 / ●●○ 0.35–0.54 / ●○○ <0.35 Trombone + sourdine straight   [62 échantillons] TechniqueNF1F2F3F4F5F6Fiab. Ordinario312265601238165820132218●●○ Flutter / Flatterzunge31248560872131419922369●●○ Fond vert = ordinario · Fond jaune = F1 dans cluster /o/ 420–550 Hz · Fond violet = technique étendue · Fiabilité: ●●● ≥0.55 / ●●○ 0.35–0.54 / ●○○ <0.35 Trombone + sourdine wah   [185 échantillons] TechniqueNF1F2F3F4F5F6Fiab. Ordinario62312630948129715772094●●○ Flutter / Flatterzunge61226516954142318952219●●○ Autres62328711996130818682137●○○ Fond vert = ordinario · Fond jaune = F1 dans cluster /o/ 420–550 Hz · Fond violet = technique étendue · Fiabilité: ●●● ≥0.55 / ●●○ 0.35–0.54 / ●○○ <0.35
\par\medskip

egin{figure}[htbp]
\centering
\includegraphics[width=\linewidth]{figures/section_ii__tuba__p48_01.jpg}
\caption{Illustration (page 48)}
\end{figure}

egin{figure}[htbp]
\centering
\includegraphics[width=\linewidth]{figures/section_ii__tuba__p48_02.jpg}
\caption{Illustration (page 48)}
\end{figure}

egin{figure}[htbp]
\centering
\includegraphics[width=\linewidth]{figures/section_ii__tuba__p48_03.png}
\caption{Illustration (page 48)}
\end{figure}

egin{figure}[htbp]
\centering
\includegraphics[width=\linewidth]{figures/section_ii__tuba__p48_04.jpg}
\caption{Illustration (page 48)}
\end{figure}

egin{figure}[htbp]
\centering
\includegraphics[width=\linewidth]{figures/section_ii__tuba__p48_05.jpg}
\caption{Illustration (page 48)}
\end{figure}

egin{figure}[htbp]
\centering
\includegraphics[width=\linewidth]{figures/section_ii__tuba__p48_06.jpg}
\caption{Illustration (page 48)}
\end{figure}

egin{figure}[htbp]
\centering
\includegraphics[width=\linewidth]{figures/section_ii__tuba__p48_07.jpg}
\caption{Illustration (page 48)}
\end{figure}

egin{figure}[htbp]
\centering
\includegraphics[width=\linewidth]{figures/section_ii__tuba__p48_08.png}
\caption{Illustration (page 48)}
\end{figure}

egin{figure}[htbp]
\centering
\includegraphics[width=\linewidth]{figures/section_ii__tuba__p48_09.jpg}
\caption{Illustration (page 48)}
\end{figure}

egin{figure}[htbp]
\centering
\includegraphics[width=\linewidth]{figures/section_ii__tuba__p48_10.jpg}
\caption{Illustration (page 48)}
\end{figure}

egin{figure}[htbp]
\centering
\includegraphics[width=\linewidth]{figures/section_ii__tuba__p48_11.png}
\caption{Illustration (page 48)}
\end{figure}

egin{figure}[htbp]
\centering
\includegraphics[width=\linewidth]{figures/section_ii__tuba__p48_12.png}
\caption{Illustration (page 48)}
\end{figure}

egin{figure}[htbp]
\centering
\includegraphics[width=\linewidth]{figures/section_ii__tuba__p48_13.png}
\caption{Illustration (page 48)}
\end{figure}

egin{figure}[htbp]
\centering
\includegraphics[width=\linewidth]{figures/section_ii__tuba__p48_14.png}
\caption{Illustration (page 48)}
\end{figure}

egin{figure}[htbp]
\centering
\includegraphics[width=\linewidth]{figures/section_ii__tuba__p48_15.png}
\caption{Illustration (page 48)}
\end{figure}

egin{figure}[htbp]
\centering
\includegraphics[width=\linewidth]{figures/section_ii__tuba__p48_16.png}
\caption{Illustration (page 48)}
\end{figure}

egin{figure}[htbp]
\centering
\includegraphics[width=\linewidth]{figures/section_ii__tuba__p48_17.png}
\caption{Illustration (page 48)}
\end{figure}

egin{figure}[htbp]
\centering
\includegraphics[width=\linewidth]{figures/section_ii__tuba__p48_18.png}
\caption{Illustration (page 48)}
\end{figure}

egin{figure}[htbp]
\centering
\includegraphics[width=\linewidth]{figures/section_ii__tuba__p48_19.png}
\caption{Illustration (page 48)}
\end{figure}

egin{figure}[htbp]
\centering
\includegraphics[width=\linewidth]{figures/section_ii__tuba__p48_20.jpg}
\caption{Illustration (page 48)}
\end{figure}

egin{figure}[htbp]
\centering
\includegraphics[width=\linewidth]{figures/section_ii__tuba__p48_21.jpg}
\caption{Illustration (page 48)}
\end{figure}

egin{figure}[htbp]
\centering
\includegraphics[width=\linewidth]{figures/section_ii__tuba__p48_22.jpg}
\caption{Illustration (page 48)}
\end{figure}

egin{figure}[htbp]
\centering
\includegraphics[width=\linewidth]{figures/section_ii__tuba__p48_23.jpg}
\caption{Illustration (page 48)}
\end{figure}

egin{figure}[htbp]
\centering
\includegraphics[width=\linewidth]{figures/section_ii__tuba__p48_24.jpg}
\caption{Illustration (page 48)}
\end{figure}

egin{figure}[htbp]
\centering
\includegraphics[width=\linewidth]{figures/section_ii__tuba__p48_25.jpg}
\caption{Illustration (page 48)}
\end{figure}

egin{figure}[htbp]
\centering
\includegraphics[width=\linewidth]{figures/section_ii__tuba__p48_26.jpg}
\caption{Illustration (page 48)}
\end{figure}

egin{figure}[htbp]
\centering
\includegraphics[width=\linewidth]{figures/section_ii__tuba__p48_27.png}
\caption{Illustration (page 48)}
\end{figure}

egin{figure}[htbp]
\centering
\includegraphics[width=\linewidth]{figures/section_ii__tuba__p48_28.png}
\caption{Illustration (page 48)}
\end{figure}

egin{figure}[htbp]
\centering
\includegraphics[width=\linewidth]{figures/section_ii__tuba__p48_29.png}
\caption{Illustration (page 48)}
\end{figure}

egin{figure}[htbp]
\centering
\includegraphics[width=\linewidth]{figures/section_ii__tuba__p48_30.png}
\caption{Illustration (page 48)}
\end{figure}

egin{figure}[htbp]
\centering
\includegraphics[width=\linewidth]{figures/section_ii__tuba__p48_31.png}
\caption{Illustration (page 48)}
\end{figure}

egin{figure}[htbp]
\centering
\includegraphics[width=\linewidth]{figures/section_ii__tuba__p48_32.jpg}
\caption{Illustration (page 48)}
\end{figure}

egin{figure}[htbp]
\centering
\includegraphics[width=\linewidth]{figures/section_ii__tuba__p48_33.png}
\caption{Illustration (page 48)}
\end{figure}
\subsection{Trompette + sourdine cup   [64 échantillons]}
Trompette + sourdine cup   [64 échantillons] TechniqueNF1F2F3F4F5F6Fiab. Ordinario34144318412336353147165308●○○ Flutter / Flatterzunge30142118632390349937585243●○○ Fond vert = ordinario · Fond jaune = F1 dans cluster /o/ 420–550 Hz · Fond violet = technique étendue · Fiabilité: ●●● ≥0.55 / ●●○ 0.35–0.54 / ●○○ <0.35 Trompette + sourdine harmon   [65 échantillons] TechniqueNF1F2F3F4F5F6Fiab. Ordinario34235833053962459749855588●○○ Flutter / Flatterzunge31211029823478419947375082●○○ Fond vert = ordinario · Fond jaune = F1 dans cluster /o/ 420–550 Hz · Fond violet = technique étendue · Fiabilité: ●●● ≥0.55 / ●●○ 0.35–0.54 / ●○○ <0.35 Trompette + sourdine straight   [65 échantillons] TechniqueNF1F2F3F4F5F6Fiab. Ordinario34109816582218279935213973●○○ Flutter / Flatterzunge31104415941992231526813273●○○ Fond vert = ordinario · Fond jaune = F1 dans cluster /o/ 420–550 Hz · Fond violet = technique étendue · Fiabilité: ●●● ≥0.55 / ●●○ 0.35–0.54 / ●○○ <0.35 Trompette + sourdine wah   [163 échantillons] TechniqueNF1F2F3F4F5F6Fiab. Ordinario6757010531563195126304004●○○ Flutter / Flatterzunge304849901443176622072606●○○ Autres6698514641863241737095055●○○ Fond vert = ordinario · Fond jaune = F1 dans cluster /o/ 420–550 Hz · Fond violet = technique étendue · Fiabilité: ●●● ≥0.55 / ●●○ 0.35–0.54 / ●○○ <0.35 Cordes — sourdines Violon + sourdine   [614 échantillons] TechniqueNF1F2F3F4F5F6Fiab. Ordinario2153668832616352142204770●●○ Vibrato2063668402304347841994716●●● Trémolo1933668511733316539734640●●○ Fond vert = ordinario · Fond jaune = F1 dans cluster /o/ 420–550 Hz · Fond violet = technique étendue · Fiabilité: ●●● ≥0.55 / ●●○ 0.35–0.54 / ●○○ <0.35 Violon + sourdine piombo   [122 échantillons] TechniqueNF1F2F3F4F5F6Fiab. Ordinario4034411632207332741775491●○○ Vibrato4034413142358345641455545●○○ Trémolo4244113892229312241564974●○○ Fond vert = ordinario · Fond jaune = F1 dans cluster /o/ 420–550 Hz · Fond violet = technique étendue · Fiabilité: ●●● ≥0.55 / ●●○ 0.35–0.54 / ●○○ <0.35
\par\medskip

egin{figure}[htbp]
\centering
\includegraphics[width=\linewidth]{figures/section_ii__trompette__p49_01.jpg}
\caption{Illustration (page 49)}
\end{figure}

egin{figure}[htbp]
\centering
\includegraphics[width=\linewidth]{figures/section_ii__trompette__p49_02.jpg}
\caption{Illustration (page 49)}
\end{figure}

egin{figure}[htbp]
\centering
\includegraphics[width=\linewidth]{figures/section_ii__trompette__p49_03.png}
\caption{Illustration (page 49)}
\end{figure}

egin{figure}[htbp]
\centering
\includegraphics[width=\linewidth]{figures/section_ii__trompette__p49_04.jpg}
\caption{Illustration (page 49)}
\end{figure}

egin{figure}[htbp]
\centering
\includegraphics[width=\linewidth]{figures/section_ii__trompette__p49_05.jpg}
\caption{Illustration (page 49)}
\end{figure}

egin{figure}[htbp]
\centering
\includegraphics[width=\linewidth]{figures/section_ii__trompette__p49_06.jpg}
\caption{Illustration (page 49)}
\end{figure}

egin{figure}[htbp]
\centering
\includegraphics[width=\linewidth]{figures/section_ii__trompette__p49_07.jpg}
\caption{Illustration (page 49)}
\end{figure}

egin{figure}[htbp]
\centering
\includegraphics[width=\linewidth]{figures/section_ii__trompette__p49_08.png}
\caption{Illustration (page 49)}
\end{figure}

egin{figure}[htbp]
\centering
\includegraphics[width=\linewidth]{figures/section_ii__trompette__p49_09.jpg}
\caption{Illustration (page 49)}
\end{figure}

egin{figure}[htbp]
\centering
\includegraphics[width=\linewidth]{figures/section_ii__trompette__p49_10.jpg}
\caption{Illustration (page 49)}
\end{figure}

egin{figure}[htbp]
\centering
\includegraphics[width=\linewidth]{figures/section_ii__trompette__p49_11.png}
\caption{Illustration (page 49)}
\end{figure}

egin{figure}[htbp]
\centering
\includegraphics[width=\linewidth]{figures/section_ii__trompette__p49_12.png}
\caption{Illustration (page 49)}
\end{figure}

egin{figure}[htbp]
\centering
\includegraphics[width=\linewidth]{figures/section_ii__trompette__p49_13.png}
\caption{Illustration (page 49)}
\end{figure}

egin{figure}[htbp]
\centering
\includegraphics[width=\linewidth]{figures/section_ii__trompette__p49_14.png}
\caption{Illustration (page 49)}
\end{figure}

egin{figure}[htbp]
\centering
\includegraphics[width=\linewidth]{figures/section_ii__trompette__p49_15.png}
\caption{Illustration (page 49)}
\end{figure}

egin{figure}[htbp]
\centering
\includegraphics[width=\linewidth]{figures/section_ii__trompette__p49_16.png}
\caption{Illustration (page 49)}
\end{figure}

egin{figure}[htbp]
\centering
\includegraphics[width=\linewidth]{figures/section_ii__trompette__p49_17.png}
\caption{Illustration (page 49)}
\end{figure}

egin{figure}[htbp]
\centering
\includegraphics[width=\linewidth]{figures/section_ii__trompette__p49_18.png}
\caption{Illustration (page 49)}
\end{figure}

egin{figure}[htbp]
\centering
\includegraphics[width=\linewidth]{figures/section_ii__trompette__p49_19.png}
\caption{Illustration (page 49)}
\end{figure}

egin{figure}[htbp]
\centering
\includegraphics[width=\linewidth]{figures/section_ii__trompette__p49_20.jpg}
\caption{Illustration (page 49)}
\end{figure}

egin{figure}[htbp]
\centering
\includegraphics[width=\linewidth]{figures/section_ii__trompette__p49_21.jpg}
\caption{Illustration (page 49)}
\end{figure}

egin{figure}[htbp]
\centering
\includegraphics[width=\linewidth]{figures/section_ii__trompette__p49_22.jpg}
\caption{Illustration (page 49)}
\end{figure}

egin{figure}[htbp]
\centering
\includegraphics[width=\linewidth]{figures/section_ii__trompette__p49_23.jpg}
\caption{Illustration (page 49)}
\end{figure}

egin{figure}[htbp]
\centering
\includegraphics[width=\linewidth]{figures/section_ii__trompette__p49_24.jpg}
\caption{Illustration (page 49)}
\end{figure}

egin{figure}[htbp]
\centering
\includegraphics[width=\linewidth]{figures/section_ii__trompette__p49_25.jpg}
\caption{Illustration (page 49)}
\end{figure}

egin{figure}[htbp]
\centering
\includegraphics[width=\linewidth]{figures/section_ii__trompette__p49_26.jpg}
\caption{Illustration (page 49)}
\end{figure}

egin{figure}[htbp]
\centering
\includegraphics[width=\linewidth]{figures/section_ii__trompette__p49_27.png}
\caption{Illustration (page 49)}
\end{figure}

egin{figure}[htbp]
\centering
\includegraphics[width=\linewidth]{figures/section_ii__trompette__p49_28.png}
\caption{Illustration (page 49)}
\end{figure}

egin{figure}[htbp]
\centering
\includegraphics[width=\linewidth]{figures/section_ii__trompette__p49_29.png}
\caption{Illustration (page 49)}
\end{figure}

egin{figure}[htbp]
\centering
\includegraphics[width=\linewidth]{figures/section_ii__trompette__p49_30.png}
\caption{Illustration (page 49)}
\end{figure}

egin{figure}[htbp]
\centering
\includegraphics[width=\linewidth]{figures/section_ii__trompette__p49_31.png}
\caption{Illustration (page 49)}
\end{figure}

egin{figure}[htbp]
\centering
\includegraphics[width=\linewidth]{figures/section_ii__trompette__p49_32.jpg}
\caption{Illustration (page 49)}
\end{figure}

egin{figure}[htbp]
\centering
\includegraphics[width=\linewidth]{figures/section_ii__trompette__p49_33.png}
\caption{Illustration (page 49)}
\end{figure}
\subsection{Alto + sourdine   [380 échantillons]}
Alto + sourdine   [380 échantillons] TechniqueNF1F2F3F4F5F6Fiab. Ordinario1683447211098164721753144●●○ Vibrato1563347111077160421532972●●● Trémolo563447321098163621963101●●○ Fond vert = ordinario · Fond jaune = F1 dans cluster /o/ 420–550 Hz · Fond violet = technique étendue · Fiabilité: ●●● ≥0.55 / ●●○ 0.35–0.54 / ●○○ <0.35 Alto + sourdine piombo   [117 échantillons] TechniqueNF1F2F3F4F5F6Fiab. Ordinario382267001249248741885297●●○ Vibrato362155711077178731335254●●● Trémolo432157431174232633054694●●○ Fond vert = ordinario · Fond jaune = F1 dans cluster /o/ 420–550 Hz · Fond violet = technique étendue · Fiabilité: ●●● ≥0.55 / ●●○ 0.35–0.54 / ●○○ <0.35 Violoncelle + sourdine   [352 échantillons] TechniqueNF1F2F3F4F5F6Fiab. Ordinario1562055381174189522184091●○○ Vibrato1441945061034177621102939●●● Trémolo522155171098154020672939●●○ Fond vert = ordinario · Fond jaune = F1 dans cluster /o/ 420–550 Hz · Fond violet = technique étendue · Fiabilité: ●●● ≥0.55 / ●●○ 0.35–0.54 / ●○○ <0.35 Violoncelle + sourdine piombo   [149 échantillons] TechniqueNF1F2F3F4F5F6Fiab. Ordinario521727321507354242535739●●○ Vibrato451725811044235839415760●●○ Trémolo521946031044173325094619●●○ Fond vert = ordinario · Fond jaune = F1 dans cluster /o/ 420–550 Hz · Fond violet = technique étendue · Fiabilité: ●●● ≥0.55 / ●●○ 0.35–0.54 / ●○○ <0.35 Contrebasse + sourdine   [351 échantillons] TechniqueNF1F2F3F4F5F6Fiab. Ordinario156162474840113015501906●●● Vibrato143172474840116315611938●●● Trémolo52183474840109815402067●●● Fond vert = ordinario · Fond jaune = F1 dans cluster /o/ 420–550 Hz · Fond violet = technique étendue · Fiabilité: ●●● ≥0.55 / ●●○ 0.35–0.54 / ●○○ <0.35 Bois — sourdines Basson + sourdine   [41 échantillons] TechniqueNF1F2F3F4F5F6Fiab. Ordinario414848831249155018304436●●○
\par\medskip

egin{figure}[htbp]
\centering
\includegraphics[width=\linewidth]{figures/section_ii__alto__p50_01.jpg}
\caption{Illustration (page 50)}
\end{figure}

egin{figure}[htbp]
\centering
\includegraphics[width=\linewidth]{figures/section_ii__alto__p50_02.jpg}
\caption{Illustration (page 50)}
\end{figure}

egin{figure}[htbp]
\centering
\includegraphics[width=\linewidth]{figures/section_ii__alto__p50_03.png}
\caption{Illustration (page 50)}
\end{figure}

egin{figure}[htbp]
\centering
\includegraphics[width=\linewidth]{figures/section_ii__alto__p50_04.jpg}
\caption{Illustration (page 50)}
\end{figure}

egin{figure}[htbp]
\centering
\includegraphics[width=\linewidth]{figures/section_ii__alto__p50_05.jpg}
\caption{Illustration (page 50)}
\end{figure}

egin{figure}[htbp]
\centering
\includegraphics[width=\linewidth]{figures/section_ii__alto__p50_06.jpg}
\caption{Illustration (page 50)}
\end{figure}

egin{figure}[htbp]
\centering
\includegraphics[width=\linewidth]{figures/section_ii__alto__p50_07.jpg}
\caption{Illustration (page 50)}
\end{figure}

egin{figure}[htbp]
\centering
\includegraphics[width=\linewidth]{figures/section_ii__alto__p50_08.png}
\caption{Illustration (page 50)}
\end{figure}

egin{figure}[htbp]
\centering
\includegraphics[width=\linewidth]{figures/section_ii__alto__p50_09.jpg}
\caption{Illustration (page 50)}
\end{figure}

egin{figure}[htbp]
\centering
\includegraphics[width=\linewidth]{figures/section_ii__alto__p50_10.jpg}
\caption{Illustration (page 50)}
\end{figure}

egin{figure}[htbp]
\centering
\includegraphics[width=\linewidth]{figures/section_ii__alto__p50_11.png}
\caption{Illustration (page 50)}
\end{figure}

egin{figure}[htbp]
\centering
\includegraphics[width=\linewidth]{figures/section_ii__alto__p50_12.png}
\caption{Illustration (page 50)}
\end{figure}

egin{figure}[htbp]
\centering
\includegraphics[width=\linewidth]{figures/section_ii__alto__p50_13.png}
\caption{Illustration (page 50)}
\end{figure}

egin{figure}[htbp]
\centering
\includegraphics[width=\linewidth]{figures/section_ii__alto__p50_14.png}
\caption{Illustration (page 50)}
\end{figure}

egin{figure}[htbp]
\centering
\includegraphics[width=\linewidth]{figures/section_ii__alto__p50_15.png}
\caption{Illustration (page 50)}
\end{figure}

egin{figure}[htbp]
\centering
\includegraphics[width=\linewidth]{figures/section_ii__alto__p50_16.png}
\caption{Illustration (page 50)}
\end{figure}

egin{figure}[htbp]
\centering
\includegraphics[width=\linewidth]{figures/section_ii__alto__p50_17.png}
\caption{Illustration (page 50)}
\end{figure}

egin{figure}[htbp]
\centering
\includegraphics[width=\linewidth]{figures/section_ii__alto__p50_18.png}
\caption{Illustration (page 50)}
\end{figure}

egin{figure}[htbp]
\centering
\includegraphics[width=\linewidth]{figures/section_ii__alto__p50_19.png}
\caption{Illustration (page 50)}
\end{figure}

egin{figure}[htbp]
\centering
\includegraphics[width=\linewidth]{figures/section_ii__alto__p50_20.jpg}
\caption{Illustration (page 50)}
\end{figure}

egin{figure}[htbp]
\centering
\includegraphics[width=\linewidth]{figures/section_ii__alto__p50_21.jpg}
\caption{Illustration (page 50)}
\end{figure}

egin{figure}[htbp]
\centering
\includegraphics[width=\linewidth]{figures/section_ii__alto__p50_22.jpg}
\caption{Illustration (page 50)}
\end{figure}

egin{figure}[htbp]
\centering
\includegraphics[width=\linewidth]{figures/section_ii__alto__p50_23.jpg}
\caption{Illustration (page 50)}
\end{figure}

egin{figure}[htbp]
\centering
\includegraphics[width=\linewidth]{figures/section_ii__alto__p50_24.jpg}
\caption{Illustration (page 50)}
\end{figure}

egin{figure}[htbp]
\centering
\includegraphics[width=\linewidth]{figures/section_ii__alto__p50_25.jpg}
\caption{Illustration (page 50)}
\end{figure}

egin{figure}[htbp]
\centering
\includegraphics[width=\linewidth]{figures/section_ii__alto__p50_26.jpg}
\caption{Illustration (page 50)}
\end{figure}

egin{figure}[htbp]
\centering
\includegraphics[width=\linewidth]{figures/section_ii__alto__p50_27.png}
\caption{Illustration (page 50)}
\end{figure}

egin{figure}[htbp]
\centering
\includegraphics[width=\linewidth]{figures/section_ii__alto__p50_28.png}
\caption{Illustration (page 50)}
\end{figure}

egin{figure}[htbp]
\centering
\includegraphics[width=\linewidth]{figures/section_ii__alto__p50_29.png}
\caption{Illustration (page 50)}
\end{figure}

egin{figure}[htbp]
\centering
\includegraphics[width=\linewidth]{figures/section_ii__alto__p50_30.png}
\caption{Illustration (page 50)}
\end{figure}

egin{figure}[htbp]
\centering
\includegraphics[width=\linewidth]{figures/section_ii__alto__p50_31.png}
\caption{Illustration (page 50)}
\end{figure}

egin{figure}[htbp]
\centering
\includegraphics[width=\linewidth]{figures/section_ii__alto__p50_32.jpg}
\caption{Illustration (page 50)}
\end{figure}

egin{figure}[htbp]
\centering
\includegraphics[width=\linewidth]{figures/section_ii__alto__p50_33.png}
\caption{Illustration (page 50)}
\end{figure}
\subsection{Hautbois + sourdine   [36 échantillons]}
Fond vert = ordinario · Fond jaune = F1 dans cluster /o/ 420–550 Hz · Fond violet = technique étendue · Fiabilité: ●●● ≥0.55 / ●●○ 0.35–0.54 / ●○○ <0.35 Hautbois + sourdine   [36 échantillons] TechniqueNF1F2F3F4F5F6Fiab. Ordinario3671112811658200324873521●○○ Fond vert = ordinario · Fond jaune = F1 dans cluster /o/ 420–550 Hz · Fond violet = technique étendue · Fiabilité: ●●● ≥0.55 / ●●○ 0.35–0.54 / ●○○ <0.35
\par\medskip

egin{figure}[htbp]
\centering
\includegraphics[width=\linewidth]{figures/section_ii__hautbois__p51_01.jpg}
\caption{Illustration (page 51)}
\end{figure}

egin{figure}[htbp]
\centering
\includegraphics[width=\linewidth]{figures/section_ii__hautbois__p51_02.jpg}
\caption{Illustration (page 51)}
\end{figure}

egin{figure}[htbp]
\centering
\includegraphics[width=\linewidth]{figures/section_ii__hautbois__p51_03.png}
\caption{Illustration (page 51)}
\end{figure}

egin{figure}[htbp]
\centering
\includegraphics[width=\linewidth]{figures/section_ii__hautbois__p51_04.jpg}
\caption{Illustration (page 51)}
\end{figure}

egin{figure}[htbp]
\centering
\includegraphics[width=\linewidth]{figures/section_ii__hautbois__p51_05.jpg}
\caption{Illustration (page 51)}
\end{figure}

egin{figure}[htbp]
\centering
\includegraphics[width=\linewidth]{figures/section_ii__hautbois__p51_06.jpg}
\caption{Illustration (page 51)}
\end{figure}

egin{figure}[htbp]
\centering
\includegraphics[width=\linewidth]{figures/section_ii__hautbois__p51_07.jpg}
\caption{Illustration (page 51)}
\end{figure}

egin{figure}[htbp]
\centering
\includegraphics[width=\linewidth]{figures/section_ii__hautbois__p51_08.png}
\caption{Illustration (page 51)}
\end{figure}

egin{figure}[htbp]
\centering
\includegraphics[width=\linewidth]{figures/section_ii__hautbois__p51_09.jpg}
\caption{Illustration (page 51)}
\end{figure}

egin{figure}[htbp]
\centering
\includegraphics[width=\linewidth]{figures/section_ii__hautbois__p51_10.jpg}
\caption{Illustration (page 51)}
\end{figure}

egin{figure}[htbp]
\centering
\includegraphics[width=\linewidth]{figures/section_ii__hautbois__p51_11.png}
\caption{Illustration (page 51)}
\end{figure}

egin{figure}[htbp]
\centering
\includegraphics[width=\linewidth]{figures/section_ii__hautbois__p51_12.png}
\caption{Illustration (page 51)}
\end{figure}

egin{figure}[htbp]
\centering
\includegraphics[width=\linewidth]{figures/section_ii__hautbois__p51_13.png}
\caption{Illustration (page 51)}
\end{figure}

egin{figure}[htbp]
\centering
\includegraphics[width=\linewidth]{figures/section_ii__hautbois__p51_14.png}
\caption{Illustration (page 51)}
\end{figure}

egin{figure}[htbp]
\centering
\includegraphics[width=\linewidth]{figures/section_ii__hautbois__p51_15.png}
\caption{Illustration (page 51)}
\end{figure}

egin{figure}[htbp]
\centering
\includegraphics[width=\linewidth]{figures/section_ii__hautbois__p51_16.png}
\caption{Illustration (page 51)}
\end{figure}

egin{figure}[htbp]
\centering
\includegraphics[width=\linewidth]{figures/section_ii__hautbois__p51_17.png}
\caption{Illustration (page 51)}
\end{figure}

egin{figure}[htbp]
\centering
\includegraphics[width=\linewidth]{figures/section_ii__hautbois__p51_18.png}
\caption{Illustration (page 51)}
\end{figure}

egin{figure}[htbp]
\centering
\includegraphics[width=\linewidth]{figures/section_ii__hautbois__p51_19.png}
\caption{Illustration (page 51)}
\end{figure}

egin{figure}[htbp]
\centering
\includegraphics[width=\linewidth]{figures/section_ii__hautbois__p51_20.jpg}
\caption{Illustration (page 51)}
\end{figure}

egin{figure}[htbp]
\centering
\includegraphics[width=\linewidth]{figures/section_ii__hautbois__p51_21.jpg}
\caption{Illustration (page 51)}
\end{figure}

egin{figure}[htbp]
\centering
\includegraphics[width=\linewidth]{figures/section_ii__hautbois__p51_22.jpg}
\caption{Illustration (page 51)}
\end{figure}

egin{figure}[htbp]
\centering
\includegraphics[width=\linewidth]{figures/section_ii__hautbois__p51_23.jpg}
\caption{Illustration (page 51)}
\end{figure}

egin{figure}[htbp]
\centering
\includegraphics[width=\linewidth]{figures/section_ii__hautbois__p51_24.jpg}
\caption{Illustration (page 51)}
\end{figure}

egin{figure}[htbp]
\centering
\includegraphics[width=\linewidth]{figures/section_ii__hautbois__p51_25.jpg}
\caption{Illustration (page 51)}
\end{figure}

egin{figure}[htbp]
\centering
\includegraphics[width=\linewidth]{figures/section_ii__hautbois__p51_26.jpg}
\caption{Illustration (page 51)}
\end{figure}

egin{figure}[htbp]
\centering
\includegraphics[width=\linewidth]{figures/section_ii__hautbois__p51_27.png}
\caption{Illustration (page 51)}
\end{figure}

egin{figure}[htbp]
\centering
\includegraphics[width=\linewidth]{figures/section_ii__hautbois__p51_28.png}
\caption{Illustration (page 51)}
\end{figure}

egin{figure}[htbp]
\centering
\includegraphics[width=\linewidth]{figures/section_ii__hautbois__p51_29.png}
\caption{Illustration (page 51)}
\end{figure}

egin{figure}[htbp]
\centering
\includegraphics[width=\linewidth]{figures/section_ii__hautbois__p51_30.png}
\caption{Illustration (page 51)}
\end{figure}

egin{figure}[htbp]
\centering
\includegraphics[width=\linewidth]{figures/section_ii__hautbois__p51_31.png}
\caption{Illustration (page 51)}
\end{figure}

egin{figure}[htbp]
\centering
\includegraphics[width=\linewidth]{figures/section_ii__hautbois__p51_32.jpg}
\caption{Illustration (page 51)}
\end{figure}

egin{figure}[htbp]
\centering
\includegraphics[width=\linewidth]{figures/section_ii__hautbois__p51_33.png}
\caption{Illustration (page 51)}
\end{figure}
II.5 — Synthèse: F1 spectral ordinario — Tous instruments Tableau récapitulatif des F1 à F6 spectraux (ordinario pondéré) pour les 16 instruments principaux, triés par F1 croissant. Les cases jaunes indiquent les F1 dans le cluster de convergence /o/ (420–550 Hz). InstrumentNF1F2F3F4F5F6Fp Contrebasse637176481852119014922209132 4 Guitare240194612981148022353021100 4 Violoncelle5662095221139177221172541102 5 Tuba basse2242354731003129716021956123 9 Harpe2412696241034145423694156136 0 Trombone ténor226292573871134717332064121 8 Alto7493697921463212828033359122 1 Cor en Fa2704197351188167021952885738 Accordéon6894209261486208926703962144 0 Violon67945012302177296937884546893 Basson2454698841217165820292825107 9 Clarinette Sib28748911611597190525663284141 2 Saxophone alto29253610311599207825823065144 0 Flûte45668011361557196323852925153 5 Hautbois20975312751716228828313851148 5 Trompette en Do16990214421940247730633737104 6 Piccolo11 1110 9199 2263 8333 8876 Flûte basse39302668958140 0133 8 Flûte c.basse37334678103 4142 1109 2 Cor anglais12 8452894124 9166 9113 5 Clarinette Mib45678154 0196 0234 7126 6 Cl. basse Sib17 3323808126 0164 7120 4 Cl. c.basse Sib45323937124 9164 7109 0
\par\medskip

egin{figure}[htbp]
\centering
\includegraphics[width=\linewidth]{figures/section_ii__contrebasse__p52_01.jpg}
\caption{Illustration (page 52)}
\end{figure}

egin{figure}[htbp]
\centering
\includegraphics[width=\linewidth]{figures/section_ii__contrebasse__p52_02.jpg}
\caption{Illustration (page 52)}
\end{figure}

egin{figure}[htbp]
\centering
\includegraphics[width=\linewidth]{figures/section_ii__contrebasse__p52_03.png}
\caption{Illustration (page 52)}
\end{figure}

egin{figure}[htbp]
\centering
\includegraphics[width=\linewidth]{figures/section_ii__contrebasse__p52_04.jpg}
\caption{Illustration (page 52)}
\end{figure}

egin{figure}[htbp]
\centering
\includegraphics[width=\linewidth]{figures/section_ii__contrebasse__p52_05.jpg}
\caption{Illustration (page 52)}
\end{figure}

egin{figure}[htbp]
\centering
\includegraphics[width=\linewidth]{figures/section_ii__contrebasse__p52_06.jpg}
\caption{Illustration (page 52)}
\end{figure}

egin{figure}[htbp]
\centering
\includegraphics[width=\linewidth]{figures/section_ii__contrebasse__p52_07.jpg}
\caption{Illustration (page 52)}
\end{figure}

egin{figure}[htbp]
\centering
\includegraphics[width=\linewidth]{figures/section_ii__contrebasse__p52_08.png}
\caption{Illustration (page 52)}
\end{figure}

egin{figure}[htbp]
\centering
\includegraphics[width=\linewidth]{figures/section_ii__contrebasse__p52_09.jpg}
\caption{Illustration (page 52)}
\end{figure}

egin{figure}[htbp]
\centering
\includegraphics[width=\linewidth]{figures/section_ii__contrebasse__p52_10.jpg}
\caption{Illustration (page 52)}
\end{figure}

egin{figure}[htbp]
\centering
\includegraphics[width=\linewidth]{figures/section_ii__contrebasse__p52_11.png}
\caption{Illustration (page 52)}
\end{figure}

egin{figure}[htbp]
\centering
\includegraphics[width=\linewidth]{figures/section_ii__contrebasse__p52_12.png}
\caption{Illustration (page 52)}
\end{figure}

egin{figure}[htbp]
\centering
\includegraphics[width=\linewidth]{figures/section_ii__contrebasse__p52_13.png}
\caption{Illustration (page 52)}
\end{figure}

egin{figure}[htbp]
\centering
\includegraphics[width=\linewidth]{figures/section_ii__contrebasse__p52_14.png}
\caption{Illustration (page 52)}
\end{figure}

egin{figure}[htbp]
\centering
\includegraphics[width=\linewidth]{figures/section_ii__contrebasse__p52_15.png}
\caption{Illustration (page 52)}
\end{figure}

egin{figure}[htbp]
\centering
\includegraphics[width=\linewidth]{figures/section_ii__contrebasse__p52_16.png}
\caption{Illustration (page 52)}
\end{figure}

egin{figure}[htbp]
\centering
\includegraphics[width=\linewidth]{figures/section_ii__contrebasse__p52_17.png}
\caption{Illustration (page 52)}
\end{figure}

egin{figure}[htbp]
\centering
\includegraphics[width=\linewidth]{figures/section_ii__contrebasse__p52_18.png}
\caption{Illustration (page 52)}
\end{figure}

egin{figure}[htbp]
\centering
\includegraphics[width=\linewidth]{figures/section_ii__contrebasse__p52_19.png}
\caption{Illustration (page 52)}
\end{figure}

egin{figure}[htbp]
\centering
\includegraphics[width=\linewidth]{figures/section_ii__contrebasse__p52_20.jpg}
\caption{Illustration (page 52)}
\end{figure}

egin{figure}[htbp]
\centering
\includegraphics[width=\linewidth]{figures/section_ii__contrebasse__p52_21.jpg}
\caption{Illustration (page 52)}
\end{figure}

egin{figure}[htbp]
\centering
\includegraphics[width=\linewidth]{figures/section_ii__contrebasse__p52_22.jpg}
\caption{Illustration (page 52)}
\end{figure}

egin{figure}[htbp]
\centering
\includegraphics[width=\linewidth]{figures/section_ii__contrebasse__p52_23.jpg}
\caption{Illustration (page 52)}
\end{figure}

egin{figure}[htbp]
\centering
\includegraphics[width=\linewidth]{figures/section_ii__contrebasse__p52_24.jpg}
\caption{Illustration (page 52)}
\end{figure}

egin{figure}[htbp]
\centering
\includegraphics[width=\linewidth]{figures/section_ii__contrebasse__p52_25.jpg}
\caption{Illustration (page 52)}
\end{figure}

egin{figure}[htbp]
\centering
\includegraphics[width=\linewidth]{figures/section_ii__contrebasse__p52_26.jpg}
\caption{Illustration (page 52)}
\end{figure}

egin{figure}[htbp]
\centering
\includegraphics[width=\linewidth]{figures/section_ii__contrebasse__p52_27.png}
\caption{Illustration (page 52)}
\end{figure}

egin{figure}[htbp]
\centering
\includegraphics[width=\linewidth]{figures/section_ii__contrebasse__p52_28.png}
\caption{Illustration (page 52)}
\end{figure}

egin{figure}[htbp]
\centering
\includegraphics[width=\linewidth]{figures/section_ii__contrebasse__p52_29.png}
\caption{Illustration (page 52)}
\end{figure}

egin{figure}[htbp]
\centering
\includegraphics[width=\linewidth]{figures/section_ii__contrebasse__p52_30.png}
\caption{Illustration (page 52)}
\end{figure}

egin{figure}[htbp]
\centering
\includegraphics[width=\linewidth]{figures/section_ii__contrebasse__p52_31.png}
\caption{Illustration (page 52)}
\end{figure}

egin{figure}[htbp]
\centering
\includegraphics[width=\linewidth]{figures/section_ii__contrebasse__p52_32.jpg}
\caption{Illustration (page 52)}
\end{figure}

egin{figure}[htbp]
\centering
\includegraphics[width=\linewidth]{figures/section_ii__contrebasse__p52_33.png}
\caption{Illustration (page 52)}
\end{figure}
\subsection{Contrebasson*10}
Contrebasson*10 6226657101 2123 8127 9 Trb. basse14 4258538818117 4133 5 Tuba c.basse13 5226463732126 0118 2 Cases jaunes: F1 spectral dans le cluster de convergence /o/ (420–550 Hz)
\par\medskip

egin{figure}[htbp]
\centering
\includegraphics[width=\linewidth]{figures/section_ii__contrebasson__p53_01.jpg}
\caption{Illustration (page 53)}
\end{figure}

egin{figure}[htbp]
\centering
\includegraphics[width=\linewidth]{figures/section_ii__contrebasson__p53_02.jpg}
\caption{Illustration (page 53)}
\end{figure}

egin{figure}[htbp]
\centering
\includegraphics[width=\linewidth]{figures/section_ii__contrebasson__p53_03.png}
\caption{Illustration (page 53)}
\end{figure}

egin{figure}[htbp]
\centering
\includegraphics[width=\linewidth]{figures/section_ii__contrebasson__p53_04.jpg}
\caption{Illustration (page 53)}
\end{figure}

egin{figure}[htbp]
\centering
\includegraphics[width=\linewidth]{figures/section_ii__contrebasson__p53_05.jpg}
\caption{Illustration (page 53)}
\end{figure}

egin{figure}[htbp]
\centering
\includegraphics[width=\linewidth]{figures/section_ii__contrebasson__p53_06.jpg}
\caption{Illustration (page 53)}
\end{figure}

egin{figure}[htbp]
\centering
\includegraphics[width=\linewidth]{figures/section_ii__contrebasson__p53_07.jpg}
\caption{Illustration (page 53)}
\end{figure}

egin{figure}[htbp]
\centering
\includegraphics[width=\linewidth]{figures/section_ii__contrebasson__p53_08.png}
\caption{Illustration (page 53)}
\end{figure}

egin{figure}[htbp]
\centering
\includegraphics[width=\linewidth]{figures/section_ii__contrebasson__p53_09.jpg}
\caption{Illustration (page 53)}
\end{figure}

egin{figure}[htbp]
\centering
\includegraphics[width=\linewidth]{figures/section_ii__contrebasson__p53_10.jpg}
\caption{Illustration (page 53)}
\end{figure}

egin{figure}[htbp]
\centering
\includegraphics[width=\linewidth]{figures/section_ii__contrebasson__p53_11.png}
\caption{Illustration (page 53)}
\end{figure}

egin{figure}[htbp]
\centering
\includegraphics[width=\linewidth]{figures/section_ii__contrebasson__p53_12.png}
\caption{Illustration (page 53)}
\end{figure}

egin{figure}[htbp]
\centering
\includegraphics[width=\linewidth]{figures/section_ii__contrebasson__p53_13.png}
\caption{Illustration (page 53)}
\end{figure}

egin{figure}[htbp]
\centering
\includegraphics[width=\linewidth]{figures/section_ii__contrebasson__p53_14.png}
\caption{Illustration (page 53)}
\end{figure}

egin{figure}[htbp]
\centering
\includegraphics[width=\linewidth]{figures/section_ii__contrebasson__p53_15.png}
\caption{Illustration (page 53)}
\end{figure}

egin{figure}[htbp]
\centering
\includegraphics[width=\linewidth]{figures/section_ii__contrebasson__p53_16.png}
\caption{Illustration (page 53)}
\end{figure}

egin{figure}[htbp]
\centering
\includegraphics[width=\linewidth]{figures/section_ii__contrebasson__p53_17.png}
\caption{Illustration (page 53)}
\end{figure}

egin{figure}[htbp]
\centering
\includegraphics[width=\linewidth]{figures/section_ii__contrebasson__p53_18.png}
\caption{Illustration (page 53)}
\end{figure}

egin{figure}[htbp]
\centering
\includegraphics[width=\linewidth]{figures/section_ii__contrebasson__p53_19.png}
\caption{Illustration (page 53)}
\end{figure}

egin{figure}[htbp]
\centering
\includegraphics[width=\linewidth]{figures/section_ii__contrebasson__p53_20.jpg}
\caption{Illustration (page 53)}
\end{figure}

egin{figure}[htbp]
\centering
\includegraphics[width=\linewidth]{figures/section_ii__contrebasson__p53_21.jpg}
\caption{Illustration (page 53)}
\end{figure}

egin{figure}[htbp]
\centering
\includegraphics[width=\linewidth]{figures/section_ii__contrebasson__p53_22.jpg}
\caption{Illustration (page 53)}
\end{figure}

egin{figure}[htbp]
\centering
\includegraphics[width=\linewidth]{figures/section_ii__contrebasson__p53_23.jpg}
\caption{Illustration (page 53)}
\end{figure}

egin{figure}[htbp]
\centering
\includegraphics[width=\linewidth]{figures/section_ii__contrebasson__p53_24.jpg}
\caption{Illustration (page 53)}
\end{figure}

egin{figure}[htbp]
\centering
\includegraphics[width=\linewidth]{figures/section_ii__contrebasson__p53_25.jpg}
\caption{Illustration (page 53)}
\end{figure}

egin{figure}[htbp]
\centering
\includegraphics[width=\linewidth]{figures/section_ii__contrebasson__p53_26.jpg}
\caption{Illustration (page 53)}
\end{figure}

egin{figure}[htbp]
\centering
\includegraphics[width=\linewidth]{figures/section_ii__contrebasson__p53_27.png}
\caption{Illustration (page 53)}
\end{figure}

egin{figure}[htbp]
\centering
\includegraphics[width=\linewidth]{figures/section_ii__contrebasson__p53_28.png}
\caption{Illustration (page 53)}
\end{figure}

egin{figure}[htbp]
\centering
\includegraphics[width=\linewidth]{figures/section_ii__contrebasson__p53_29.png}
\caption{Illustration (page 53)}
\end{figure}

egin{figure}[htbp]
\centering
\includegraphics[width=\linewidth]{figures/section_ii__contrebasson__p53_30.png}
\caption{Illustration (page 53)}
\end{figure}

egin{figure}[htbp]
\centering
\includegraphics[width=\linewidth]{figures/section_ii__contrebasson__p53_31.png}
\caption{Illustration (page 53)}
\end{figure}

egin{figure}[htbp]
\centering
\includegraphics[width=\linewidth]{figures/section_ii__contrebasson__p53_32.jpg}
\caption{Illustration (page 53)}
\end{figure}

egin{figure}[htbp]
\centering
\includegraphics[width=\linewidth]{figures/section_ii__contrebasson__p53_33.png}
\caption{Illustration (page 53)}
\end{figure}

\end{document}
